%%%%%%%%%%%%%%%%%%%%%%%%%%%%%%%%%%%%%%%%%%%%%%%%%%%%%%%%%%%%%%%%%%%%%%
% Lớp 11 | Hàm số mũ và hàm số lôgarít | Phép tính lôgarít
% Dạng 1D6?2-1: Tính giá trị biểu thức chứa lôgarít
% Số câu: 216  (sinh tự động bởi tach_cau.py; từ nay thầy sửa trực tiếp file này)
%%%%%%%%%%%%%%%%%%%%%%%%%%%%%%%%%%%%%%%%%%%%%%%%%%%%%%%%%%%%%%%%%%%%%%

%%%%%%%%%%%%%%%% TRẮC NGHIỆM 4 PHƯƠNG ÁN %%%%%%%%%%%%%%%%

% nguồn: Ok_Toan11_TracNghiem_Lan2
\begin{ex}%[1D6N2-1]%[ma:b785c053]
	 Cho $a>0$ và $a\ne1$, khi đó $\log_{a}\sqrt[2\,021]{a^{2\,022}}$ bằng
	\choice 
		 {$2\,021$}
		{\True $\dfrac{2\,022}{2\,021}$} 
		 {$\dfrac{2\,021}{2\,022}$}
		{$2\,022$} 
	\loigiai{Với $a>0$  và $a\ne1$ ta có $\log_{a}\sqrt[2\,021]{a^{2\,022}}=\log_{a}a^{\tfrac{2\,022}{2\,021}}=\dfrac{2\,022}{2\,021}\log_{a}a=\dfrac{2\,022}{2\,021}$.}
\end{ex}

% nguồn: Ok_Toan11_TracNghiem_Lan2
\begin{ex}%[1D6N2-1]%[ma:4d5c3535]
	Cho $a$ là số thực dương, $a\ne 1$, khi đó $a^{\log_a5}$ bằng 
	\choice 
		{$\log_5a$} 
		{$\log_a5$} 
		{$a^5$} 
		{\True $5$} 
	\loigiai{ 
		Áp dụng công thức  $a^{\log_ab}=b$ ta có $a^{\log_a5}=5$.
	 }
\end{ex}

% nguồn: Ok_Toan11_TracNghiem_Lan2
\begin{ex}%[1D6N2-1]%[ma:09fec02d]
	Cho $a>0$ và $a\ne1$. Khi đó $\log_{\sqrt{a}}a^{2}$ bằng 
	\choice 
		 {$1$}
		 {$\dfrac{1}{4}$}
		 {\True $4$}
		 {$2$}
	\loigiai{Với $a>0$ và $a\ne1$, ta có $\log_{\sqrt{a}}a^{2}=\log_{\sqrt{a}}\left(\sqrt{a}\right)^{4}=4\log_{\sqrt{a}}\sqrt{a}=4 \cdot 1=4$.} 
\end{ex}

% nguồn: Ok_Toan11_TracNghiem_Lan2
\begin{ex}%[1D6N2-1]%[ma:ffaf2671]
	Với $a$ là số thực dương tùy ý, $\log_{4}\left(4a\right)$  bằng
	 \choice
		{$1-\log_{4}a$} 
		{\True  $1+\log_{4}a$}
		{$4-\log_{4}a$} 
		 {$4+\log_{4}a$}
	\loigiai{Với $a>0$  ta có $\log_{4}\left(4a\right)=\log_{4}4+\log_{4}a=1+\log_{4}a$.}
\end{ex}

% nguồn: Ok_Toan11_TracNghiem_Lan2
\begin{ex}%[1D6N2-1]%[ma:6abf4167]
	Cho $a>0$, $a\ne1$, biểu thức $D=\log_{a^{3}}a$ có giá trị bằng bao  nhiêu?
	 \choice
		 {\True $\dfrac{1}{3}$}
		{$3$} 
		{$-\dfrac{1}{3}$} 
		 {$-3$}
	\loigiai{Ta  có $D=\log_{a^{3}}a=\dfrac{1}{3}\log_{a}a=\dfrac{1}{3}$.}
\end{ex}

% nguồn: Ok_Toan11_TracNghiem_Lan2
\begin{ex}%[1D6N2-1]%[ma:e8e2e771]
	Với mọi số thực  $a$ dương khác $1$, $\log_{a}\sqrt[3]{a}$ bằng
	 \choice
		{\True $\dfrac{1}{3}$} 
		{$3$} 
		 {$-3$}
		{$0$} 
	 \loigiai{Ta có $\log_{a}\sqrt[3]{a}=\log_{a}a^{\tfrac{1}{3}}=\dfrac{1}{3}\log_{a}a=\dfrac{1}{3}$.}
\end{ex}

% nguồn: Ok_Toan11_TracNghiem_Lan2
\begin{ex}%[1D6N2-1]%[ma:e2d7905d]
	Giá trị của biểu  thức $4^{\log_{2}\sqrt{3}}$ bằng
	\choice 
		{$\sqrt{3}$} 
		 {\True $3$}
		 {$2^{\sqrt{3}}$}
		{$2\sqrt{3}$} 
	 \loigiai{
		Ta  có $4^{\log_{2}\sqrt{3}}=\left(2^{2}\right)^{\log_{2}\sqrt{3}}=\left(2^{\log_{2}\sqrt{3}}\right)^{2}=\left(\sqrt{3}\right)^{2}=3$.
	 }
\end{ex}

% nguồn: Ok_Toan11_TracNghiem_Lan2
\begin{ex}%[1D6N2-1]%[ma:15d2e082]
	Cho $a>0$ và $a\ne1$, khi  đó $\log_{a}\sqrt[5]{a}$ bằng
	 \choice
		 {\True $\dfrac{1}{5}$}
		 {$-5$}
		{$5$} 
		 {$-\dfrac{1}{5}$}
	\loigiai{Ta  có $\log_{a}\sqrt[5]{a}=\log_{a}\left(a\right)^{\tfrac{1}{5}}=\dfrac{1}{5}$.}
\end{ex}

% nguồn: Ok_Toan11_TracNghiem_Lan2
\begin{ex}%[1D6N2-1]%[ma:4f7a7235]
	Cho $a>0$ và $a\ne1$, khi đó $\left(a^{3}\right)^{\log_{\sqrt{a}}2}$ bằng 
	\choice 
		 {\True $64$}
		 {$\sqrt{8}$}
		 {$12$}
		{$2\sqrt{3}$} 
	 \loigiai{
		 Ta có $\left(a^{3}\right)^{\log_{\sqrt{a}}2}=\left(a^{3}\right)^{2\log_{a}2}=a^{6\log_{a}2}=\left(a^{\log_{a}2}\right)^{6}=2^{6}=64$.
	} 
\end{ex}

% nguồn: Ok_Toan11_TracNghiem_Lan2
\begin{ex}%[1D6N2-1]%[ma:33eae83e]
	Giá trị của $\log_{a}\dfrac{1}{\sqrt{a^{3}}}$ với $a > 0$ và $a  \ne 1$ bằng
	 \choice
		 {\True $-\dfrac{3}{2}$}
		 {$\dfrac{3}{2}$}
		 {$-\dfrac{2}{3}$}
		 {$\dfrac{2}{3}$}
	 \loigiai{Ta có $\log_{a}\dfrac{1}{\sqrt{a^{3}}}=\log_{a}a^{-\tfrac{3}{2}}=-\dfrac{3}{2}$.}
\end{ex}

% nguồn: Ok_Toan11_TracNghiem_Lan2
\begin{ex}%[1D6N2-1]%[ma:934c5192]
	Với $a$ là số  thực dương tùy ý, $\log_2a^5$ bằng
	\choice 
		{$5+\log_2a$} 
		{$\dfrac{1}{5}+\log_2a$} 
		 {$\dfrac{1}{5}{\log_2}a$}
		 {\True $5\log_2a$}
	 \loigiai{
		 Ta có $\log_2a^5=5\log_2a$.
	 }
\end{ex}

% nguồn: Ok_Toan11_TracNghiem_Lan2
\begin{ex}%[1D6N2-1]%[ma:7105b67a]
	Với $a$, $b$ là các số thực dương tùy ý  và $a\ne 1$, $\log_{a^3}b$.
	 \choice
		{$3+\log_ab$} 
		 {$3\log_ab$}
		{$\dfrac{1}{3}+\log_ab$} 
		{\True  $\dfrac{1}{3}\log_ab$}
	\loigiai{ 
		Ta có $\log_{a^3}b=\dfrac{1}{3}\log_ab$. 
	 }
\end{ex}

% nguồn: Ok_Toan11_TracNghiem_Lan2
\begin{ex}%[1D6N2-1]%[ma:97ceed15]
	 Tính giá trị biểu thức $\log_2\left(2\sqrt{2}\right)$
	\choice 
		{$-2$} 
		 {\True $\dfrac{3}{2}$}
		{$3$} 
		 {$-3$}
	 \loigiai{
		Ta có $\log_2\left(2\sqrt{2}\right)=\log_2\left(2\cdot  2^{\tfrac{1}{2}}\right)=\log_2\left(2^{\tfrac{3}{2}}\right)=\dfrac{3}{2}$.
	} 
\end{ex}

% nguồn: Ok_Toan11_TracNghiem_Lan2
\begin{ex}%[1D6N2-1]%[ma:9ddb4ed4]
	Giá trị của  $\log_{2}\dfrac{1}{16}$ bằng
	 \choice
		{$4$} 
		 {$\dfrac{1}{4}$}
		 {$\dfrac{1}{8}$}
		{\True  $-4$}
	\loigiai{Ta có $\log_{2}\dfrac{1}{16}=\log_{2}2^{-4}=-4$.} 
\end{ex}

% nguồn: Ok_TracNghiemLop11_Lan1
\begin{ex}%[1D6N2-1]%[ma:fea95851]
	Cho  $ a $ là số dương khác $ 1 $. Giá trị của $\log_{a^3} a^2$ là
	 \choice
		{\True $\dfrac{2}{3}$} 
		 {$\dfrac{3}{2}$}
		 {$-\dfrac{2}{3}$}
		{$-\dfrac{3}{2}$} 
	 \loigiai{
		 Ta có $\log_{a^3} a^2=\dfrac{2}{3}\log_aa=\dfrac{2}{3}$.
	 }
\end{ex}

% nguồn: Ok_TracNghiemLop11_Lan1
\begin{ex}%[1D6N2-1]%[ma:83f0ab0f]
	 Với $a$ là số thực dương tùy ý, $\log_5 a^2$ bằng
	 \choice
		 {\True $2\log_5 a$}
		{$2+\log_5 a$} 
		 {$\dfrac{1}{2}+\log_5 a$}
		 {$\dfrac{1}{2} \log_5 a$}
	\loigiai{ 
		Ta có  $\log_5 a^2=2\log_5 a$, $\forall a>0$.
	} 
\end{ex}

% nguồn: Ok_TracNghiemLop11_Lan1
\begin{ex}%[1D6N2-1]%[ma:b849ab34]
	Với $a$ là số thực dương tùy ý, $\log_4\left(\dfrac{4}{a}\right)$  bằng?
	\choice 
		{\True  $1-\log_4a$}
		{$\dfrac{1}{\log_4a}$} 
		 {$1+\log_4a$}
		{$4-\log_4a$} 
	\loigiai{$\log_4\left(\dfrac{4}{a}\right)=\log_4  4-\log_4 a=1-\log_4 a.$
	} 
\end{ex}

% nguồn: Ok_TracNghiemLop11_Lan1
\begin{ex}%[1D6N2-1]%[ma:1a38491b]
	 Tính giá trị biểu thức $P=\log_{\tfrac{1}{2}} 4+\log_2 8+2$.
	\choice{$1$}{\True  $3$}{$\log_2 3$}{$0$}
	\loigiai{ 
		Ta có $P=\log_{\tfrac{1}{2}} 4+\log_2 8+2=\log_{2^{(-1)}} 2^2+\log_2 2^3+2=-2\log_2 2+3\log_2  2+2=3$.
	} 
\end{ex}

% nguồn: Ok_TracNghiemLop11_Lan1
\begin{ex}%[1D6N2-1]%[ma:449eb869]
	Với các số  thực dương $a$, $b$ bất kì. Mệnh đề nào sau đây đúng?
	 \choice
		{$\log  (ab)=\log a.\log b$}
		{$\log \dfrac{a}{b}=\log b-\log a$} 
		{$\log \dfrac{a}{b}=\dfrac{\log a}{\log b}$} 
		{\True $\log (ab)=\log  a+\log b$}
	 \loigiai{
		Với các số thực dương  $a$, $b$ bất kì, ta có
		\begin{enumerate}
			\item $\log \dfrac{a}{b}=\log a-\log  b$.
			\item $\log (ab)=\log a+\log b$. 
		\end{enumerate}
	 }
\end{ex}

% nguồn: Ok_TracNghiemLop11_Lan1
\begin{ex}%[1D6N2-1]%[ma:172cd757]
	Cho $\log_a b=2$. Tính $P=\log_a\left(a^3b^4\right)$. 
	 \choice
		{$P=7$} 
		 {$P=10$}
		 {$P=24$}
		{\True  $P=11$}
	 \loigiai{
		Ta có $P=\log_a\left(a^3 b^4\right)=\log_a a^3+\log_a b^4=3+4\log_ab=3+4\cdot  2=11$.
	 }
\end{ex}

% nguồn: Ok_TracNghiemLop11_Lan1
\begin{ex}%[1D6N2-1]%[ma:550ce795]
	Với mọi $a$ dương khác $1$, chọn mệnh đề đúng trong các  mệnh đề sau
	 \choice
		{\True $\log _a a=1$} 
		{$\log _a a^2=2+a$} 
		{$a^{\log _a 1}=a$} 
		{$\log _a 1=1$} 
	 \loigiai{
		Khẳng định  đúng $\log _a a=1$.
	 }
\end{ex}

% nguồn: Ok_TracNghiemLop11_Lan1
\begin{ex}%[1D6N2-1]%[ma:86534d20]
	Cho $a$ là số  thực dương, khác $1$. Cho các mệnh đề sau
	\begin{multicols}{4}
		\begin{itemize}
			\item[(I)] $\log_a  a=a$
			 \item[(II)] $\log_a a^2=2$
			\item[(III)] $\log_a 1=0$ 
			 \item[(IV)] $\log_a 2=2a$
		\end{itemize}
	\end{multicols}
	Trong các mệnh đề  trên, có bao nhiêu mệnh đề \textbf{sai}?
	 \choice
		{\True  $2$}
		{$3$} 
		 {$1$}
		 {$4$}
	\loigiai{ 
		Ta  có
		\begin{itemize}
			\item $\log_a a=1$ 
			\item  $\log_a a^2=2$
			\item  $\log_a 1=0$
			\item $\log_a 2=\log_a  2$
		\end{itemize}
		Mệnh đề  (I) và (IV) sai.
	 }
\end{ex}

% nguồn: Ok_TracNghiemLop11_Lan1
\begin{ex}%[1D6N2-1]%[ma:ef0bb5c5]
	Cho $a$ là số  thực dương khác $1$. Mệnh đề nào dưới đây đúng với mọi số dương $x$, $y$?
	 \choice
		{\True  $\log_a\dfrac{x}{y}=\log_ax-\log_ay$}
		 {$\log_a\dfrac{x}{y}=\log_a(x-y)$}
		{$\log_a\dfrac{x}{y}=\log_ax+\log_ay$} 
		{$\log_a\dfrac{x}{y}=\dfrac{\log_ax}{\log_ay}$} 
	 \loigiai{
		Ta có  $\log_a\dfrac{x}{y}=\log_ax-\log_ay$.
	} 
\end{ex}

% nguồn: Ok_TracNghiemLop11_Lan1
\begin{ex}%[1D6N2-1]%[ma:bf08720e]
	Giá trị của biểu thức  $4^{\log_{\sqrt{2}} 3}$ là
	\choice 
		{$\dfrac{1}{3}$} 
		{$  3 $}
		{\True $ 81 $} 
		{$ 9 $} 
	 \loigiai{
		Ta có $4^{\log_{\sqrt{2}} 3}=4^{2\log_23}=4^{\log_29}=(2^{\log_29})^2=9^2=81$. 
	 }
\end{ex}

% nguồn: Ok_TracNghiemLop11_Lan1
\begin{ex}%[1D6N2-1]%[ma:f3e7eca5]
	Cho các số thực dương $a$,  $b$, $N$ với $a\ne 1$, $b\ne 1$. Mệnh đề nào sau đây đúng?
	\choice 
		{$\log_a N=\log _b N \cdot \log_b  a$}
		{$\log_a N=\dfrac{\log_b a}{\log_b N}$} 
		 {$\log_a N=\log_b(aN)$}
		{\True $\log_a N=\dfrac{\log_b N}{\log_b a}$} 
	\loigiai{ 
		 Ta có $\log_a N=\dfrac{\log_b N}{\log_b a}$.
	 }
\end{ex}

% nguồn: Ok_TracNghiemLop11_Lan1
\begin{ex}%[1D6N2-1]%[ma:7a5ccfdb]
	Với các số thực dương $a, b$  bất kì. Mệnh đề nào sau đây đúng?
	 \choice
		{$\log (a b)=\log  a \cdot \log b$}
		{$\log \dfrac{a}{b}=\log b-\log a$} 
		{$\log \dfrac{a}{b}=\dfrac{\log a}{\log  b}$}
		{\True $\log (a b)=\log a+\log b$} 
	\loigiai{ 
		$\log (a b)=\log a+\log b$ là mệnh  đề đúng.
	 }
\end{ex}

% nguồn: Ok_TracNghiemLop11_Lan1
\begin{ex}%[1D6N2-1]%[ma:33a2bbe6]
	 Với $a$ là số thực dương tùy ý, $4\log \sqrt{a}$ bằng
	 \choice
		{$-4\log a$} 
		 {$8\log a$}
		 {\True $2\log a$}
		{$-2\log  a$}
	 \loigiai{
	Ta có $4\log\sqrt{a}=4\log{a^{\frac{1}{2}}}=2\log a$.} 
\end{ex}

% nguồn: Ok_TracNghiemLop11_Lan1
\begin{ex}%[1D6N2-1]%[ma:e6ce3d1c]
	Cho  $a$ là số thực dương $a\ne 1$ và $\log _{\sqrt[3]{a}} a^3$. Mệnh đề nào sau đây đúng?
	\choice 
		{$P=\dfrac{1}{3}$} 
		 {$P=3$}
		{$P=1$} 
		{\True $P=9$} 
	\loigiai{ 
		$\log _{\sqrt[3]{a}}  a^3=\log _{a^{\tfrac{1}{3}}} a^3=9$.
	} 
\end{ex}

% nguồn: Ok_TracNghiemLop11_Lan1
\begin{ex}%[1D6N2-1]%[ma:0a403002]
	Cho $a>0$ và $a\ne 1$, khi đó $\log_a\sqrt{a}$  bằng
	\choice 
		 {$2$}
		{$-2$} 
		{$-\dfrac{1}{2}$} 
		{\True $\dfrac{1}{2}$} 
	\loigiai{ 
	Với $a>0$  và $a\ne 1$, ta có $\log_a\sqrt{a}=\log_aa^{\frac{1}{2}}=\dfrac{1}{2}\log_aa=\dfrac{1}{2}$.}
\end{ex}

% nguồn: Ok_TracNghiemLop11_Lan1
\begin{ex}%[1D6N2-1]%[ma:51356a55]
	Với mọi số thực dương $a$, $b$, $x$, $y$ và $a$,  $b \ne 1$, mệnh đề nào sau đây \textbf{sai}?
	\choice 
		{\True $\log_{a}{\dfrac{1}{x}}=\dfrac{1}{\log_{a}{x}}$} 
		 {$\log_{a}(xy)=\log_{a}{x}+\log_{a}{y}$}
		{$\log_{b}{a}.\log_{a}{x}=\log_{b}{x}$} 
		{$\log_{a}{\dfrac{x}{y}}=\log_{a}{x}-\log_{a}{y}$} 
	 \loigiai{
		Với  mọi số thực dương $a$, $b$, $x$, $y$ và $a$, $b \ne 1$.\\
		Ta có $\log_{a}{\dfrac{1}{x}}=\log_{a}{x^{-1}}=-\log_{a}{x}\ne  \dfrac{1}{\log_{a}{x}}$.\\
	 }
\end{ex}

% nguồn: Ok_TracNghiemLop11_Lan1
\begin{ex}%[1D6N2-1]%[ma:3c164bdb]
	 Cho $a > 0$ và $a \neq 1$, khi đó $\log _a \sqrt[5]{a}$ bằng
	 \choice
		 {\True $\dfrac{1}{5}$}
		{$-\dfrac{1}{5}$} 
		 {$5$}
		 {$-5$}
	 \loigiai{
		Ta có $\log _a \sqrt[5]{a}=\log  _a a^{\tfrac{1}{5}} = \dfrac{1}{5}$.
	 }
\end{ex}

% nguồn: Ok_TracNghiemLop11_Lan1
\begin{ex}%[1D6N2-1]%[ma:629baf1b]
	Cho $ a $ là số thực dương khác $ 1 $. Khi đó giá trị  của $ a^{\log_a 2024}$ bằng
	\choice 
		 {\True $2024$}
		 {$2024a$}
		 {$2024^a$}
		{$a^{2024}$} 
	\loigiai{ 
		Ta  có $ a^{\log_a 2024}=2024$.
	 }
\end{ex}

% nguồn: Ok_TracNghiemLop11_Lan1
\begin{ex}%[1D6N2-1]%[ma:d53e20bc]
	Cho $a > 0, a \neq 1$. Giá trị  của biểu thức $D=\log _{a^3} a$ bằng
	\choice 
		{$-3$} 
		 {$-\dfrac{1}{3}$}
		{$3$} 
		 {\True $\dfrac{1}{3}$}
	 \loigiai{
		Ta có $D = \log _{a^3} a = \dfrac{1}{3}  \log_a a = \dfrac{1}{3}.$
	} 
\end{ex}

% nguồn: Ok_TracNghiemLop11_Lan1
\begin{ex}%[1D6N2-1]%[ma:f1ea9b23]
	Cho các số thực dương $a$, $M$, $N$ với $a \ne 1$. Mệnh  đề nào sau đây đúng?
	 \choice
		{$\log_a(MN)=\log_a M-\log_a N$}{$\log_a(MN)=\dfrac{\log_a  M}{\log_a N}$}{\True $\log_a(MN)=\log_a M+\log_a N$}{$\log_a(MN)=\log_a M \cdot \log_a N$}
	 \loigiai{
		Ta  có $\log_a(MN)=\log_a M+\log_a N$.
	} 
\end{ex}

% nguồn: Ok_Toan11_TracNghiem_Lan2
\begin{ex}%[1D6H2-1]%[ma:a6d53c1e]
	Xét các số thực dương $a$, $b$ thỏa mãn $\log_5a=5$ và $\log_3b=\dfrac{2}{3}$. Tính giá trị  biểu thức $I=2\log_{6}\left[\log_5(5a)\right]+\log_{\tfrac{1}{9}}b^3$.
	\choice 
		{$I=3$} 
		{$I=-2$} 
		 {\True $I=1$}
		 {$I=2\log_{6}5+1$}
	 \loigiai{
		Ta có \allowdisplaybreaks\begin{align*}
		I&=2\log_{6}\left[\log_5(5a)\right]+\log_{\tfrac{1}{9}}b^3\\ 
		 &=2\log_{6}\left(1+\log_5a\right)-\dfrac{3}{2}\log_3b\\
		 &=2\log_{6}6-\dfrac{3}{2}\cdot\dfrac{2}{3}=2-1=1.
	\end{align*}
 }
\end{ex}

% nguồn: Ok_Toan11_TracNghiem_Lan2
\begin{ex}%[1D6H2-1]%[ma:d9eff443]
	Giá trị  của biểu thức $P=\log_28+\log_{\sqrt{3}}9$ là
	\choice 
		{$6$} 
		{\True $7$} 
		{$8$} 
		 {$4$}
	\loigiai{ 
		Ta có  $P=\log_28+\log_{\sqrt{3}}9=\log_22^3+\log_{\tfrac{1}{3^2}}3^2=3\log_22+4\log_33=3+4=7$
	 }
\end{ex}

% nguồn: Ok_Toan11_TracNghiem_Lan2
\begin{ex}%[1D6H2-1]%[ma:e9a5ae92]
	Cho  $\log_ab=3$, $\log_ac=-4$. Khi đó $P=\log_a\left(\dfrac{a^3\sqrt{c}}{b^2}\right)$ bằng bao nhiêu?
	\choice 
		{$P=7$} 
		 {$P=-1$}
		{$P=11$} 
		 {$P=-5$}
	\loigiai{ Ta có \begin{align*}
		 P&=\log_a\left(\dfrac{a^3\sqrt{c}}{b^2}\right)=\log_aa^3\sqrt{c}-\log_ab^2=\log_aa^3+\log_ac^{\tfrac{1}{2}}-\log_ab^2\\
		 &=3+\dfrac{1}{2}\log_ac-2\log_ab=3+\dfrac{1}{2}\cdot(-4)-2\cdot 3=-5.
	\end{align*}
 }
\end{ex}

% nguồn: Ok_Toan11_TracNghiem_Lan2
\begin{ex}%[1D6H2-1]%[ma:a48b4c61]
	Cho $\log_a3=5$.  Tính $P=\log_a\left(3a^5\right)$.
	 \choice
		{\True  $P=10$}
		 {$P=25$}
		{$P=12$} 
		{$P=125$} 
	 \loigiai{
		 Ta có $\log_a\left(3a^5\right)=\log_a3+\log_aa^5=5+5=10$.
	} 
\end{ex}

% nguồn: Ok_Toan11_TracNghiem_Lan2
\begin{ex}%[1D6H2-1]%[ma:f8cb8c60]
	Cho số thực $a>0$, $a\ne 1$, $a\ne  \dfrac{1}{27}$ và số thực $x$ thỏa mãn $\log_a3=x$. Tính $\log_{27a}9$ theo $x$.
	\choice 
		 {$\dfrac{2x}{x+3}$}
		 {\True $\dfrac{2x}{3x+1}$}
		 {$2(3x+1)$}
		{$\dfrac{2}{3x+1}$} 
	 \loigiai{
		Ta có \begin{align*}
		 \log_{27a}9&=\dfrac{\log_39}{\log_327a}\\
		 &=\dfrac{2}{\log_327+\log_3a}\\
		&=\dfrac{2}{3+\log_3a}\\ 
		&=\dfrac{2}{\dfrac{1}{\log_a3}+3}\\ 
		 &=\dfrac{2}{\dfrac{1}{x}+3}=\dfrac{2x}{3x+1}.
	\end{align*}
 }
\end{ex}

% nguồn: Ok_Toan11_TracNghiem_Lan2
\begin{ex}%[1D6H2-1]%[ma:051f2a9e]
	 Cho $a$, $b$ là các số thực dương khác $1$ thỏa mãn $\log_ab=\sqrt{3}$. Giá trị của $\log_{\tfrac{\sqrt{b}}{a}}\left(\dfrac{\sqrt[3]{b}}{\sqrt{a}}\right)$ là
	\choice 
		 {$-\sqrt{3}$}
		 {$-2\sqrt{3}$}
		 {$\sqrt{3}$}
		{\True  $-\dfrac{1}{\sqrt{3}}$}
	\loigiai{ 
		Ta có \begin{align*}
		\log_{\tfrac{\sqrt{b}}{a}}\left(\dfrac{\sqrt[3]{b}}{\sqrt{a}}\right)&=\dfrac{\log_a\left(\tfrac{\sqrt[3]{b}}{\sqrt{a}}\right)}{\log_a\left(\tfrac{\sqrt{b}}{a}\right)}\\ 
		&=\dfrac{\log_a\sqrt[3]{b}-\log_a\sqrt{a}}{\log_a\sqrt{b}-\log_aa}=\dfrac{\tfrac{1}{3}\log_a b-\tfrac{1}{2}}{\tfrac{1}{2}\log_ab-1}\\ 
		 &=\dfrac{\tfrac{1}{3}\cdot\sqrt{3}-\tfrac{1}{2}}{\tfrac{1}{2}\cdot\sqrt{3}-1}=-\dfrac{1}{\sqrt{3}}.
	\end{align*}
} 
\end{ex}

% nguồn: Ok_Toan11_TracNghiem_Lan2
\begin{ex}%[1D6H2-1]%[ma:33f57b2d]
	Với mọi số  thực $a$ dương và $a\ne 1$, $\log_{a^3}(3a)$ bằng
	\choice 
		 {$\log_a3-1$}
		{$1$} 
		{$3(\log_a3+1)$} 
		{\True $\dfrac{1}{3}(\log_a3+1)$} 
	\loigiai{ 
		 Ta có \[\log_{a^3}(3a)=\dfrac{1}{3} \log_a(3a)=\dfrac{1}{3}(\log_a3+\log_aa)=\dfrac{1}{3}(\log_a3+1).\]
	 }
\end{ex}

% nguồn: Ok_Toan11_TracNghiem_Lan2
\begin{ex}%[1D6H2-1]%[ma:15281cd8]
	Cho $a$ và $b$ là  hai số thực dương thỏa mãn $a^3b^2=32$. Giá trị của $3\log_2a+2\log_2b$ bằng
	\choice 
		 {$4$}
		 {\True $5$}
		{$2$} 
		 {$32$}
	 \loigiai{
		Ta có  $\log_2a^3b^2=\log_232\Leftrightarrow 3\log_2a+2\log_2b=5$
	 }
\end{ex}

% nguồn: Ok_Toan11_TracNghiem_Lan2
\begin{ex}%[1D6H2-1]%[ma:66b3fbf7]
	Với mọi số thực $a$ dương,  $\log_2\dfrac{a}{2}$ bằng
	 \choice
		 {$\dfrac{1}{2}\log_2a$}
		{$\log_2a+1$} 
		{\True $\log_2a-1$} 
		 {$\log_2a-2$}
	 \loigiai{
		Ta  có $\log_2\dfrac{a}{2}=\log_2a-\log_22=\log_2a-1$.
	 }
\end{ex}

% nguồn: Ok_Toan11_TracNghiem_Lan2
\begin{ex}%[1D6H2-1]%[ma:6e7d6ffc]
	Với $a$ là  số thực dương tùy ý, $\log_4(a^3)$ bằng
	 \choice
		 {$3\log_3a$}
		 {$\dfrac{2}{3}\log_2a$}
		{\True  $\dfrac{3}{2}\log_2a$}
		{$3+\log_4a$} 
	\loigiai{ 
		Ta có $\log_4(a^3)=\log_{2^2}(a^3)=\dfrac{3}{2}\log_2a$. 
	 }
\end{ex}

% nguồn: Ok_Toan11_TracNghiem_Lan2
\begin{ex}%[1D6H2-1]%[ma:c948ac21]
	Với $a$, $b$ là hai số dương tùy ý, $\log\left(a^2b^3\right)$ bằng 
	\choice 
		{$\dfrac{1}{2}\log  a+\dfrac{1}{3}\log b$}
		{$  2\log a+\log b$}
		{\True $  2\log a+3\log b$}
		{$\log a+3\log  b$}
	 \loigiai{
		Ta có 
		\[\log\left(a^2b^3\right)=\log{a^2}+\log{b^3}=2\log\left|  a\right|+3\log b=2\log a+3\log b.\]
	} 
\end{ex}

% nguồn: Ok_Toan11_TracNghiem_Lan2
\begin{ex}%[1D6H2-1]%[ma:d86c3c17]
	Cho $\log_{a}b=2$ với $a, b$ là các số thực dương và $a$ khác 1. Giá trị biểu thức $T=\log_{a^{2}}b^{6}+\log_{a}\sqrt{b}$ bằng 
	\choice 
		 {$8$}
		{\True  $7$}
		 {$5$}
		 {$6$}
	\loigiai{Ta có $T=\log_{a^{2}}b^{6}+\log_{a}\sqrt{b}=3\log_{a}b+\dfrac{1}{2}\log_{a}b=\dfrac{7}{2}\log_{a}b=7$.} 
\end{ex}

% nguồn: Ok_Toan11_TracNghiem_Lan2
\begin{ex}%[1D6H2-1]%[ma:0b2645a4]
	Cho $a$, $b$ là các số thực dương và $a \neq 1$ thỏa mãn $\log _a\left(a^2 b^2\right)=1$. Giá  trị của $\log _{a^3} b$ bằng
	 \choice
		 {\True $-\dfrac{1}{6}$}
		 {$\dfrac{1}{3}$}
		{$\dfrac{1}{6}$} 
		{$1$} 
	\loigiai{ 
		Ta có $$\log _a\left(a^2 b^2\right)=1 \Leftrightarrow \log _a a^2+\log _a b^2=1 \Leftrightarrow 2+2 \log  _a b=1 \Leftrightarrow \log _a b=-\dfrac{1}{2}.$$
		Vậy $\log _{a^3}  b=\dfrac{1}{3} \log _a b=-\dfrac{1}{6}$.
	} 
\end{ex}

% nguồn: Ok_Toan11_TracNghiem_Lan2
\begin{ex}%[1D6H2-1]%[ma:e0b0fe90]
	Với  $a$, $b$ là hai số dương tùy ý thì $\log\left(a^{3}b^{2}\right)$ có giá trị bằng biểu thức nào sau đây?
	\choice 
		{$3\log  a+\dfrac{1}{2}\log b$}
		{$2\log a+3\log b$} 
		{\True $3\log a+2\log  b$}
		 {$3\left(\log a+\dfrac{1}{2}\log b\right)$}
	\loigiai{Do $a$, $b$  là hai số dương nên ta có $\log\left(a^{3}b^{2}\right)=\log a^{3}+\log b^{2}=3\log a+2\log b$.}
\end{ex}

% nguồn: Ok_Toan11_TracNghiem_Lan2
\begin{ex}%[1D6H2-1]%[ma:e8f238c4]
	Cho $a$, $b$ là các số thực dương khác $1$ thoả mãn $\log_2a=2$  và $\log_4b=3$. Giá trị biểu thức $P=\log_a\left(a^2b\right)$ bằng
	\choice 
		 {$P=10$}
		 {\True $P=5$}
		{$P=2$} 
		 {$P=1$}
	\loigiai{ 
		Với $a$, $b$ là các số thực dương khác  $1$ ta có
		\begin{itemize}
			\item  $\log_2a=2\Leftrightarrow a=2^2=4$.
			\item $\log_4b=3\Leftrightarrow b=4^3$. 
		\end{itemize}
		Vậy $P=\log_a\left(a^2b\right)=\log_4\left(4^24^3\right)=5$. 
	 }
\end{ex}

% nguồn: Ok_Toan11_TracNghiem_Lan2
\begin{ex}%[1D6H2-1]%[ma:e9e384b1]
	Với mọi số thực $a$ dương, $\log_3\left(\dfrac{a^2}{9}\right)$ bằng 
	\choice 
		{\True $2(\log_3a-1)$} 
		{$\dfrac{1}{2} \log_3a$} 
		 {$\log_3a-1$}
		{$\log_3a+2$} 
	 \loigiai{
		Ta  có \[\log_3\left(\dfrac{a^2}{9}\right)=\log_3(a^2)-\log_39=2\log_3a-2=2(\log_3a-1).\]
	} 
\end{ex}

% nguồn: Ok_Toan11_TracNghiem_Lan2
\begin{ex}%[1D6H2-1]%[ma:42b5f642]
	Cho $a>0$ và $a\ne 1$, khi đó  $\log_{a^2}\sqrt[4]{a^3}$ bằng
	 \choice
		{$-\dfrac{3}{8}$} 
		 {$\dfrac{3}{4}$}
		 {\True $\dfrac{3}{8}$}
		 {$\dfrac{3}{2}$}
	 \loigiai{
		Ta có  $\log_{a^2}\sqrt[4]{a^3}=\dfrac{1}{2}\cdot \log_a{a^{\tfrac{3}{4}}}=\dfrac{1}{2}\cdot \dfrac{3}{4}\cdot \log_aa=\dfrac{3}{8}$.
	} 
\end{ex}

% nguồn: Ok_Toan11_TracNghiem_Lan2
\begin{ex}%[1D6H2-1]%[ma:338129ac]
	 Giá trị của biểu thức $4^{\log_2\sqrt 3}$ bằng
	\choice 
		{$\sqrt  3$}
		{\True $3$} 
		 {$2^{\sqrt 3}$}
		{$2\sqrt  3$}
	\loigiai{ 
		 Ta có
		\[4^{\log_2\sqrt 3}=\left(2^2\right)^{\log_2\sqrt 3}=\left(2^{\log_2\sqrt 3}\right)^2=\left(\sqrt 3\right)^2=3.\] 
	} 
\end{ex}

% nguồn: Ok_Toan11_TracNghiem_Lan2
\begin{ex}%[1D6H2-1]%[ma:a566ef34]
	Với $a$ là số  thực dương tùy ý, $\log_22a$ bằng
	\choice 
		{\True $1+\log_2a$} 
		{$1-\log_2a$} 
		{$2-\log_2a$} 
		 {$2+\log_2a$}
	\loigiai{ 
		Ta có $\log_22a=\log_22+\log_2a=1+\log_2a$. 
	} 
\end{ex}

% nguồn: Ok_Toan11_TracNghiem_Lan2
\begin{ex}%[1D6H2-1]%[ma:da686503]
	Cho $a>0$ và $a\ne 1$, khi đó $\log_a\sqrt[2021]{a^{2022}}$ bằng 
	 \choice
		 {$2021$}
		 {\True $\dfrac{2022}{2021}$}
		 {$\dfrac{2021}{2022}$}
		 {$2022$}
	 \loigiai{
		Với $a>0$  và $a\ne 1$ ta có $\log_a\sqrt[202]{a^{2022}}=\log_aa^{\tfrac{2022}{2021}}=\dfrac{2022}{2021}\cdot\log_aa=\dfrac{2022}{2021}$.
	} 
\end{ex}

% nguồn: Ok_Toan11_TracNghiem_Lan2
\begin{ex}%[1D6H2-1]%[ma:da85ccd2]
	Tính giá trị của biểu thức $P=2^{\log_{z}a}+\log_a\left(a^b\right)$ với $a>0$,  $a\ne 1$.
	\choice 
		 {$P=2^{a}+b$}
		{$P=a-b$} 
		{$P=2a+b$} 
		{\True $P=a+b$} 
	\loigiai{ 
		Ta có  $P=2^{\log_2a}+\log_a\left(a^b\right)=a+b$.
	 }
\end{ex}

% nguồn: Ok_Toan11_TracNghiem_Lan2
\begin{ex}%[1D6H2-1]%[ma:44e4f204]
	Cho $\log_{a}b=2$;  $\log_{a}c=3$. Tính $Q=\log_{a}\left(b^{3}c\right)$.
	\choice 
		{$Q=4$} 
		{\True $Q=9$} 
		 {$Q=10$}
		 {$Q=12$}
	\loigiai{Ta  có $Q=\log_{a}\left(b^{3}c\right)=3\log_{a}b+\log_{a}c=3 \cdot 2+3=9$.}
\end{ex}

% nguồn: Ok_Toan11_TracNghiem_Lan2
\begin{ex}%[1D6H2-1]%[ma:36313c11]
	Cho  $a$ là số thực dương khác $5$. Tính $I=\log_{\tfrac{a}{5}}\left(\dfrac{a^{3}}{125}\right)$.
	\choice 
		 {$I=-\dfrac{1}{3}$}
		{$I=-3$} 
		 {$I=\dfrac{1}{3}$}
		 {\True $I=3$}
	\loigiai{Ta có $I=\log_{\tfrac{a}{5}}\left(\dfrac{a^{3}}{125}\right)=\log_{\tfrac{a}{5}}\left(\dfrac{a}{5}\right)^{3}=3$.} 
\end{ex}

% nguồn: Ok_Toan11_TracNghiem_Lan2
\begin{ex}%[1D6H2-1]%[ma:fcc29ffb]
	Với $a$ là số  thực dương tùy ý, $\log\left(100a^3\right)$ bằng
	\choice 
		 {$6\log a$}
		{\True $2+3\log a$} 
		 {$\dfrac{1}{2}+\dfrac{1}{3}\log a$}
		{$3+3\log  a$}
	\loigiai{ 
		Ta có \[\log(100a^3)=\log 100+\log a^3=\log 10^2+\log a^3=2+3\log a.\] 
	 }
\end{ex}

% nguồn: Ok_Toan11_TracNghiem_Lan2
\begin{ex}%[1D6H2-1]%[ma:b224542e]
	Với mọi $a$, $b$ thỏa mãn  $\log_3\left(3a^2\right)+\log_3b^3=4$, khẳng định nào dưới đây đúng?
	\choice 
		 {$a^2b^3=81$}
		 {$a^3+b^3+1=81$}
		{\True $a^2b^3=27$} 
		 {$a^2+b^3=27$}
	 \loigiai{
		 Ta có
		\[\log_33a^2+\log_3b^3=4\Leftrightarrow \log_3\left(3a^2\cdot b^3\right)=4\Leftrightarrow  3a^2b^3=81\Leftrightarrow a^2b^3=27.\]
	} 
\end{ex}

% nguồn: Ok_Toan11_TracNghiem_Lan2
\begin{ex}%[1D6H2-1]%[ma:5a43e560]
	Tính giá trị  biểu thức $3^{2\log_3 5}$
	 \choice
		{$10$} 
		 {\True $25$}
		 {$7$}
		 {$20$}
	\loigiai{ 
		Ta  có $3^{2\log_3 5}=3^{\log_35^2}=5^2=25$.
	 }
\end{ex}

% nguồn: Ok_Toan11_TracNghiem_Lan2
\begin{ex}%[1D6H2-1]%[ma:49c4561b]
	Cho $a,b$ là  các số thực dương $\left(a,b\ne1\right)$ và $\log_{a}b=16$. Tính giá trị của biểu thức $P=\log_{a}\sqrt{b}$.
	\choice 
		{$256$} 
		{$4$} 
		 {$23$}
		{\True $8$} 
	\loigiai{Theo tính chất  của logarit ta có $P=\log_{a}\sqrt{b}=\log_{a}\left(b\right)^{\tfrac{1}{2}}=\dfrac{1}{2}\log_{a}b=\dfrac{1}{2} \cdot 16=8$.}
\end{ex}

% nguồn: Ok_Toan11_TracNghiem_Lan2
\begin{ex}%[1D6H2-1]%[ma:8829f20f]
	 Với $a$ là số thực dương tùy ý, $\log_3(3a^2)$ bằng
	 \choice
		 {$3+2\log_3a$}
		 {$1+\dfrac{1}{2} \log_3a$}
		 {\True $1+2\log_3a$}
		 {$1-2\log_3a$}
	\loigiai{ 
		Ta  có $\log_3\left(3a^2\right)=\log_33+\log_3a^2=1+2\log_3a$.
	 }
\end{ex}

% nguồn: Ok_Toan11_TracNghiem_Lan2
\begin{ex}%[1D6H2-1]%[ma:a5ea8b8a]
	Cho hai số tự nhiên $x, y$ thỏa mãn $x \log _{28} 2+y \log _{28} 7=2$. Giá trị của $x+y$  bằng
	 \choice
		{$5$} 
		{\True $6$} 
		{$4$} 
		{$8$} 
	 \loigiai{
		Ta  có
		 \allowdisplaybreaks
		\begin{eqnarray*}
			x \log _{28} 2+y \log _{28} 7=2  &\Leftrightarrow& \log _{28}\left(2^x 7^y\right)=2 \Leftrightarrow 2^x 7^y=28^2 \\&\Leftrightarrow& 2^x 7^y=\left(2^2 7\right)^2 \Leftrightarrow 2^x 7^y=2^4 7^2.
		\end{eqnarray*}
		Vì $x$, $y$ là số tự nhiên nên $x=4$,  $y=2 \Rightarrow x+y=6$.
	 }
\end{ex}

% nguồn: Ok_Toan11_TracNghiem_Lan2
\begin{ex}%[1D6H2-1]%[ma:786340fa]
	Với mọi số  thực $a$ dương, $\log_2\dfrac{a^2}{4}$ bằng
	 \choice
		 {\True $2\left(\log_2a-1\right)$}
		 {$\log_2a-2$}
		 {$\log_2a-1$}
		{$2\log_2a-1$} 
	\loigiai{ 
		Ta có  $\log_2\dfrac{a^2}{4}=\log_2a^2-\log_24=2\log_2a-2=2\left(\log_2a-1\right)$.
	 }
\end{ex}

% nguồn: Ok_Toan11_TracNghiem_Lan2
\begin{ex}%[1D6H2-1]%[ma:818ab86c]
	 Cho $P=\sqrt[10]{3\sqrt[5]{27\sqrt[2]{243}}}$. Tính $\log_3P$.
	 \choice
		 {$\dfrac{45}{28}$}
		{\True $\dfrac{21}{100}$} 
		 {$\dfrac{45}{56}$}
		{$\dfrac{13}{100}$} 
	\loigiai{ 
		Ta có 
		\[P=\sqrt[10]{3\sqrt[5]{27\sqrt[2]{243}}}=3^{\tfrac{1}{10}}\cdot 27^{\tfrac{1}{10}\cdot \tfrac{1}{5}}\cdot 243^{\tfrac{1}{10}\cdot \tfrac{1}{5}\cdot  \tfrac{1}{2}}=3^{\tfrac{21}{100}}.\]
		Suy ra $\log_3P=\log_33^{\tfrac{21}{100}}=\dfrac{21}{100}$. 
	} 
\end{ex}

% nguồn: Ok_Toan11_TracNghiem_Lan2
\begin{ex}%[1D6H2-1]%[ma:73edcdac]
	Tính giá trị của biểu thức $P=2^{\log_{2}a}+\log_{a}\left(a^{b}\right)$  $\left(a>0,a\ne1\right)$.
	 \choice
		{$P=2^{a}+b$} 
		 {$P=a-b$}
		{$P=2a+b$} 
		 {\True $P=a+b$}
	\loigiai{Ta có  $P=2^{\log_{2}a}+\log_{a}\left(a^{b}\right)=a+b$.}
\end{ex}

% nguồn: Ok_Toan11_TracNghiem_Lan2
\begin{ex}%[1D6H2-1]%[ma:bca63dbf]
	Cho $a=\log _7 5$, $b=\log _3  5$. Biểu thức $M=\log _{21} 5$ bằng
	\choice 
		 {\True $\dfrac{a b}{a+b}$}
		 {$a b$}
		{$\dfrac{1}{a b}$} 
		 {$\dfrac{a+b}{a b}$}
	 \loigiai{
		 Ta có
		\allowdisplaybreaks 
		\begin{eqnarray*}
			M&=&\log _{21} 5=\dfrac{1}{\log _5 21}=\dfrac{1}{\log _5(7.3)}=\dfrac{1}{\log _5 7+\log _5  3}\\&=&\dfrac{1}{\dfrac{1}{\log _{7} 5}+\dfrac{1}{\log _3 5}}=\dfrac{1}{\dfrac{1}{a}+\dfrac{1}{b}}=\dfrac{a b}{a+b}
			. 
		\end{eqnarray*}
	 }
\end{ex}

% nguồn: Ok_Toan11_TracNghiem_Lan2
\begin{ex}%[1D6H2-1]%[ma:14f6545f]
	Cho hàm số  $f\left(x\right)=\log_{2}x$. Với $x>0$, giá trị của biểu thức $P=f\left(\dfrac{6}{x}\right)+f\left(\dfrac{8x}{3}\right)$ bằng
	\choice 
		 {$P=2$}
		{$P=1$} 
		{\True $P=4$} 
		 {$P=3$}
	\loigiai{Ta có  $P=f\left(\dfrac{6}{x}\right)+f\left(\dfrac{8x}{3}\right)=f\left(\dfrac{6}{x} \cdot \dfrac{8x}{3}\right)=f\left(16\right)=4$.}
\end{ex}

% nguồn: Ok_Toan11_TracNghiem_Lan2
\begin{ex}%[1D6H2-1]%[ma:1a61ac02]
	Cho $a, b$ là các số thực dương khác $1$ thỏa mãn $\log_{a}b=\sqrt{3}$. Giá  trị của $\log_{\tfrac{\sqrt{b}}{a}}\left(\dfrac{\sqrt[3]{b}}{\sqrt{a}}\right)$ là
	 \choice
		 {$-\sqrt{3}$}
		{$-2\sqrt{3}$} 
		{$\sqrt{3}$} 
		 {\True $-\dfrac{1}{\sqrt{3}}$}
	\loigiai{Ta có  $\log_{\tfrac{\sqrt{b}}{a}}\left(\dfrac{\sqrt[3]{b}}{\sqrt{a}}\right)=\dfrac{\log_{a}\left(\dfrac{\sqrt[3]{b}}{\sqrt{a}}\right)}{\log_{a}\dfrac{\sqrt{b}}{a}}=\dfrac{\log_{a}\sqrt[3]{b}-\log_{a}\sqrt{a}}{\log_{a}\sqrt{b}-\log_{a}a}=\dfrac{\dfrac{1}{3}\log_{a}b-\dfrac{1}{2}}{\dfrac{1}{2}\log_{a}b-1}=\dfrac{\dfrac{\sqrt{3}}{3}-\dfrac{1}{2}}{\dfrac{\sqrt{3}}{2}-1}=-\dfrac{1}{\sqrt{3}}$.}
\end{ex}

% nguồn: Ok_Toan11_TracNghiem_Lan2
\begin{ex}%[1D6H2-1]%[ma:8c68fed8]
	Với mọi số thực  $a$ dương, $\log_2^2a^2$ bằng
	\choice 
		{$  2\log_2^2a$}
		 {\True $ 4\log_2^2a$}
		 {$ 2\log_2a^2$}
		{$ 4\log_2a$} 
	 \loigiai{
		 Ta có $\log_2^2a^2=\left(2\log_2a\right)^2=4\log_2^2a$.
	} 
\end{ex}

% nguồn: Ok_Toan11_TracNghiem_Lan2
\begin{ex}%[1D6H2-1]%[ma:50e75a43]
	Với  $a$, $b$ là các số thực dương tùy ý và $a\ne 1$. Ta có $\log_a\left(a^3b\right)$ bằng
	\choice 
		 {$3\cdot \log_ab$}
		{$\dfrac{1}{3}\cdot  \log_ab$}
		{$\dfrac{1}{3}+\log_ab$} 
		{\True  $3+\log_ab$}
	\loigiai{ 
		Ta có 
		\[\log_a\left(a^3b\right)=\log_a{a^3}+\log_ab=3+\log_ab\text{, với }a,b>0;a\ne 1.\] 
	} 
\end{ex}

% nguồn: Ok_Toan11_TracNghiem_Lan2
\begin{ex}%[1D6H2-1]%[ma:5f14b734]
	Với mọi số  thực $a$ dương, $\log_5(5a)$ bằng
	 \choice
		{\True $1+\log_5a$} 
		{$\log_5a$} 
		 {$5\log_5a$}
		{$1-\log_5a$} 
	 \loigiai{
		Suy ra $\log_5(5a)=\log_55+\log_5a=1+\log_5a$. 
	} 
\end{ex}

% nguồn: Ok_Toan11_TracNghiem_Lan2
\begin{ex}%[1D6H2-1]%[ma:6104483f]
	Cho $a>0$ và  $a\ne 1$, khi đó $\log_a\sqrt[4]{a}$ bằng
	\choice 
		 {$4$}
		{\True $\dfrac{1}{4}$} 
		 {$-\dfrac{1}{4}$}
		{$-4$} 
	\loigiai{ 
		Ta có  $\log_a\sqrt[4]{a}=\log_aa^{\tfrac{1}{4}}=\dfrac{1}{4}$.
	 }
\end{ex}

% nguồn: Ok_Toan11_TracNghiem_Lan2
\begin{ex}%[1D6H2-1]%[ma:8fe74349]
	Cho số thực  $a>0$; $a\ne 1$, $a\ne\dfrac{1}{27}$ và số thực $x$ thỏa mãn $\log_a3=x$. Tính $\log_{27a}9$ theo $x$.
	 \choice
		 {$\dfrac{2x}{x+3}$}
		{\True  $\dfrac{2x}{3x+1}$}
		{$2\left(3x+1\right)$} 
		{$\dfrac{2}{3x+1}$} 
	\loigiai{ 
		Ta có 
		 \allowdisplaybreaks
		\begin{eqnarray*}
			\log_{27a}9&=&\dfrac{\log_39}{\log_327a}=\dfrac{2}{\log_327+\log_3a}=\dfrac{2}{3+\log_3a}=\dfrac{2}{\dfrac{1}{\log_a3}+3}\\ 
			 &=&\dfrac{2}{\dfrac{1}{x}+3}=\dfrac{2x}{3x+1}.
		\end{eqnarray*}
	} 
\end{ex}

% nguồn: Ok_Toan11_TracNghiem_Lan2
\begin{ex}%[1D6H2-1]%[ma:2be1b546]
	Với $a$ là số thực  dương tùy ý, $\log_{\tfrac{1}{2}}\left(8a\right)$ bằng
	\choice 
		 {$\dfrac{1}{2}+\log_2a$}
		 {$-3+\log_2a$}
		{$-\left(\log_2a\right)^3$} 
		{\True  $-3-\log_2a$}
	 \loigiai{
		 Ta có \[\log_{\tfrac{1}{2}}(8a)=-(\log_28+\log_2a)=-(\log_22^3+\log_2a)=-3-\log_2a.\]
	 }
\end{ex}

% nguồn: Ok_Toan11_TracNghiem_Lan2
\begin{ex}%[1D6H2-1]%[ma:43c713c3]
	Cho $a$, $b$ là các  số thực dương thỏa mãn $\log_ab=2$. Tính giá trị của biểu thức $P=\log_{\tfrac{\sqrt a}{b}}\left(a\cdot \sqrt[3]{b}\right)$.
	\choice 
		{$P=\dfrac{2}{15}$} 
		 {$P=-\dfrac{2}{9}$}
		{\True  $P=-\dfrac{10}{9}$}
		{$P=\dfrac{2}{3}$} 
	\loigiai{ 
		Ta có 
		\allowdisplaybreaks 
		\begin{eqnarray*}
			P&=&\log_{\tfrac{\sqrt a}{b}}\left(a\cdot \sqrt[3]{b}\right)=\dfrac{\log_a\left(a\cdot \sqrt[3]{b}\right)}{\log_a\left(\dfrac{\sqrt  a}{b}\right)}=\dfrac{\log_aa+\log_a{b^{\tfrac{1}{3}}}}{\log_a{a^{\tfrac{1}{2}}}-\log_ab}\\
			&=&\dfrac{1+\dfrac{1}{3}  \log_ab}{\tfrac{1}{2}-\log_ab}=\dfrac{1+\dfrac{1}{3}\cdot 2}{\dfrac{1}{2}-2}=-\dfrac{10}{9}.
		\end{eqnarray*}
	} 
\end{ex}

% nguồn: Ok_Toan11_TracNghiem_Lan2
\begin{ex}%[1D6H2-1]%[ma:98001e99]
	 Với $a$ là số thực dương tùy ý, $\log_5\dfrac{25}{a}$ bằng
	\choice 
		{\True  $2-\log_5a$}
		{$\dfrac{5}{\log_5a}$} 
		{$5-\log_5a$} 
		 {$\dfrac{2}{\log_5a}$}
	\loigiai{ 
		 Ta có \[\log_5\dfrac{25}{a}=\log_525-\log_5a=\log_55^2-\log_5a=2-\log_5a.\]
	} 
\end{ex}

% nguồn: Ok_Toan11_TracNghiem_Lan2
\begin{ex}%[1D6H2-1]%[ma:5a123a43]
	Cho $\log _2 5=a$, $ \log _5 3=b$. Tính $\log _5 24$ theo $a$ và  $b$.
	\choice 
		{\True $\log  _5 24=\dfrac{3+a b}{a}$}
		{$\log _5  24=\dfrac{a+3 b}{a}$}
		{$\log _5 24=\dfrac{a+b}{3 a b}$} 
		{$\log _5 24=\dfrac{3  a+b}{b}$}
	\loigiai{Ta có $$\log _5 24=\log _5\left(3 \cdot 2^3\right)=\log _5 3+3 \log _5 2=\log  _5 3+\dfrac{3}{\log _2 5}=b+\dfrac{3}{a}=\dfrac{a b+3}{a}.$$
	} 
\end{ex}

% nguồn: Ok_Toan11_TracNghiem_Lan2
\begin{ex}%[1D6H2-1]%[ma:5b35a2db]
	Cho $a$, $b$ là các số thực dương và $a$ khác 1, thỏa mãn $\log_{a^{2}}\left(\dfrac{a^{3}}{\sqrt[5]{b^{3}}}\right)=3$. Giá trị của biểu  thức $\log_{a}b$ bằng
	 \choice
		 {\True $-5$}
		{$5$} 
		{$\dfrac{1}{5}$} 
		 {$-\dfrac{1}{5}$}
	\loigiai{Ta  có $\log_{a^{2}}\left(\dfrac{a^{3}}{\sqrt[5]{b^{3}}}\right)=3\Leftrightarrow\dfrac{1}{2}\left(\log_{a}a^{3}-\log_{a}b^{\tfrac{3}{5}}\right)=3\Leftrightarrow3-\dfrac{3}{5}\log_{a}b=6\Leftrightarrow \log_{a}b=-5$.}
\end{ex}

% nguồn: Ok_Toan11_TracNghiem_Lan2
\begin{ex}%[1D6H2-1]%[ma:548e08be]
	 Cho $\log _2 3=a$. Giá trị của biểu thức $P=\log _6 12$ tính theo $a$ bằng
	\choice 
		 {$\dfrac{a}{2+a}$}
		 {$\dfrac{1+a}{2+a}$}
		{$\dfrac{a}{1+a}$} 
		{\True $\dfrac{2+a}{1+a}$} 
	 \loigiai{
		Ta có $$P=\log _6 12=\dfrac{\log _2 12}{\log _2 6}=\dfrac{\log _2(4.3)}{\log _2(2\cdot 3)}=\dfrac{2+\log _2 3}{1+\log _2 3}=\dfrac{2+a}{1+a}.$$ 
	} 
\end{ex}

% nguồn: Ok_Toan11_TracNghiem_Lan2
\begin{ex}%[1D6H2-1]%[ma:c0f3f47f]
	Với $a$ là số thực dương tùy ý, $\log_3\left(\dfrac{3}{a}\right)$ bằng 
	 \choice
		{\True $1-\log_3a$} 
		 {$3-\log_3a$}
		{$\dfrac{1}{\log_3a}$} 
		{$1+\log_3a$} 
	\loigiai{ 
		 Ta có $\log_3\left(\dfrac{3}{a}\right)=\log_33-\log_3a=1-\log_3a$.
	 }
\end{ex}

% nguồn: Ok_Toan11_TracNghiem_Lan2
\begin{ex}%[1D6H2-1]%[ma:6e702435]
	Cho $a$  là số thực dương khác $2$. Tính $I=\log_{\tfrac{a}{2}}\left(\dfrac{a^{2}}{4}\right)$.
	\choice 
		{$I=\dfrac{1}{2}$} 
		{$I=-\dfrac{1}{2}$} 
		{\True  $I=2$}
		 {$I=-2$}
	\loigiai{Ta  có $I=\log_{\tfrac{a}{2}}\left(\dfrac{a^{2}}{4}\right)=\log_{\tfrac{a}{2}}\left(\dfrac{a}{2}\right)^{2}=2\log_{\tfrac{a}{2}}\left(\dfrac{a}{2}\right)=2$.}
\end{ex}

% nguồn: Ok_Toan11_TracNghiem_Lan2
\begin{ex}%[1D6H2-1]%[ma:95ab68f8]
	Với $a$  là số thực dương tùy ý, $\log_5(5a)$ bằng
	 \choice
		 {$5+\log_5a$}
		 {$5-\log_5a$}
		 {\True $1+\log_5a$}
		 {$1-\log_5a$}
	\loigiai{ 
		Ta có $\log_5(5a)=\log_55+\log_5a=1+\log_5a$. 
	 }
\end{ex}

% nguồn: Ok_Toan11_TracNghiem_Lan2
\begin{ex}%[1D6H2-1]%[ma:730d27f9]
	Cho $a$ là  số thực dương. Khi đó $\log_48a^3$ bằng
	\choice 
		 {$\dfrac{3}{2}+\log_2a$}
		{\True  $\dfrac{3}{2}+\dfrac{3}{2}{\log_2}a$}
		{$2+3\log_2a$} 
		{$6+6\log_2a$} 
	 \loigiai{
		Ta có \[\log_48a^3=\log_48+\log_4a^3=\log_{2^2}{2^3}+3\log_{2^2}a=\dfrac{3}{2}+\dfrac{3}{2}{\log_2}a.\] 
	 }
\end{ex}

% nguồn: Ok_Toan11_TracNghiem_Lan2
\begin{ex}%[1D6H2-1]%[ma:81fbe539]
	Cho $a>0$, $a\ne 1$, khi đó  $\log_a\left(a\cdot \sqrt[3]{a}\right)$ bằng
	\choice 
		{$ 4$} 
		{\True $\dfrac{4}{3}$} 
		 {$\dfrac{-4}{3}$}
		{$\dfrac{1}{3}$} 
	 \loigiai{
		Với $a>0$, $a\ne 1$  nên ta có
		\[\log_a\left(a\cdot  \sqrt[3]{a}\right)=\log_aa+\log_a\sqrt[3]{a}=1+\dfrac{1}{3}=\dfrac{4}{3}.\]
		Vậy $\log_a\left(a\cdot \sqrt[3]{a}\right)=\dfrac{4}{3}$. 
	 }
\end{ex}

% nguồn: Ok_Toan11_TracNghiem_Lan2
\begin{ex}%[1D6H2-1]%[ma:50a6f5ed]
	Cho $\log _{15} 30=\dfrac{1+a \log 3}{b \log 3+c \log 5}$, với $a,  b, c$ là các số nguyên. Giá trị của $a b+c$ bằng
	 \choice
		 {$4$}
		{$3$} 
		{$1$} 
		{\True  $2$}
	\loigiai{ 
		Ta có $$\log _{15} 30=\dfrac{\log 30}{\log 15}=\dfrac{\log 10+\log 3}{\log 3+\log 5}=\dfrac{1+\log  3}{\log 3+\log 5}.$$
		Suy ra $a=1$, $b=1$,  $c=1 \Rightarrow a b+c=2$.
	 }
\end{ex}

% nguồn: Ok_Toan11_TracNghiem_Lan2
\begin{ex}%[1D6H2-1]%[ma:9f7d83c8]
	Với $a$,  $b$ là các số thực dương bất kỳ, $\log_2\dfrac{a}{b^4}$ bằng
	\choice 
		{$\log_2a-\log_2\left(4b\right)$} 
		 {$\dfrac{1}{4}{\log_2}\dfrac{a}{b}$}
		 {$ 2\log_2\dfrac{a}{b}$}
		{\True $\log_2a-4\log_2b$} 
	 \loigiai{
		Ta  có $\log_2\dfrac{a}{b^4}=\log_2a-\log_2b^4=\log_2a-4\log_2b$.
	 }
\end{ex}

% nguồn: Ok_Toan11_TracNghiem_Lan2
\begin{ex}%[1D6H2-1]%[ma:5df06609]
	Giả sử $a$, $b$ là  các số thực dương tùy ý thỏa mãn $a^2b^3=4^4$. Mệnh đề nào sau đây đúng?
	\choice 
		{\True $2\log_2a+3\log_2b=8$} 
		 {$2\log_2a-3\log_2b=8$}
		 {$2\log_2a-3\log_2b=4$}
		{$2\log_2a+3\log_2b=4$} 
	\loigiai{ 
		Ta có $\log_2a^2b^3=\log_24^4\Leftrightarrow  2\log_2a+3\log_2b=8$.
	 }
\end{ex}

% nguồn: Ok_Toan11_TracNghiem_Lan2
\begin{ex}%[1D6H2-1]%[ma:2053b8e7]
	Với $a$, $b$ là hai số  thực dương tùy ý, $\log_3\left(ab^3\right)$ bằng
	\choice 
		{$\log_3a+\dfrac{1}{3}\log_3b$} 
		{$3(\log_3a+\log_3b)$} 
		{\True  $\log_3a+3\log_3b$}
		{$3\log_3a+\log_3b$} 
	 \loigiai{
		Ta  có $\log_3\left(ab^3\right)=\log_3a+\log_3b^3=\log_3a+3\log_3b$.
	} 
\end{ex}

% nguồn: Ok_Toan11_TracNghiem_Lan2
\begin{ex}%[1D6H2-1]%[ma:6fc767b3]
	Cho $\log_ab=2$  và $\log_ac=3$. Tính $P=\log_a\left(b^2c^3\right)$.
	 \choice
		 {\True $P=13$}
		{$P=31$} 
		 {$P=30$}
		{$P=108$} 
	 \loigiai{
		 Ta có $\log_a\left(b^2c^3\right)=2\log_ab+3\log_ac=2\cdot 2+3\cdot 3=13$.
	} 
\end{ex}

% nguồn: Ok_Toan11_TracNghiem_Lan2
\begin{ex}%[1D6H2-1]%[ma:61700195]
	Cho các số dương $a$, $b$ khác $ 1$ sao cho $\log_{16}\sqrt[3]{a}=\log_{a^2}\sqrt[9]{b}=\log_b2$. Giá trị  của $\dfrac{b}{a^3}$ bằng
	\choice 
		{\True $1$} 
		 {$2$}
		{$4$} 
		{$8$} 
	\loigiai{ 
		 Đặt $\log_{16}\sqrt[3]{a}=\log_{a^2}\sqrt[9]{b}=\log_b2=t\Leftrightarrow \heva{&a=16^{3t}\\&b=a^{18t}\\&2=b^t}\Leftrightarrow \heva{&a=16^{3t}\\&b=16^{54t^2}\\&2=16^{54t^3}.}$\hfill (*)\\
		Từ $(*)$ suy ra $2=2^{216t^3}\Leftrightarrow 216t^3=1\Leftrightarrow t=\dfrac{1}{6}$.\\ 
		Suy ra $a=4$, $b=4^3$.\\ 
		Vậy  $\dfrac{b}{a^3}=1$.
	} 
\end{ex}

% nguồn: Ok_Toan11_TracNghiem_Lan2
\begin{ex}%[1D6H2-1]%[ma:0741e686]
	Cho $a$,  $b$ là các số thực dương và $a\ne 1$ thoả mãn $\log_a\left(a^2b\right)=\dfrac{1}{2}$. Giá trị của $\log_{a^2}b$ bằng
	 \choice
		{\True  $-\dfrac{3}{4}$}
		 {$-\dfrac{3}{2}$}
		{$-3$} 
		{$\dfrac{3}{4}$} 
	 \loigiai{
		Từ  giả thiết ta có
		\[\log_a\left(a^2b\right)=\dfrac{1}{2}\Leftrightarrow \log_a{a^2}+\log_ab=\dfrac{1}{2}\Leftrightarrow \log_ab=\dfrac{1}{2}-2=-\dfrac{3}{2}.\] 
		 Vậy $\log_{a^2}b=\dfrac{1}{2} \log_ab=\dfrac{1}{2}\cdot \left(-\dfrac{3}{2}\right)=-\dfrac{3}{4}$.
	} 
\end{ex}

% nguồn: Ok_Toan11_TracNghiem_Lan2
\begin{ex}%[1D6H2-1]%[ma:5b15fdbe]
	Với $a$, $b$ là hai số thực dương tùy ý, khi đó $\ln (\mathrm{e}a^3b^2)$ bằng 
	 \choice
		{$2\ln a+3\ln  b$}
		{$3\ln a+2\ln  b$}
		{\True $1+3\ln a+2\ln b$} 
		 {$1+6\ln a\cdot \ln b$}
	 \loigiai{
		Ta có  $\ln (\mathrm{e}a^3b^2)=\ln \mathrm{e}+\ln (a^3)+\ln (b^2)=1+3\ln a+2\ln b$.
	} 
\end{ex}

% nguồn: Ok_Toan11_TracNghiem_Lan2
\begin{ex}%[1D6H2-1]%[ma:828bd11f]
	Cho $a$ là số thực dương khác  $1$. Tính $P=\log_a\left(a^4\cdot a^{\tfrac{1}{3}}\right)$.
	\choice 
		{$P=\dfrac{3}{4}$} 
		 {$P=12$}
		{$P=1$} 
		{\True $P=\dfrac{13}{3}$} 
	 \loigiai{
		Ta có  $P=\log_a\left(a^4\cdot a^{\tfrac{1}{3}}\right)=\log_a\left(a^{\tfrac{13}{3}}\right)=\dfrac{13}{3}$.
	 }
\end{ex}

% nguồn: Ok_TracNghiemLop11_Lan1
\begin{ex}%[1D6H2-1]%[ma:09519010]
	Cho $a$  là số thực dương khác $1$. Giá trị của $\log _{a^3} a^2$ bằng
	\choice 
		 {\True $\dfrac{2}{3}$}
		{$-\dfrac{3}{2}$} 
		{$-\dfrac{2}{3}$} 
		 {$\dfrac{3}{2}$}
	\loigiai{ 
		Ta có $\log _{a^3} a^2 =  \dfrac{1}{3} \log _{a} a^2 = \dfrac{2}{3} \log _{a} a = \dfrac{2}{3}$.
	 }
\end{ex}

% nguồn: Ok_TracNghiemLop11_Lan1
\begin{ex}%[1D6H2-1]%[ma:0a7e7350]
	Chọn mệnh đề \textbf{sai}? 
	 \choice
		{\True $1+\log_2 2< \log_3 3$} 
		 {$\log_2 \sqrt{2} < \log_2 2$}
		{$\log_2 8> \log_2  4$}
		 {$\log_2 \dfrac{1}{2} > \log_2 \dfrac{1}{4}$}
	\loigiai{ 
		Vì\\ $\left. \begin{array}{l}1+\log_2 2=1+1=2 \\ \log_3 3=1
\end{array}\right\}\Rightarrow
1+\log_2 2> \log_3 3$ nên $1+\log_2 2< \log_3 3$ là khẳng  định sai.\\
Vì\\ $\left. \begin{array}{l} \log_2 \sqrt{2}=\log_2 2^{\tfrac{1}{2}}=\dfrac{1}{2} \\
\log_2 2=1 
\end{array}\right\}\Rightarrow \log_2 \sqrt{2} < \log_2 2$ nên $\log_2 \sqrt{2} < \log_2 2$ là khẳng định đúng.\\
Vì\\ $\left. \begin{array}{l} \log_2 8=\log_2 2^3=3 \\
\log_2 4=\log_2 2^2=2 
\end{array}\right\}\Rightarrow \log_2 8> \log_2 4$ nên $\log_2 8> \log_2 4$ là khẳng định đúng.\\
Vì\\ $\left. \begin{array}{l} \log_2 \dfrac{1}{2}=\log_2 2^{-1}=-1 \\
\log_2 \dfrac{1}{4}=\log_2  2^{-2}=-2
\end{array}\right\}\Rightarrow \log_2 \dfrac{1}{2} > \log_2 \dfrac{1}{4}$ nên $\log_2 \dfrac{1}{2} > \log_2 \dfrac{1}{4}$ là khẳng định đúng.
 }
\end{ex}

% nguồn: Ok_TracNghiemLop11_Lan1
\begin{ex}%[1D6H2-1]%[ma:bf7b0948]
	Cho $a$ là số thực dương  khác $4$. Tính $I=\log _{\tfrac{a}{4}} \left(\dfrac{a^3}{64} \right)$.
	\choice 
		{\True  $I=3$}
		{$I=\dfrac{1}{3}$} 
		 {$I=-3$}
		{$I=-\dfrac{1}{3}$} 
	\loigiai{ 
		Ta  có $I=\log _{\tfrac{a}{4}} \left(\dfrac{a^3}{64} \right) = \log _{\tfrac{a}{4}} \left(\dfrac{a}{4} \right)^3 = 3$.
	 }
\end{ex}

% nguồn: Ok_TracNghiemLop11_Lan1
\begin{ex}%[1D6H2-1]%[ma:7e0bc04c]
	 Cho $\log _{30} 3=a$, $\log _{30} 5=b$. Tính $\log _{30} 1350$ theo $a,b$; $\log _{30} 1350$ bằng
	 \choice
		 {$2a+b$}
		 {\True $2a+b+1$}
		 {$2a+b-1$}
		{$2a+b-2$} 
	\loigiai{ 
		Ta có  $1350=30\cdot45=30\cdot9\cdot5=30\cdot3^2\cdot5$\\
		Nên $\begin{aligned}
		\log _{30} 1350 &= \log  _{30} 30\cdot3^2\cdot5\\
		& = \log _{30} 30+\log _{30} 3^2+\log _{30} 5\\ 
		&  = 1+2\log _{30} 3+\log _{30} 5\\
		&  = 1+2a+b.
	\end{aligned}$
 }
\end{ex}

% nguồn: Ok_TracNghiemLop11_Lan1
\begin{ex}%[1D6H2-1]%[ma:c2502d53]
	Với $a$ là số thực dương tùy ý, $\log_5 a^3$  bằng
	\choice 
		{$\dfrac{1}{3}\log_5  a$}
		{$\dfrac{1}{3}+\log_5  a$}
		{$3+\log_5  a$}
		{\True $3\log_5  a$}
	\loigiai{ 
		 Ta có $\log_5 a^3=3\log_5 a$.
	} 
\end{ex}

% nguồn: Ok_TracNghiemLop11_Lan1
\begin{ex}%[1D6H2-1]%[ma:51d7933a]
	Cho $a$ là số thực dương tùy ý khác  $1$. Mệnh đề nào dưới đây đúng?
	 \choice
		{$\log_2 a=\log_a 2$} 
		{$\log_2  a=\dfrac{1}{\log_2 a}$}
		 {\True $\log_2 a=\dfrac{1}{\log_a 2}$}
		{$\log_2  a=-\log_a 2$}
	\loigiai{ 
		Áp dụng công thức đổi  cơ số ta có $\log_2 a=\dfrac{1}{\log_a 2}$.
	} 
\end{ex}

% nguồn: Ok_TracNghiemLop11_Lan1
\begin{ex}%[1D6H2-1]%[ma:cecb2195]
	Với $t$ là số thực dương tùy ý,  $\log_9(9t)$ bằng
	 \choice
		{\True $1+\log_9 t$}{$1+t$}{$t$}{$1-\log_9 t$} 
	 \loigiai{
		Ta có $\log_9(9t)=\log_9 9+\log_9 t=1+\log_9 t$. 
	 }
\end{ex}

% nguồn: Ok_TracNghiemLop11_Lan1
\begin{ex}%[1D6H2-1]%[ma:bb49e655]
	Cho $a=\log_2 3$  và $b=\log_2 5$. Biểu diễn $\log 60$ theo $a$, $b$?
	 \choice
		 {\True $\log 60=\dfrac{a+b+2}{1+b}$}
		{$\log 60=\dfrac{a+b+2}{1+a}$} 
		{$\log  60=\dfrac{a+b+1}{1+a}$}
		 {$\log 60=\dfrac{a+b+1}{1+b}$}
	\loigiai{ 
		 Ta có $\log 60=\dfrac{\log_2 60}{\log_2 10}=\dfrac{\log_2 (3\cdot 5\cdot 4)}{\log_2 (2\cdot 5)}=\dfrac{\log_2 3+\log_2 5+2}{l+\log_2 5}=\dfrac{a+b+2}{1+b}$.
	 }
\end{ex}

% nguồn: Ok_TracNghiemLop11_Lan1
\begin{ex}%[1D6H2-1]%[ma:90e67521]
	Cho $0<a\neq 1$, $x>0$, $y>0$. Khẳng định nào sau  đây \textbf{sai}?
	\choice 
		{$\log_a x^{\alpha}=\alpha\log_a  x$}
		{$\log_a(xy)=\log_a x+\log_a  y$}
		{$\log_a \sqrt{x}=\dfrac{1}{2}\log_a x$} 
		{\True  $\log_{\sqrt{a}}x=\dfrac{1}{2}\log_a x$}
	\loigiai{Do  $\log_{\sqrt{a}}x=2\log_a x$.}
\end{ex}

% nguồn: Ok_TracNghiemLop11_Lan1
\begin{ex}%[1D6H2-1]%[ma:ccb1285f]
	 Với các số thực dương $a$, $b$ bất kỳ. Mệnh đề nào đúng?
	\choice 
		 {\True $\log_2 \dfrac{2a^3}{b}=1+3\log_2 a-\log_2 b$}
		{$\log_2 \dfrac{2a^3}{b}=1+\dfrac{1}{3}  \log_2 a-\log_2 b$}
		{$\log_2 \dfrac{2a^3}{b}=1+3\log_2  a+\log_2 b$}
		{$\log_2 \dfrac{2a^3}{b}=1+\dfrac{1}{3}  \log_2 a+\log_2 b$}
	 \loigiai{
		Ta có:  $\log_2 \dfrac{2a^3}{b}=\log_2 2+\log_2 a^3-\log_2 b=1+3\log_2 a-\log_2 b$.
	} 
\end{ex}

% nguồn: Ok_TracNghiemLop11_Lan1
\begin{ex}%[1D6H2-1]%[ma:1991b094]
	Nếu $\log_2 3=a$ thì $\log  _{72} 108$ bằng
	\choice 
		{$\dfrac{2+a}{3+a}$} 
		 {\True $\dfrac{2+3a}{3+2a}$}
		 {$\dfrac{ab}{a+b}$}
		{$\dfrac{2+3a}{2+2a}$} 
	\loigiai{ 
		Ta có $\begin{aligned}
		 \log _{72} 108 &= \dfrac{\log_2 108}{\log_2 72}\\
		& = \dfrac{\log_2 \left(2^2.3^3\right)}{\log_2 \left(2^3.3^2\right)}\\ 
		& = \dfrac{2+3\log_2 3}{3+2\log_2 3}\\ 
		 & = \dfrac{2+3a}{3+2a}.
	\end{aligned}
 $}
\end{ex}

% nguồn: Ok_TracNghiemLop11_Lan1
\begin{ex}%[1D6H2-1]%[ma:87cf5412]
	Giá  trị của biểu thức $4^{\log _{\sqrt{2}} 3}$ bằng
	 \choice
		{$9$} 
		{$\dfrac{1}{3}$} 
		 {$3$}
		{\True $81$} 
	\loigiai{ 
		Ta có $4^{\log _{\sqrt{2}} 3} = \sqrt{2}^{4\log_{\sqrt{2}} 3}=\left(  \sqrt{2}^{\log_{\sqrt{2}} 3}\right)^4 = 3^4 = 81$.
	 }
\end{ex}

% nguồn: Ok_TracNghiemLop11_Lan1
\begin{ex}%[1D6H2-1]%[ma:7b63d337]
	Cho $\log_3  a=2$ và $\log_2 b=\dfrac{1}{2}$. Tính $I=2\log_3 \left[\log_3 (3a)\right]+\log _{\tfrac{1}{4}} b^2$.
	 \choice
		{$I=\dfrac{5}{4}$} 
		{$I=0$} 
		{\True  $I=4$}
		 {$I=\dfrac{3}{2}$}
	 \loigiai{
		 Ta có
		\[I=2\log_3 \left[\log_3 (3a)\right]+\log _{\tfrac{1}{4}} b^2=2\log_3 \left(\log_3 3+\log_3 a\right)+2\log  _{2^{-2}} b=2-\dfrac{1}{2}=\dfrac{3}{2}.\]
	} 
\end{ex}

% nguồn: Ok_TracNghiemLop11_Lan1
\begin{ex}%[1D6H2-1]%[ma:7a88cb48]
	Tính giá trị biểu thức $A=\log_3  18-\log_3 12+\dfrac{1}{2}\log_3 36$.
	 \choice
		 {$0$}
		{$4$} 
		{$\log_3 42$} 
		 {\True$2$}
	 \loigiai{
		$\begin{aligned}
		A&=\log_3  18-\log_3 12+\dfrac{1}{2}\log_3 36\\
		&=\log_3 (9\cdot2)-\log_3(4\cdot3)+\dfrac{1}{2}\log_3  (9\cdot4)\\
		&=2\log_3 3+\log_3  2-\left(2\log_3 2+\log_3 3\right) +\dfrac{1}{2}\left(2\log_3 2+ 2\log_3 3\right)\\
		&=2 
	\end{aligned}$
} 
\end{ex}

% nguồn: Ok_TracNghiemLop11_Lan1
\begin{ex}%[1D6H2-1]%[ma:844cebe3]
	Với các số thực dương $a$, $b$ bất kì. Mệnh đề nào dưới đây đúng? 
	\choice 
		{$\log_2\left(\dfrac{2a^3}{b}\right)=1+3\log_2 a+\log_2  b$}
		{$\log_2\left(\dfrac{2a^3}{b}\right)=1+\dfrac{1}{3}\log_2  a+\log_2 b$}
		 {$\log_2\left(\dfrac{2a^3}{b}\right)=1+3\log_2 a-\log_2 b$}
		{$\log_2\left(\dfrac{2a^3}{b}\right)=1+\dfrac{1}{3}\log_2  a-\log_2 b$}
	 \loigiai{
		Ta có \[\log_2\left(\dfrac{2a^3}{b}\right)=\log_2 (2a^3)-\log_2 (b)=\log_2 2+\log_2 a^3-\log_2 b=1+3\log_2 a+\log_2 b.\] 
	 }
\end{ex}

% nguồn: Ok_TracNghiemLop11_Lan1
\begin{ex}%[1D6H2-1]%[ma:b16efaf1]
	Giá trị  của biểu thức $P=4^{\log_2 \sqrt{3}}$ bằng
	\choice 
		 {$\sqrt{3}$}
		 {\True $3$}
		{$2^{\sqrt{3}}$} 
		{$2\sqrt{3}$} 
	 \loigiai{Ta có $P=4^{\log_2 \sqrt{3}}=\left( 2^2\right)^{\log_2 \sqrt{3}} = \left( 2^{\log_2 \sqrt{3}}\right)^{2}=\left(\sqrt{3}\right)^2=3$.}
\end{ex}

% nguồn: Ok_TracNghiemLop11_Lan1
\begin{ex}%[1D6H2-1]%[ma:2d1e9e7d]
	Cho $a>0$, $a \ne 1$, biểu thức $P=\log_{a^3}a$ có giá trị bằng 
	\choice{$-3$}{\True  $\dfrac{1}{3}$}{$3$}{$-\dfrac{1}{3}$}
	\loigiai{ 
		Ta  có $P=\dfrac{1}{3}\log_aa=\dfrac{1}{3}$.
	 }
\end{ex}

% nguồn: Ok_TracNghiemLop11_Lan1
\begin{ex}%[1D6H2-1]%[ma:68c44e01]
	 Cho $\log_6 3=a$, $\log_6 5=b$. Biểu diễn $\log _{54} 50$ theo $a$, $b$?
	 \choice
		 {\True $\dfrac{2b-a+1}{2a+1}$}
		 {$\dfrac{2a-b+1}{2a+1}$}
		{$\dfrac{2b+1+a}{2a+1}$} 
		{$\dfrac{2a+b+1}{2a+1}$} 
	\loigiai{ 
		Ta có $\log _{54} 50=\dfrac{\log_6 \left(5^2\cdot \dfrac{6}{3} \right)}{\log_6 \left(3^2\cdot 6\right)}=\dfrac{2\log_6 5+\log_6 6-\log_6 3}{2\log_6 3+\log_6  6}=\dfrac{2b+1-a}{2a+1}$.
	 }
\end{ex}

% nguồn: Ok_TracNghiemLop11_Lan1
\begin{ex}%[1D6H2-1]%[ma:0c1b1efe]
	Cho $a, b$  là các số thực dương thỏa mãn $a\ne 1$, $a\ne \sqrt{b}$ và $\log_ab=\sqrt{3}$. Tính $P=\log_{\tfrac{\sqrt{b}}{a}}\sqrt{\dfrac{b}{a}}$.
	\choice 
		 {$P=-5+3\sqrt{3}$}
		 {$P=-1+\sqrt{3}$}
		{\True \True $P=-1-\sqrt{3}$} 
		{$P=-5-3\sqrt{3}$} 
	 \loigiai{
		\begin{itemize}
			 \item \textbf{Cách 1:} Phương pháp tự luận.
			 $$P=\dfrac{\log_a\sqrt{\dfrac{b}{a}}}{\log_a\dfrac{\sqrt{b}}{a}}=\dfrac{\dfrac{1}{2}\left(\log_ab-1\right)}{\log_a\sqrt{b}-1}=\dfrac{\dfrac{1}{2}\left(\sqrt{3}-1\right)}{\dfrac{1}{2}\log_ab-1}=\dfrac{\sqrt{3}-1}{\sqrt{3}-2}=-1-\sqrt{3}.$$
			\item \textbf{Cách 2:} Phương  pháp trắc nghiệm. \\
			 Chọn $a=2$, $b=2^{\sqrt{3}}$. Bấm máy tính ta được $P=-1-\sqrt{3}$.
		\end{itemize}
	} 
\end{ex}

% nguồn: Ok_TracNghiemLop11_Lan1
\begin{ex}%[1D6H2-1]%[ma:6c5405dd]
	Với $a$  là số thực dương tùy ý, $\log_7 (7a)$ bằng
	\choice 
		 {$1-\log_7 a$}
		{\True  $1+\log_7 a$}
		{$1+a$} 
		{$a$} 
	\loigiai{ 
		Ta có  $\log_7 (7a)=\log_7 7+\log_7 a=1+\log_7 a$.
	} 
\end{ex}

% nguồn: Ok_TracNghiemLop11_Lan1
\begin{ex}%[1D6H2-1]%[ma:e1c80ac3]
	Với $a$, $b$ là các số thực dương tùy ý thỏa mãn $a\neq 1$ và $\log_a  b=2$, giá trị của $\log_{a^2}(ab^2)$ bằng
	\choice 
		 {$2$}
		{$\dfrac{3}{2}$} 
		 {$\dfrac{1}{2}$}
		 {\True $\dfrac{5}{2}$}
	 \loigiai{
		Ta có $\log_{a^2}(ab^2)=\dfrac{\log_a (ab^2)}{\log_a a^2}=\dfrac{1+\log_a b^2}{2}=\dfrac{1+2\log_a b}{2}=\dfrac{1+2\cdot  2}{2}=\dfrac{5}{2}$.
	 }
\end{ex}

% nguồn: Ok_TracNghiemLop11_Lan1
\begin{ex}%[1D6H2-1]%[ma:004aa1d0]
	Tính $P=\log  _{a^2} \left(a^{10} b^2\right)+\log _{\sqrt{a}} \left(\dfrac{a}{\sqrt{b}} \right)+\log _{\sqrt[3]b} b^{-2}$
	với $0< a\ne 1$ và  $0< b\ne 1$.
	 \choice
		{$P=2$} 
		{\True $P=1$} 
		{$P=\sqrt{3}$} 
		{$P=\sqrt{3}$} 
	 \loigiai{
		Ta  có:
		\begin{eqnarray*}
			P&=&\log _{a^2} \left(a^{10} b^2\right)+\log _{\sqrt{a}} \left(\dfrac{a}{\sqrt{b}} \right)+\log  _{\sqrt[3]b} b^{-2}\\
			&=&\left(5+\log_a b\right)+2\log_a \left(\dfrac{a}{\sqrt{b}}  \right)-6\\
			 &=&5+\log_a b+2-\log_a b-6=1.
		\end{eqnarray*}
	} 
\end{ex}

% nguồn: Ok_TracNghiemLop11_Lan1
\begin{ex}%[1D6H2-1]%[ma:c7c2febc]
	Với các số thực  dương $x$, $y$ tùy ý, đặt $\log_3 x= \alpha$, $\log_3 y=\beta$. Mệnh đề nào dưới đây đúng?
	 \choice
		 {$\log_{27}\left(\dfrac{\sqrt{x}}{y}\right)^3=\dfrac{\alpha}{2}+\beta$}
		{$\log_{27}\left(\dfrac{\sqrt{x}}{y}\right)^3=9\left(\dfrac{\alpha}{2}+\beta\right)$} 
		{\True  $\log_{27}\left(\dfrac{\sqrt{x}}{y}\right)^3=\dfrac{\alpha}{2}-\beta$}
		{  $\log_{27}\left(\dfrac{\sqrt{x}}{y}\right)^3=9\left(\dfrac{\alpha}{2}-\beta\right)$}
	\loigiai{ 
		 Ta có \[\log_{27}\left(\dfrac{\sqrt{x}}{y}\right)^3=\dfrac{3}{2}\log_{27}x-3\log_{27}y=\dfrac{1}{2}\log_3 x-\log_3 y=\dfrac{\alpha}{2}-\beta.\]
	} 
\end{ex}

% nguồn: Ok_TracNghiemLop11_Lan1
\begin{ex}%[1D6H2-1]%[ma:c2215389]
	Cho $a$, $b$ là các số thực dương thỏa mãn  $a\neq 1$, $a\neq \sqrt{b}$ và $\log_a b=\sqrt{3}$. Tính $P=\log_{\frac{\sqrt{b}}{a}}\sqrt{\dfrac{b}{a}}$.
	\choice 
		 {$P=-5+3\sqrt{3}$}
		 {$P=-1+\sqrt{3}$}
		{\True $P=-1-\sqrt{3}$} 
		{$P=-5-3\sqrt{3}$} 
	 \loigiai{
		Ta có \begin{eqnarray*}
		P&=&\dfrac{\log_a  \sqrt{\dfrac{b}{a}}}{\log_a \dfrac{\sqrt{b}}{a}}\\
		&=&\dfrac{\dfrac{1}{2}\log_a \dfrac{b}{a}}{\log_a  \sqrt{b}-\log_a a}\\
		 &=&\dfrac{\dfrac{1}{2}\left(\log_a b-\log_a a\right)}{\dfrac{1}{2}\log_a b-1}\\
		 &=&\dfrac{\dfrac{1}{2}\left(\sqrt{3}-1\right)}{\dfrac{1}{2}\cdot\sqrt{3}-1}\\
		&=&\dfrac{\sqrt{3}-1}{\sqrt{3}-2}\\&=&  -1-\sqrt{3}.
	\end{eqnarray*}
} 
\end{ex}

% nguồn: Ok_TracNghiemLop11_Lan1
\begin{ex}%[1D6H2-1]%[ma:6ef82fc5]
	Giá trị của biểu thức $M=\log_2 2+\log_2 4+\log_2  8+\cdots+\log_2 256$ bằng
	\choice 
		{$48$} 
		 {$56$}
		 {\True $36$}
		{$8\log_2 256$} 
	\loigiai{ 
		Ta có: $\begin{aligned}
		M&=\log_2  2+\log_2 4+\log_2 8+\cdots+\log_2 256\\
		& =\log_2\left(2^1\cdot2^2\cdot2^3\cdots 2^8\right)\\ 
		&=\log_2\left(2^{1+2+3+\cdots+8}  \right)\\
		 &=\left(1+2+3+\cdots+8\right)\cdot\log_22\\
		 &=1+2+3+\cdots+8=36.
	\end{aligned}$
	 }
\end{ex}

% nguồn: Ok_Toan11_TracNghiem_Lan2
\begin{ex}%[1D6V2-1]%[ma:03911a8e]
	Cho $a, b$ là các số thực dương thỏa mãn $\log_{a}b=2$. Tính giá trị của biểu  thức $P=\log_{\tfrac{\sqrt{a}}{b}}\left(a\cdot\sqrt[3]{b}\right)$.
	\choice 
		{$P=\dfrac{2}{15}$} 
		 {$P=-\dfrac{2}{9}$}
		{\True  $P=-\dfrac{10}{9}$}
		 {$P=\dfrac{2}{3}$}
	\loigiai{ 
		 Ta có
		\begin{eqnarray*}
			 P&=&\log_{\tfrac{\sqrt{a}}{b}}\left(a\sqrt[3]{b}\right)=\dfrac{\log_{a}\left(a\cdot\sqrt[3]{b}\right)}{\log_{a}\left(\dfrac{\sqrt{a}}{b}\right)}=\dfrac{\log_{a}a+\log_{a}b^{\tfrac{1}{3}}}{\log_{a}a^{\tfrac{1}{2}}-\log_{a}b}=\dfrac{1+\dfrac{1}{3}\log_{a}b}{\dfrac{1}{2}-\log_{a}b}\\
			&=&\dfrac{1+\dfrac{1}{3} \cdot  2}{\dfrac{1}{2}-2}=-\dfrac{10}{9}.
		\end{eqnarray*}
	} 
\end{ex}

% nguồn: Ok_Toan11_TracNghiem_Lan2
\begin{ex}%[1D6V2-1]%[ma:85eec760]
	Cho $x$, $y$ là hai số thực dương, $x\ne 1$ thỏa mãn $\log_{\sqrt{x}}y=\dfrac{2y}{5}$, $\log_{25}x=\dfrac{5}{2y}$. Tính giá trị  của $P=y^2-2x^2$.
	\choice 
		{$P=1$} 
		 {$P=0$}
		{\True $P=-25$} 
		{$P=25$} 
	 \loigiai{
		Ta  có
		 \allowdisplaybreaks
		\begin{eqnarray*}
			\heva{&\log_{\sqrt x}y=\dfrac{2y}{5}\\&\log_{25}x=\dfrac{5}{2y}}&\Leftrightarrow&  \heva{&\log_xy^2=\dfrac{2y}{5}\\&\log_x25=\dfrac{2y}{5}}\Leftrightarrow \heva{&y^2=25\\&\log_{25}x=\dfrac{5}{2y}}\\
			 &\Leftrightarrow&\heva{&y=5\\&\log_{25}x=\dfrac{1}{2}}\Leftrightarrow \heva{&y=5\\&x=5.}
		\end{eqnarray*}
		Vậy  $P=y^2-2x^2=-25$.
	} 
\end{ex}

% nguồn: Ok_Toan11_TracNghiem_Lan2
\begin{ex}%[1D6V2-1]%[ma:94c0b2a9]
	Cho $x$,  $y$ là hai số thực dương, $x \ne 1$ thỏa mãn
	$\log_{\sqrt{x}} y = \dfrac{2y}{5}$, $\log_{25} x = \dfrac{5}{2y}$. Giá trị của $P =  y^2 - 2x^2$ là
	\choice 
		{ $P  = 1$ }
		{ $P = 0$  }
		{  \True $P = -25$ }
		{ $P =  25$ }
	\loigiai{ 
		Điều kiện  $0<x\ne 1$, $y>0$.\\
		Ta có 
		 $$
		\heva{ 
		&  \log_{\sqrt{x}} y = \dfrac{2y}{5} \\
		 & \log_{25} x = \dfrac{5}{2y}
		 }
		\Leftrightarrow\heva{ 
		 & \log_x y^2 = \dfrac{2y}{5} \\
		& \log_x 25 =  \dfrac{2y}{5}
		 }
		 \Leftrightarrow\heva{
		 & y^2 = 25 \\
		& \log_{25}  x = \dfrac{5}{2y}
		} 
		 \Leftrightarrow\heva{
		&  y = 5 \\
		&  \log_{25} x = \dfrac{1}{2}
		 }
		\Leftrightarrow\heva{ 
		 & y = 5 \\
		& x = 5. 
		} 
		$$ 
		 Vậy $P = y^2 - 2x^2 = -25$.
	 }
\end{ex}

% nguồn: Ok_Toan11_TracNghiem_Lan2
\begin{ex}%[1D6V2-1]%[ma:79628995]
	Cho các số thực $a$, $b$ thỏa mãn $a > b > 1$ và $\dfrac{1}{\log_b a} + \dfrac{1}{\log_a  b} = \sqrt{2020}$. Giá trị của biểu thức
	$P = \dfrac{1}{\log_{ab} b} - \dfrac{1}{\log_{ab} a}$ là 
	\choice 
		{ $\sqrt{2014}$  }
		{ $\sqrt{2016}$  }
		{ $\sqrt{2018}$ } 
		{ \True $\sqrt{2020}$ } 
	\loigiai{ 
		Do $a > b > 1$ nên $\log_a b  > 0$, $\log_b a > 0$ và $\log_b a > \log_a b$.\\
		Ta có 
		\begin{eqnarray*}
			&&\dfrac{1}{\log_b  a} + \dfrac{1}{\log_a b} = \sqrt{2020}\\
			&\Leftrightarrow& \log_a b  + \log_b a = \sqrt{2020}\\
			&\Leftrightarrow& \log_a^2 b + \log_b^2 a + 2 =  2020\\
			&\Leftrightarrow& \log_a^2 b + \log_b^2 a = 2018\quad  (*)
		\end{eqnarray*}
		Khi đó 
		 $$
		P = \log_b  ab - \log_a ab
		= \log_b a + \log_b  b - (\log_a a + \log_a b)
		= \log_b a - \log_a  b.
		$$ 
		Suy  ra
		$$ 
		P^2  = (\log_b a - \log_a b)^2
		= \log_a^2 b  + \log_b^2 a - 2
		 = 2018 - 2 = 2016.
		 $$
		 Vậy $P = \sqrt{2016}$.
	} 
\end{ex}

% nguồn: Ok_Toan11_TracNghiem_Lan2
\begin{ex}%[1D6V2-1]%[ma:e74cc181]
	Cho $x$, $y$ là hai số thực dương, $x\ne 1$ thỏa mãn $\log_{\sqrt[3]{x}} y = \dfrac{3y}{8}$, $\log_{\sqrt{2}} x = \dfrac{32}{y}$. Giá trị của $P  = x^2 - y^2$ là
	\choice 
		{ $P =  120$ }
		{  $P = 132$ }
		{ \True $P =  240$ }
		 { $P = 340$ }
	 \loigiai{
		 Điều kiện $0<x\ne 1$, $y>0$.\\
		Ta có 
		 $$
		\log_{\sqrt{2}}  x = \dfrac{32}{y}
		 \Leftrightarrow
		2\log_2 x  = \dfrac{32}{y}
		\Leftrightarrow 
		y  = 16\log_x 2 \quad(*)
		$$ 
		 Mặt khác, ta có
		\begin{eqnarray*}
			&&\log_{\sqrt[3]{x}} y  = \dfrac{3y}{8}\\
			&\Leftrightarrow& 3\log_x y =  \dfrac{3y}{8}\\
			&\Leftrightarrow&  \log_x y = \dfrac{y}{8}\\
			 &\Leftrightarrow& \log_x (16\log_x 2) = 2\log_x 2\\
			&\Leftrightarrow& 16\log_x 2 =  4\\
			 &\Leftrightarrow& \log_2 x = 4\\
			&\Leftrightarrow& x =  16.
		\end{eqnarray*}
		Thay $x=16$  vào (*), ta được $y = 4$.\\
		Vậy $P = x^2 - y^2 = 16^2 -  4^2 = 240$.
	 }
\end{ex}

% nguồn: Ok_Toan11_TracNghiem_Lan2
\begin{ex}%[1D6V2-1]%[ma:18c35e97]
	Cho $a$, $b$ là các số  thực dương lớn hơn $1$ thỏa mãn $\log_ab=2$. Tính giá trị biểu thức $P=\log_{a^2}b+\log_{ab^2}b^5$.
	 \choice
		 {\True $P=3$}
		{$P=4$} 
		 {$P=2$}
		{$P=5$} 
	\loigiai{ 
		Ta có \begin{align*}
		P&=\log_{a^2}b+\log_{ab^2}b^5\\ 
		&=\log_{a^2}b+\dfrac{1}{\log_{b^5}\left(ab^2\right)}\\ 
		&=\dfrac{1}{2}\log_ab+\dfrac{1}{\log_{b^5}a+\log_{b^5}b^2}\\ 
		&=\dfrac{1}{2}\log_ab+\dfrac{1}{\tfrac{1}{5}\log_ba+\tfrac{2}{5}}\\ 
		&=\dfrac{1}{2}\cdot 2+\dfrac{1}{\tfrac{1}{5}\cdot\tfrac{1}{2}+\tfrac{2}{5}}=1+2=3 
	\end{align*}
 }
\end{ex}

% nguồn: Ok_Toan11_TracNghiem_Lan2
\begin{ex}%[1D6V2-1]%[ma:89b812f1]
	Cho các số thực $a$,  $b$ thỏa mãn $a>b>1$và $\dfrac{1}{\log_ba}+\dfrac{1}{\log_ab}=\sqrt{2\,022}$. Giá trị của biểu thức $P=\dfrac{1}{\log_{ab}b}-\dfrac{1}{\log_{ab}a}$ bằng
	\choice 
		{\True  $\sqrt{2\,018}$}
		 {$\sqrt{2\,020}$}
		{$\sqrt{2\,016}$} 
		 {$\sqrt{2\,022}$}
	 \loigiai{
		Ta  có
		 \[\dfrac{1}{\log_ba}+\dfrac{1}{\log_ab}=\sqrt{2\,022}\Leftrightarrow \log_ab+\log_ba=\sqrt{2\,022}.\quad\quad \hfill (*)\]
		Khi đó 
		 \[P=\dfrac{1}{\log_{ab}b}-\dfrac{1}{\log_{ab}a}=\log_b(ab)-\log_a(ab)=\log_ba-\log_ab.\]
		Đặt  $t=\log_ab$ thì (*) trở thành
		\allowdisplaybreaks 
		\begin{eqnarray*}
			 t+\dfrac{1}{t}=\sqrt{2\,022}&\Leftrightarrow&t^2-\sqrt{2\,022}t+1=0\\
			&\Leftrightarrow&\hoac{& 
			t=\dfrac{\sqrt{2\,022}+\sqrt{2\,018}}{2}\Rightarrow  P=\dfrac{1}{t}-t=-\sqrt{2\,018}\\
			 &t=\dfrac{\sqrt{2\,022}-\sqrt{2\,018}}{2}\Rightarrow P=\dfrac{1}{t}-t=\sqrt{2\,018}.}
		\end{eqnarray*}
		 Mặt khác, do $a>b>1\Leftrightarrow 0<\log_ab<1$ nên $0<t<1\Rightarrow\dfrac{1}{t}>1\Rightarrow P=\dfrac{1}{t}-t>0$.\\
		Vậy  $P=\sqrt{2\,018}$.
	} 
\end{ex}

% nguồn: Ok_Toan11_TracNghiem_Lan2
\begin{ex}%[1D6V2-1]%[ma:bb869252]
	Biết $x$ và $y$ là hai số thực thỏa mãn $\log _4 x=\log _{9}  y=\log _6(x-2 y)$. Giá trị của $\dfrac{x}{y}$ bằng
	\choice 
		 {$\log _{\tfrac{2}{3}}^2 2$}
		{$1$} 
		 {\True $4$}
		{$2$} 
	\loigiai{Điều kiện $\heva{&x>0 \\&  y>0 \\& x>2 y.}$\\
		Đặt  $\log _4 x=\log _9 y=\log _6(x-2 y)=t$. Ta có
		 \allowdisplaybreaks
		\begin{eqnarray*}
			\heva{& 
			x=4^t \\& 
			 y=9^t \\&
			x-2 y=6^t}  &\Rightarrow& 4^t-2\cdot 9^t=6^t \\&\Leftrightarrow&\left(\dfrac{4}{9}\right)^t-\left(\dfrac{2}{3}\right)^t-2=0 \\&\Leftrightarrow&\heva{&
			 \left(\dfrac{2}{3}\right)^t=-1& \text { (loại) } \\&
			\left(\dfrac{2}{3}\right)^t=2 
			 &\text { (nhận) }.}
		\end{eqnarray*}
		 Khi đó $\dfrac{x}{y}=\left(\dfrac{4}{9}\right)^t=\left[\left(\dfrac{2}{3}\right)^t\right]^2=4$.
	} 
\end{ex}

% nguồn: Ok_Toan11_TracNghiem_Lan2
\begin{ex}%[1D6V2-1]%[ma:76f89383]
	Có bao nhiêu số thực dương $n \ne 1$ để $\log_{n}265$ là một  số nguyên?
	 \choice
		 {$2$}
		{$4$} 
		{$6$} 
		 {\True $8$}
	 \loigiai{
		\begin{itemize}
			\item Ta có $\log_{n}256=\log_{n}2^{8}=8\log_{n}2=\dfrac{8}{\log_{2}n}$. 
			\item Để $\log_{n}265$ là một số nguyên thì $\log_{2}n\in\left\{\pm1;\pm2;\pm4;\pm8\right\}\Leftrightarrow  n\in\left\{\dfrac{1}{2};2;\dfrac{1}{4};4;\dfrac{1}{16};16;\dfrac{1}{256};256\right\}$.
			\item Vậy có tất cả $8$ số thực dương $n\ne1$ thỏa  mãn điều kiện bài toán.
		\end{itemize}
	} 
\end{ex}

% nguồn: Ok_Toan11_TracNghiem_Lan2
\begin{ex}%[1D6V2-1]%[ma:008a2a2a]
	Đặt $a=\log _2 3$, $b=\log _5 3$. Nếu  biểu diễn $\log _6 45=\dfrac{a(m+n b)}{b(a+p)}$ thì $m+n+p$ bằng
	 \choice
		 {$-3$}
		{$3$} 
		 {$6$}
		 {\True $4$}
	 \loigiai{
		Ta  có
		 \allowdisplaybreaks
		\begin{eqnarray*}
			\log _6 45&=&\dfrac{\log  _3 45}{\log _3 6}=\dfrac{\log _3 9+\log _3 5}{\log _3 2+\log _3 3} \\
			&=&\dfrac{\log _3 3^2+\dfrac{1}{\log _5 3}}{\log _3 3+\dfrac{1}{\log _2 3}}=\dfrac{2+\dfrac{1}{b}}{1+\dfrac{1}{a}}=\dfrac{a(1+2 b)}{b(a+1)}  .
		\end{eqnarray*}
		Theo bài  ra $\dfrac{a(m+n b)}{b(a+p)}=\dfrac{a(1+2 b)}{b(a+1)}\Rightarrow m=1$, $n=2$, $p=1$.\\
		 Vậy $m+n+p=4$.
	 }
\end{ex}

% nguồn: Ok_Toan11_TracNghiem_Lan2
\begin{ex}%[1D6V2-1]%[ma:df4f6521]
	Cho $a,b$ là các số thực dương lớn hơn 1 thỏa mãn $\log_{a}b=3$. Tính  giá trị biểu thức $P=\log_{a^{2}b}a^{3}-3\log_{a^{2}}2 \cdot \log_{4}\left(\dfrac{a}{b}\right)$.
	\choice 
		{$P=\dfrac{15}{8}$} 
		{$P=\dfrac{18}{25}$} 
		{\True $P=\dfrac{21}{10}$} 
		 {$P=\dfrac{7}{5}$}
	\loigiai{ 
		\begin{itemize}
			\item  Ta có $\log_{a}b=3\Leftrightarrow b=a^{3}$.
			\item  Suy ra
			\begin{eqnarray*}
				P&=&\log_{a^{2}b}a^{3}-3\log_{a^{2}}2 \cdot \log_{4}\left(\dfrac{a}{b}\right)\\ &=&\log_{a^{5}}a^{3}-3\log_{a^{2}}2 \cdot \log_{2^{2}}\left(\dfrac{a}{a^{3}}\right)\\ &=&\dfrac{3}{5}-3 \cdot \dfrac{1}{2}\log_{a}2 \cdot \dfrac{1}{2}\log_{2}a^{-2}\\  &=&\dfrac{3}{5}-3 \cdot \dfrac{1}{2}\log_{a}2 \cdot \dfrac{1}{2} \cdot \left(-2\right)\log_{2}a\\ &=&\dfrac{3}{5}+\dfrac{3}{2}\log_{a}2 \cdot \log_{2}a=\dfrac{3}{5}+\dfrac{3}{2}=\dfrac{21}{10}.
			\end{eqnarray*}
		\end{itemize}
	 }
\end{ex}

% nguồn: Ok_Toan11_TracNghiem_Lan2
\begin{ex}%[1D6V2-1]%[ma:f42edc6e]
	 Cho hàm số
	$f(x) = \dfrac12 \log_2 \left(\dfrac{2x}{1-x}\right)$ và hai số thực $m$, $n$ thuộc  $(0;1)$ sao cho $m+n=1$. Giá trị của $f(m)+f(n)$ bằng
	\choice 
		 { $2$ }
		{ $0$  }
		{ \True $1$ } 
		{ $\dfrac12$  }
	\loigiai{ 
		Ta  có
		\begin{eqnarray*}
			 f(m)+f(n) &=& \dfrac12 \log_2\left(\dfrac{2m}{1-m}\right) + \dfrac12 \log_2
			\left(\dfrac{2n}{1-n}\right)  \\
			 &=& \dfrac{1}{2} \log_2
			 \left(\dfrac{2m}{1-m} \cdot \dfrac{2n}{1-n}\right) \\
			&=& \dfrac12  \log_2\left(\dfrac{4mn}{1-m-n+mn}\right).
		\end{eqnarray*}
		Vì $m+n=1$ nên 
		$$ 
		f(m)+f(n) 
		=  \dfrac12 \log_2 \left(\dfrac{4mn}{mn}\right)
		= \dfrac12  \log_2 4
		= 1. 
		 $$
	} 
\end{ex}

% nguồn: Ok_Toan11_TracNghiem_Lan2
\begin{ex}%[1D6V2-1]%[ma:2e236ccd]
	Có bao nhiêu số nguyên dương $n$ để  $\log_n 256$ là một số nguyên dương?
	 \choice
		{ $2$  }
		{  $3$ }
		 { \True $4$ }
		 { $1$ }
	\loigiai{ 
		Ta  có
		$ 
		 \log_n 256 = 8\log_n 2 = \dfrac{8}{\log_2 n}
		$ là  số nguyên dương. \\
		Khi  đó ta có
		 $$
		\log_2 n \in \{1,2,4,8\} 
		 \Rightarrow n \in \{2,4,16,256\}.
		 $$
		Vậy có $4$ số  nguyên dương.
	} 
\end{ex}

% nguồn: Ok_Toan11_TracNghiem_Lan2
\begin{ex}%[1D6V2-1]%[ma:8136656a]
	Cho $x$, $y$ là  hai số thực dương, $x\ne 1$ thỏa mãn $\log_{\sqrt{x}} y = \dfrac{2y}{4}$, $\log_{16} x = \dfrac{4}{2y}$. Giá trị của $P = y^2 - 2x^2$ là
	\choice 
		{ $P  = 1$ }
		{ $P = 0$ } 
		{ \True  $P = -16$ }
		{ $P = 16$ } 
	\loigiai{ 
		Điều kiện $0<x\ne 1$, $y>0$.\\ 
		Ta có 
		 $$
		\heva{ 
		& \log_{\sqrt{x}} y  = \dfrac{2y}{4} \\
		& \log_{16}  x = \dfrac{4}{2y}
		 }
		\Leftrightarrow 
		\heva{ 
		& \log_x y^2 =  \dfrac{2y}{4} \\
		&  \log_x 16 = \dfrac{2y}{4}
		 }
		 \Leftrightarrow
		 \heva{
		& y^2 = 16  \\
		& \log_{16} x =  \dfrac{4}{2y}
		 }
		\Leftrightarrow 
		\heva{ 
		& y =  4 \\
		& \log_{16} x  = \dfrac{1}{2}
		} 
		 \Leftrightarrow
		 \heva{
		 & y = 4 \\
		 & x = 4.
		} 
		 $$
		 Vậy $P = y^2 - 2x^2 = -16$.
	 }
\end{ex}

% nguồn: Ok_Toan11_TracNghiem_Lan2
\begin{ex}%[1D6V2-1]%[ma:a507b7f1]
	 Cho các số thực dương $x$, $y$, $z$ thỏa mãn đồng thời $\dfrac{1}{\log_2 x} + \dfrac{1}{\log_2 y} + \dfrac{1}{\log_2 z} = \dfrac{1}{2020}$ và
	$\log_2 (xyz) = 2020$. Giá trị của $\log_2\left[xyz(x+y+z) - xy - yz  - zx +1\right]$ bằng
	\choice 
		{ \True  $4040$ }
		{  $1010$ }
		{ $2020$ } 
		{ $2020^2$  }
	\loigiai{ 
		Đặt  $a = \log_2 x$, $b = \log_2 y$, $c = \log_2 z$.\\
		 Ta có $\dfrac{1}{a} + \dfrac{1}{b} + \dfrac{1}{c} = \dfrac{1}{2020}$, $a + b + c = 2020$.\\
		Khi  đó, ta có
		\begin{eqnarray*}
			&&\left(\dfrac{1}{a} + \dfrac{1}{b} + \dfrac{1}{c}\right)\cdot(a+b+c) = 1\\ 
			 &\Leftrightarrow& (a+b+c)\cdot(ab+ac+bc) = abc\\
			&\Leftrightarrow& a^2 b +  ab^2 + abc + b^2 c + bc^2 + ac^2 = 0\\
			&\Leftrightarrow& (a+b)\cdot(b+c)\cdot(c+a)  = 0.
		\end{eqnarray*}
		Vì vai trò $a$, $b$, $c$ như nhau nên  giả sử $a + b = 0$.\\
		Suy  ra
		 $
		c = 2020 \Rightarrow z = 2^{2020}$  và $xy = 1$.\\
		 Do đó
		\begin{eqnarray*}
			\log_2\left[xyz(x+y+z) - xy - yz - zx +1\right] &=& \log_2\left[xyz(x+y+z)  - xy - yz - zx +1\right]\\
			 &=& \log_2(z^2)\\
			 & =& 2\log_2 z \\
			&=& 4040. 
		\end{eqnarray*}
	} 
\end{ex}

% nguồn: Ok_Toan11_TracNghiem_Lan2
\begin{ex}%[1D6V2-1]%[ma:d0ea8a71]
	Cho $x$, $y$, $z$ là các  số thực lớn hơn $1$ và $w$ là số thực dương sao cho $\log_x w = 24$, $\log_y w = 40$, $\log_{xyz} w = 12$. Giá trị của $\log_z w$ bằng
	\choice 
		{ $52$  }
		 { $-60$ }
		{ \True $60$ } 
		 { $-52$ }
	 \loigiai{
		Ta  có $\log_x w = 24 \Rightarrow \log_w x = \dfrac{1}{24}$.\\
		Và  $\log_y w = 40 \Rightarrow \log_w y = \dfrac{1}{40}$.\\
		Khi đó, ta  có
		\begin{eqnarray*}
			&&\log_{xyz} w =  12\\
			&\Leftrightarrow& \dfrac{1}{\log_w (xyz)} = 12\\ 
			&\Leftrightarrow& \dfrac{1}{\log_w x  + \log_w y + \log_w z} = 12\\
			&\Leftrightarrow&\log_w x + \log_w y + \log_w  z = \dfrac{1}{12}\\
			&\Leftrightarrow& \dfrac{1}{24} + \dfrac{1}{40} + \log_w  z = \dfrac{1}{12}\\
			&\Leftrightarrow&  \log_w z = \dfrac{1}{60}\\
			&\Leftrightarrow& \log_z  w = 60.
		\end{eqnarray*}
		Vậy $\log_z w =  60$.
	} 
\end{ex}

% nguồn: Ok_Toan11_TracNghiem_Lan2
\begin{ex}%[1D6V2-1]%[ma:159513be]
	Cho $a$, $b$, $c$ là các số thực dương,  trong đó $a, b > 1$ và thỏa mãn
	$\log_a c = 3$, $\log_b c = 4$. Giá  trị biểu thức $P = \log_{ab} c$ là
	 \choice
		{ \True $P  = \dfrac{12}{7}$ }
		{ $P = \dfrac{7}{12}$ } 
		 { $P = \dfrac{1}{12}$ }
		{ $P  = 12$ }
	 \loigiai{
		Từ giả thiết ta suy ra $c \ne  1$ và
		$\log_c a = \dfrac{1}{\log_a c} = \dfrac{1}{3}$;  $\log_c b = \dfrac{1}{\log_b c} = \dfrac{1}{4}$.\\
		Khi đó,  ta có
		 $$
		P = \log_{ab}  c
		 = \dfrac{1}{\log_c(ab)}
		= \dfrac{1}{\log_c  a + \log_c b}
		= \dfrac{1}{\dfrac{1}{3} +  \dfrac{1}{4}}
		=  \dfrac{12}{7}.
		 $$
	} 
\end{ex}

% nguồn: Ok_Toan11_TracNghiem_Lan2
\begin{ex}%[1D6V2-1]%[ma:0f0cf2b3]
	Cho $x$, $y$ là hai số thực  dương, $x\ne 1$ thỏa mãn $\log_{\sqrt[3]{x}}y=\dfrac{3y}{8}$, $\log_{\sqrt 2}x=\dfrac{32}{y}$. Tính giá trị của $P=x^2-y^2$.
	\choice 
		{$P=120$} 
		{$P=132$} 
		{\True  $P=240$}
		 {$P=340$}
	\loigiai{ 
		Ta có $\log_{\sqrt  2}x=\dfrac{32}{y}\Leftrightarrow 2\log_2x=\dfrac{32}{y}\Leftrightarrow y=\dfrac{16}{\log_2x}=16\log_x2.$ \hfill (*)\\
		Ta lại có 
		 \allowdisplaybreaks
		\begin{eqnarray*}
			\log_{\sqrt[3]{x}}y=\dfrac{3y}{8}&\Leftrightarrow& 3\log_xy=\dfrac{3y}{8}\Leftrightarrow  \log_xy=\dfrac{y}{8}\\
			 &\Leftrightarrow&\log_x\left(16\log_x2\right)=2\log_x2\\ &\Leftrightarrow&\log_x\left(16\log_x2\right)=\log_x2^2\Leftrightarrow 16\log_x2=4\\
			&\Leftrightarrow&\log_x2=\dfrac{1}{4}\Leftrightarrow  \log_2x=4\Leftrightarrow x=16.
		\end{eqnarray*}
		Từ (*)  suy ra $y=4$.\\
		Vậy $P=x^2-y^2=16^2-4^2=240$. 
	} 
\end{ex}

% nguồn: Ok_Toan11_TracNghiem_Lan2
\begin{ex}%[1D6V2-1]%[ma:6493ceec]
	 Cho $\log 2=a$. Tính $A=\log \dfrac{125}{4}$ theo $a$?
	\choice 
		 {$6+7a$}
		{$2(a+5)$} 
		{$4(1+a)$} 
		{\True  $3-5a$}
	\loigiai{ 
		Ta có \begin{itemize}
		\item $A=\log \dfrac{125}{4}=\log 125-\log 4=\log 5^3-\log  2^2=3\log 5-2\log 2=3\log 5-2a$,
		\item $\log 5=\log \dfrac{10}{2}=\log  10-\log 2=1-a$.
	\end{itemize}
	 Vậy $A=3(1-a)-2a=3-5a$.
} 
\end{ex}

% nguồn: Ok_TracNghiemLop11_Lan1
\begin{ex}%[1D6V2-1]%[ma:7c7ea74c]
	Cho $\log _a x=3, \log _b x=4$ với $a, b$ là các số thực lớn  hơn $1$. Tính $P=\log _{a b}x$.
	 \choice
		{$P=12$} 
		 {\True $P=\dfrac{12}{7}$}
		 {$P=\dfrac{7}{12}$}
		 {$P=\dfrac{1}{12}$}
	 \loigiai{
		$P=\log_{ab}x=\dfrac 1{\log_x a b}=\dfrac 1{\log_x a+\log_x b}=\dfrac 1{\dfrac 13+\dfrac  14}=\dfrac{12}{7}.$
	 }
\end{ex}

% nguồn: Ok_Toan11_TracNghiem_Lan2
\begin{ex}%[1D6C2-1]%[ma:f3874cb9]
	Cho $x$, $y$, $z$ là ba số thực dương lập  thành cấp số nhân
	$\log_a x$, $\log_{\sqrt{a}} y$, $\log_{\sqrt[3]{a}} z$ lập  thành cấp số cộng
	với $a$ là số  thực dương khác $1$. Giá trị của
	$p = \dfrac{9x}{y} + \dfrac{y}{z}  + \dfrac{3z}{x}$ bằng
	 \choice
		 { \True $13$ }
		 { $3$ }
		 { $12$ }
		{ $10$  }
	 \loigiai{
		Điều kiện  $x$>0, $y>0$, $z>0$.\\
		Vì $x$, $y$, $z$ lập  thành cấp số nhân nên
		$xz = y^2$ $(1)$\\ 
		Mặt khác,  dãy
		$\log_a x$, $\log_{\sqrt{a}} y$, $\log_{\sqrt[3]{a}} z$ lập thành cấp  số cộng nên
		$$ 
		\log_a x +  \log_{\sqrt[3]{a}} z = 2\log_{\sqrt{a}} y.
		 $$
		Suy ra 
		\begin{eqnarray*}
			\log_a x + 3\log_a z = 2\cdot 2\log_a  y
			 &\Leftrightarrow&
			\log_a x + 3\log_a z = 4\log_a y\\ 
			&\Leftrightarrow&  \log_a {xz^3}=\log_a y^4\\
			&\Leftrightarrow& xz^3=y^4  \quad(2)
		\end{eqnarray*}
		Từ  $(1)$ và $(2)$ ta suy ra $x = y = z$.\\
		 Vậy
		$$ 
		 p = \dfrac{9x}{y} + \dfrac{y}{z} + \dfrac{3z}{x}
		= 9 + 1 + 3 = 13. 
		 $$
	} 
\end{ex}


%%%%%%%%%%%%%%%% ĐÚNG/SAI %%%%%%%%%%%%%%%%

% nguồn: Ok_dungsai_lop11_Lan1
\begin{ex}%[1D6N2-1]%[ma:61ff5812]
 Cho $a, b$ là hai số thực dương tùy ý và $b \neq  1$. Xác định tính đúng sai của các khẳng định sau.
 \choiceTF 
  {\True  $\ln a+\ln b=\ln (a\cdot b)$}
  {$\ln (a+b)=\ln a \cdot \ln b$} 
   {$\ln a-\ln b=\ln (a-b)$}
  {\True  $\log _b a=\dfrac{\ln a}{\ln b}$}
  \loigiai{
  Theo tính chất của logarit ta có $\ln a+\ln b=\ln (a\cdot b)$, $\log _b a=\dfrac{\ln  a}{\ln b}$ là đúng.
 } 
\end{ex}

% nguồn: Ok_dungsai_lop11_Lan1
\begin{ex}%[1D6N2-1]%[ma:bf7f5348]
  \choiceTF
  {\True Biểu thức $A=\log_6(x+5)$  có nghĩa khi $x\in(-5;+\infty)$}
   {\True Biểu thức $B=\ln(x-1)^2$ có nghĩa khi $x\ne 1$}
  {Biểu thức $C=\log^2\dfrac{1}{x-x^2}$ có nghĩa  khi $x\in (1;+\infty)$}
  {Biểu thức  $D=\log_3(2x+1)$ có nghĩa khi $x\in (0;+\infty)$}
  \loigiai{
  \begin{itemchoice}
   \itemch Đúng. Biểu thức  $A=\log_6(x+5)$ có nghĩa khi $x+5>0\Leftrightarrow x>-5\Leftrightarrow x\in (-5;+\infty)$.
   \itemch  Đúng. Biểu thức $B=\ln(x-1)^2$ có nghĩa khi $(x-1)^2>0\Leftrightarrow x\ne 1$.
   \itemch Sai. Biểu thức $C$ có nghĩa  khi $\dfrac{1}{x-x^2}>0\Leftrightarrow x(1-x)>0\Leftrightarrow 0<x<1$.
   \itemch Sai. Biểu thức  $D$ có nghĩa khi $2x+1>0\Leftrightarrow x>-\dfrac{1}{2}\Leftrightarrow x\in\left(-\dfrac{1}{2};+\infty\right)$.
  \end{itemchoice}
 } 
\end{ex}

% nguồn: Ok_dungsai_lop11_Lan1
\begin{ex}%[1D6H2-1]%[ma:a2a14b8f]
 Cho biểu thức $A=\log_2x^2+\log_{\frac{1}{2}}x^3+\log_4x$. Vậy: 
  \choiceTF
  {\True  Khi $\log_2x=1$ thì $A=-\dfrac{1}{2}$}
   {Khi $\log_2x=2$ thì $A=1$}
  {\True Khi $\log_2x=\sqrt{3}$  thì $A=-\dfrac{\sqrt{3}}{2}$}
  {Khi $\log_2x=\sqrt{2}$ thì $A=\dfrac{\sqrt{2}}{2}$} 
  \loigiai{
  Với điều kiện $x>0$, ta  có
  \begin{align*}
   A=\  &\log_2x^2+\log_{\frac{1}{2}}x^3+\log_4x\\
   =\ &2\log_2x-3\log_2+\dfrac{1}{2}\log_2x\\ 
   =\  &-\dfrac{1}{2}\log_2x.
  \end{align*}
   Từ đó
  \begin{enumerate}
    \item Khi $\log_2x=1$ thì $A=-\dfrac{1}{2}$.
   \item Khi $\log_2x=2$  thì $A=-1$
   \item  Khi $\log_2x=\sqrt{3}$ thì $A=-\dfrac{\sqrt{3}}{2}$
   \item Khi  $\log_2x=\sqrt{2}$ thì $A=-\dfrac{\sqrt{2}}{2}$
  \end{enumerate}
 } 
\end{ex}

% nguồn: Ok_dungsai_lop11_Lan1
\begin{ex}%[1D6H2-1]%[ma:72a79cd9]
 Cho  các biểu thức sau:
  $
 A=\log _{2^{2030}} 4-\dfrac{1}{1015}+\ln e^{2035}$, $B=\log _5 3 \cdot \log _2 5-\dfrac{\ln 9}{\ln  4}
 $ 
  \choiceTF
  {\True $A$ chia  hết cho $5$}
  {$A-B=2036$} 
  {\True  $A+2024B=2035$}
  {\True  $A-2024B=2035$}
  \loigiai{
  Ta có 
   \allowdisplaybreaks
  \begin{eqnarray*}
   A&=&\log _{2^{2030}} 4-\dfrac{1}{1015}+\ln  e^{2035}=\log _{2^{2030}} 2^2-\dfrac{1}{1015}+2035 \\
    &=&\dfrac{2}{2030}-\dfrac{1}{1015}+2035=2035.
  \end{eqnarray*}
  Ta  có
  \allowdisplaybreaks 
  \begin{eqnarray*}
   B&=&\log _5  3 \cdot \log _2 5-\dfrac{\ln 9}{\ln 4}=\log _2 5 \cdot \log _5 3-\log _4 9 \\
    &=&\log _2 3-\log _{2^2} 3^2=\log _2 3-\log _2 3=0.
  \end{eqnarray*}
   Khi đó
  \begin{itemize}
   \item $A=2035$  chia hết cho $5$.
   \item $A-B=2035$. 
    \item $A+2024B=2035$.
   \item $A-2024B=2035$. 
  \end{itemize}
 } 
\end{ex}

% nguồn: Ok_dungsai_lop11_Lan1
\begin{ex}%[1D6H2-1]%[ma:dc652e21]
 Cho các biểu thức sau: $P=\log _2 8+\log _3  27-\log _5 5^3; Q=\ln (2e)-\log 100$. Khi đó
 \choiceTF 
  {$P+Q=2  \ln 2$}
  {\True $Q-P=\ln  2-4$}
   {\True $3 Q+P=3 \ln 2$}
  {\True $2  Q+P=2 \ln 2+1$}
 \loigiai{ 
  Ta  có
   \allowdisplaybreaks
  $\begin{aligned}
  & P=\log _2 8+\log _3 27-\log _5 5^3 = \log _2  2^3+ \log _3 3^3 -\log _5 5^3 = 3+3-3=3.\\
  &Q=\ln (2\mathrm{e})-\log 100 = \ln 2 +  \ln \mathrm{e} -\log 10^2 = \ln 2 +1-2=\ln 2 -1.
 \end{aligned}$
 Suy  ra
  \allowdisplaybreaks
 $\begin{aligned}
 &  P+Q=2+\ln 2.\\
 & Q-P=\ln  2 - 4. \\
 &  3Q+P = 3\ln 2.\\
 & 2Q+P=2\ln 2  +1.
 \end{aligned}$
} 
\end{ex}

% nguồn: Ok_dungsai_lop11_Lan1
\begin{ex}%[1D6H2-1]%[ma:44b31e72]
 Cho biểu thức $B=2\ln\sqrt{\mathrm{e}x}-\ln \dfrac{\mathrm{e}^2}{\sqrt{x}}+\ln3\cdot\log_3\left(\mathrm{e}x^2\right)$.  Vậy:
  \choiceTF
  {\True Cho $\ln x=2$ thì $B=7$} 
  {\True Cho  $\ln x=4$ thì $B=14$}
  {Cho  $\ln x=3$ thì $B=\dfrac{21}{2}$}
  {Cho $\ln  x=6$ thì $B=18$}
  \loigiai{
   Ta có:
   $$
  \begin{aligned}
   B &  =2\ln(\mathrm{e}x)^{\frac{1}{2}}-\left(\ln\mathrm{e}^2-\ln x^{\frac{1}{2}}\right)+\ln\left(\mathrm{e}x^2\right).\\
   &  =\ln(\mathrm{e}x)-\left(2-\dfrac{1}{2}\ln x\right)+\ln \left(\mathrm{e}x^2\right)=(\ln\mathrm{e}+\ln x)-2+\dfrac{1}{2}\ln x+\ln\mathrm{e}+\ln x^2.\\
   & =1+\ln x-2+\dfrac{1}{2}\ln x+1+2 \ln x=\dfrac{7}{2}\ln  x.
  \end{aligned}
   $$
  \begin{itemchoice}
   \itemch Đúng  vì với $\ln x=2$ thì $B=\dfrac{7}{2}\cdot2=7$.
   \itemch Đúng vì với $\ln  x=4$ thì $B=\dfrac{7}{2}\cdot4=14$.
   \itemch Đúng  vì với $\ln x=3$ thì $B=\dfrac{7}{2}\cdot3=\dfrac{21}{2}$.
   \itemch Sai vì với $\ln x=6$  thì $B=\dfrac{7}{2}\cdot6=21$.
  \end{itemchoice}
  }
\end{ex}

% nguồn: Ok_dungsai_lop11_Lan1
\begin{ex}%[1D6H2-1]%[ma:d095ad22]
  Cho các biểu thức sau: $A=\log_{2^{2030}}4-\dfrac{1}{1015}+\ln{\mathrm{e}^{2035}}$; $ B=\log_53\cdot\log_25-\dfrac{\ln9}{\ln4}$.
 \choiceTF 
  {\True $A$  chia hết cho $5$}
   {$A-B=2036$}
  {\True $A+2024B=2035$} 
  {\True  $A-2024B=2035$}
 \loigiai{Ta có 
  \begin{eqnarray*}
    A&=&\log_{2^{2030}}4-\dfrac{1}{1015}+\ln{\mathrm{e}^{2035}}=\log_{2^{2030}}{2^2}-\dfrac{1}{1015}+2035\\
   &=&\dfrac{2}{2030}-\dfrac{1}{1015}+2035=2035. 
  \end{eqnarray*}
  Mặt  khác $B=\log_53\cdot\log_25-\dfrac{\ln9}{\ln4}=\log_25\cdot\log_53-\log_49=\log_23-\log_23=0$.\\
 Vậy ta có $A$ chia hết cho $5$  và $A+2024B=2035$, $A-2024B=2035$.}
\end{ex}

% nguồn: Ok_dungsai_lop11_Lan1
\begin{ex}%[1D6H2-1]%[ma:2683ce46]
 Tính được giá trị của biểu thức sau (biết $a>0$, $a\ne  1$). Vậy
 \choiceTF 
  {$A=2^{\log_23}-\log_{\sqrt{3}}3$ có $A>2$} 
  {\True $B=\ln 2\cdot  \log_24\cdot\log_43\cdot\log_32-5^{\log_5\ln 2}$ có $B=0$}
   {$C=\log_a\sqrt{a\sqrt{a\sqrt{a}}}$ có $C>1$}
  {$D=\log_a\dfrac{\sqrt{a^3}}{a\sqrt[4]{a}}$  có $D>1$}
  \loigiai{
  \begin{itemize}
   \item Ta có  $A=2^{\log_23}-\log_{\sqrt{3}}3=3-\log_{3^{\tfrac{1}{2}}}3=3-2=1$.
   \item Ta có  $B=\ln 2\cdot \log_24\cdot\log_43\cdot\log_32-5^{\log_5\ln 2}=\ln 2\cdot\log_23\cdot\log_32-\ln 2=\ln 2-\ln 2=0$.
    \item Ta có
   \allowdisplaybreaks 
   \begin{eqnarray*}
    \sqrt{a\sqrt{a\sqrt{a}}}&=&\sqrt{a\sqrt{a\cdot  a^{\tfrac{1}{2}}}}\\
    &=&\sqrt{a\sqrt{a^{\tfrac{3}{2}}}}\\ 
    &=&\sqrt{a\cdot a^{\tfrac{3}{4}}}\\ 
     &=&\sqrt{a^{\tfrac{7}{4}}}\\
     &=&a^{\tfrac{7}{8}}.
   \end{eqnarray*}
   Do đó,  $C=\log_a\sqrt{a\sqrt{a\sqrt{a}}}=\log_aa^{\tfrac{7}{8}}=\dfrac{7}{8}\log_aa=\dfrac{7}{8}$.
   \item Ta có $D=\log_a\dfrac{\sqrt{a^3}}{a\sqrt[4]{a}}=\log_a\dfrac{a^{\tfrac{3}{2}}}{a\cdot  a^{\tfrac{1}{4}}}=\log_aa^{\tfrac{3}{2}-1-\tfrac{1}{4}}=\log_a a^{\tfrac{1}{4}}=\dfrac{1}{4}$.
  \end{itemize}
 } 
\end{ex}

% nguồn: Ok_dungsai_lop11_Lan1
\begin{ex}%[1D6H2-1]%[ma:8dc1231b]
 Cho các  biểu thức $P=\log_2 8 +\log_3 27 -\log_5 {5^3}$ và $Q=\ln(2\mathrm{e})-\log 100$.
 \choiceTF 
  {$P+Q=2\ln 2$} 
   {\True $Q-P=\ln 2-4$}
  {\True  $3Q+P=3\ln 2$}
  {\True $2Q+P=2\ln 2+1$} 
 \loigiai{Ta có $P=\log_2  8 +\log_3 27 -\log_5 {5^3}=\log_2 2^3 +\log_3 3^3 -\log_5 5^3 =3+3-3=3$.\\
 $Q= \ln(2\mathrm{e})-\log 100=\ln 2+\ln \mathrm{e} -\log  10^2=\ln 2 +1-2=\ln 2-1$.
 \begin{itemchoice}
  \itemch Sai. $P+Q=3+\ln 2 -1=2+\ln 2$ 
  \itemch Đúng.  $Q-P=\ln 2-1 -3=\ln 2-4$.
  \itemch Đúng. $3Q+P=3\left(\ln 2 -1 \right)+3=3\ln 2$. 
  \itemch Đúng.  $2Q+P=2\left(\ln 2-1\right)+3=2\ln 2+1$.
 \end{itemchoice}
 } 
\end{ex}

% nguồn: Ok_dungsai_lop11_Lan1
\begin{ex}%[1D6H2-1]%[ma:2df9d14a]
 Cho $\log_ab=4$, với $a$ và $b$ là hai số dương  khác $1$.
 \choiceTF 
  {Giá trị của $\log_{a^2}b$ bằng $8$} 
  {\True Giá trị của $\log_a(a^2b)$ bằng  $6$}
  {\True Giá trị của $\log_b\sqrt[3]{a}$  bằng $\dfrac{1}{12}$}
  {Giá trị của $\log_{a^3b^2}(a^2b^3)$ bằng $4$} 
 \loigiai{ 
  \begin{itemchoice}
   \itemch Ta có $\log_{a^2}b=\dfrac{1}{2}\log_ab=\dfrac{1}{2}\cdot  4=2$. \\
   Vậy mệnh đề  đã cho sai.
    \itemch Ta có $\log_a(a^2 b)=\log_a a^2+\log_a b=2+\log_ab=2+4=6$. \\
   Vậy mệnh đề đã cho  đúng.
   \itemch Ta có $\log_b  \sqrt[3]{a}=\log_b a^{\tfrac{1}{3}}=\dfrac{1}{3}\log_b a=\dfrac{1}{3}\cdot \dfrac{1}{\log_a b}=\dfrac{1}{12}$. \\
    Vậy mệnh đề đã cho đúng.
   \itemch Ta  có
   $\log_{a^3b^2}(a^2b^3)=\dfrac{\log_a(a^2b^3)}{\log_a  (a^3b^2)}
   =\dfrac{\log_a  a^2+\log_a b^3}{\log_a a^3+\log_a b^2}
   =\dfrac{2\log_a  a+3\log_ab}{3\log_a a+2\log_ab}
   =\dfrac{2+3\cdot  4}{3+2\cdot 4}=\dfrac{14}{11}$.\\
    Vậy mệnh đề đã cho sai.
  \end{itemchoice}
 } 
\end{ex}

% nguồn: Ok_dungsai_lop11_Lan2
\begin{ex}%[1D6H2-1]%[ma:6e4e6fa5]
	Cho các  biểu thức sau: $P=\log_2 8+\log_3 27-\log_5 5^3$; $Q=\ln (2e)-\log 100$.
	\choiceTF 
		{$P+Q=2\ln  2$}
		{\True $Q-P=\ln  2-4$}
		{\True $3Q+P=3\ln  2$}
		 {\True $2Q+P=2\ln 2+1$}
	\loigiai{ 
		Ta có $P=\log_2 8+\log_3 27-\log_5  5^3=\log_2 2^3+\log_3 3^3-\log_5 5^3=3+3-3=3$.\\
		Ta có $Q=\ln (2e)-\log 100=\ln 2+\ln e-\log 10^2=\ln  2+1-2=\ln 2-1$.
		\begin{itemchoice}
			\itemch  $P+Q=2+\ln 2$.
			\itemch $Q-P=\ln 2 - 4$. 
			\itemch  $3Q+P=3\ln 2$.
			\itemch $2Q+P=2\ln 2 +  1$.
		\end{itemchoice}
	 }
\end{ex}

% nguồn: Ok_dungsai_lop11_Lan2
\begin{ex}%[1D6H2-1]%[ma:c5835ba6]
	Cho biểu thức $A=\log_2 x^2+\log _{\tfrac{1}{2}} x^3+\log_4 x$. 
	\choiceTF 
		{\True Khi  $\log_2 x=1$ thì $A=-\dfrac{1}{2}$}
		{Khi $\log_2 x=2$ thì $A=1$} 
		 {\True Khi $\log_2 x=\sqrt{3}$ thì $A=-\dfrac{\sqrt{3}}{2}$}
		{Khi $\log_2 x=\sqrt{2}$  thì $A=\dfrac{\sqrt{2}}{2}$}
	\loigiai{ 
		Ta có 
		\[\begin{aligned}
		A&=\log_2 x^2+\log _{\tfrac{1}{2}} x^3+\log_4 x  \\
		&=2\log_2 x+3\log _{2^{-1}} x+\log  _{2^2} x\\
		&=2\log_2  x-3\log_2 x+\dfrac{1}{2} \log_2 x\\
		&=-\dfrac{1}{2} \log_2 x. 
	\end{aligned}\]
	\begin{itemchoice}
		\itemch  Khi $\log_2 x=1$ thì $A=-\dfrac{1}{2}$.
		\itemch  Khi $\log_2 x=2$ thì $A=-1$.
		\itemch Khi  $\log_2 x=\sqrt{3}$ thì $A=-\dfrac{\sqrt{3}}{2}$.
		\itemch Khi $\log_2 x=\sqrt{2}$ thì $A=-\dfrac{\sqrt{2}}{2}$. 
	\end{itemchoice}
} 
\end{ex}

% nguồn: Ok_dungsai_lop11_Lan2
\begin{ex}%[1D6H2-1]%[ma:215221d8]
	Cho các biểu thức sau: $A=\log _{2^{2030}} 4-\dfrac{1}{1015}+\ln e^{2035}$; $B=\log_5 3\cdot \log_2 5-\dfrac{\ln 9}{\ln 4}$. 
	\choiceTF 
		{\True $A$ chia hết cho  $5$}
		{$A-B=2036$} 
		 {\True $A+2024B=2035$}
		 {\True $A-2024B=2035$}
	 \loigiai{
		\begin{itemchoice}
			 \itemch Ta có
			\begin{eqnarray*}
				A&=&\log _{2^{2030}}  4-\dfrac{1}{1015}+\ln e^{2035}\\
				&=&  \log _{2^{2030}} 2^2-\dfrac{1}{1015}+2035\\
				 &=& \dfrac{2}{2030}-\dfrac{1}{1015}+2035\\
				&=& 2035. 
			\end{eqnarray*}
			 Vậy $A$ chia hết cho $5$.
			\itemch Ta  có
			\begin{eqnarray*}
				B&=&\log_5  3\cdot\log_2 5-\dfrac{\ln 9}{\ln 4}\\
				&=& \log_2  5\cdot\log_5 3-\log_4 9\\
				&=& \log_2 3-\log _{2^2}  3^2=\log_2 3-\log_2 3\\
				&=&  0.
			\end{eqnarray*}
			 Vậy $A-B=2035$.
			\itemch  $A+2024B=2035$.
			\itemch  $A-2024B=2035$.
		\end{itemchoice}
	} 
\end{ex}

% nguồn: Ok_dungsai_lop11_Lan2
\begin{ex}%[1D6H2-1]%[ma:21c4ab8d]
	Cho biểu thức $B=2\ln \sqrt{ex}-\ln \dfrac{e^2}{\sqrt{x}}+\ln 3\cdot  \log_3 (ex^2)$.
	\choiceTF 
		{\True Cho $\ln x=2$  thì $B=7$}
		{\True Cho  $\ln x=4$ thì $B=14$}
		{Cho $\ln  x=3$ thì $B=\dfrac{15}{2}$}
		{Cho $\ln  x=6$ thì $B=18$}
	\loigiai{ 
		Ta có 
		\[\begin{aligned}
		B&=2\ln (ex)^{\tfrac{1}{2}}-\left(\ln  e^2-\ln x^{\tfrac{1}{2}} \right)+\ln (ex^2) \\
		 & =\ln (ex)-\left(2-\dfrac{1}{2} \ln x\right)+\ln (ex^2)\\
		&=(\ln e+\ln x)-2+\dfrac{1}{2} \ln  x+\ln e+\ln x^2 \\
		&=1+\ln x-2+\dfrac{1}{2} \ln x+1+2\ln  x\\
		 &=\dfrac{7}{2} \ln x.
	\end{aligned}\]
	\begin{itemchoice}
		\itemch Cho $\ln x=2$  thì $B=7$.
		\itemch Cho $\ln x=4$ thì $B=14$. 
		\itemch  Cho $\ln x=3$ thì $B=\dfrac{21}{2}$.
		\itemch Cho  $\ln x=6$ thì $B=21$.
	\end{itemchoice}
 }
\end{ex}

% nguồn: Ok_dungsai_lop11_Lan2
\begin{ex}%[1D6V2-1]%[ma:ab327072]
	Cho $a=\log _{12}  18$, $b=\log _{24} 54$.
	\choiceTF 
		{\True  $\log_2 3=\dfrac{2a-1}{2-a}$}
		{\True  $\log_2 3=\dfrac{3b-1}{3-b}$}
		{\True $ab+5\left(a-b\right)=1$} 
		{\True  $3\cdot12^a=24^b$}
	\loigiai{ 
		\begin{itemchoice}
			\itemch Ta có $a=\log _{12} 18=\dfrac{\log_2 18}{\log_2 12}=\dfrac{1+2\log_2  3}{2+\log_2 3} \Leftrightarrow \log_2 3=\dfrac{2a-1}{2-a}$.
			\itemch $b=\log _{24} 54=\dfrac{\log_2 54}{\log_2 24}=\dfrac{1+3\log_2 3}{3+\log_2 3} \Leftrightarrow  \log_2 3=\dfrac{3b-1}{3-b}.$
			\itemch  Ta có $\dfrac{2a-1}{2-a}=\dfrac{3b-1}{3-b} \Leftrightarrow 5\left(a-b\right)+ab=1$.
			\itemch  $a=\log _{12} 18,b=\log _{24} 54\Rightarrow 12^a=18; 24^b=54\Rightarrow 3\cdot 12^a=24^b$.
		\end{itemchoice}
	} 
\end{ex}

% nguồn: Ok_dungsai_lop11_Lan2
\begin{ex}%[1D6V2-1]%[ma:6bb197fa]
	Gọi $x$, $y$ là các số thực dương thỏa mãn điều kiện $\log_9 x=\log _{12} y=\log _{16} \left(x+y\right)=t$ và $\dfrac{x}{y}=\dfrac{-a+\sqrt{b}}{2}$,  với $a, \; b$ là hai số nguyên dương.
	 \choiceTF
		 {$9^t+12^t+16^t=0$}
		{$\dfrac{x}{y}=\left(\dfrac{4}{3} \right)^t$} 
		{\True  $\left(\dfrac{3}{4} \right)^t=\dfrac{-1+\sqrt{5}}{2}$}
		{\True $\dfrac{y}{x}=\dfrac{2}{-1+\sqrt{5}}$} 
	\loigiai{ 
		Điều kiện $\heva{&x >  0 \\&y > 0.}$
		\begin{itemchoice}
			\itemch Ta có $t=\log_9 x=\log _{12} y=\log _{16}  \left(x+y\right)\Rightarrow \heva{&x=9^t \\&y=12^t \\&x+y=16^t.}$ \\
			Suy ra  $9^t+12^t=16^t \Leftrightarrow 9^t+12^t-16^t=0$.
			 \itemch Ta có $\dfrac{x}{y}=\dfrac{9^t}{12^t}=\left(\dfrac{3}{4} \right)^t$.
			\itemch  Ta có $9^t+12^t-16^t=0 \Leftrightarrow \dfrac{9^t}{16^t}+\dfrac{12^t}{16^t}-1=0\Leftrightarrow \left(\left(\dfrac{3}{4} \right)^t \right)^2+\left(\dfrac{3}{4} \right)^t-1=0$
			 \[\Rightarrow \hoac{&\left(\dfrac{3}{4} \right)^t=\dfrac{-1+\sqrt{5}}{2} \\ &\left(\dfrac{3}{4} \right)^t=\dfrac{-1-\sqrt{5}}{2}.}\]
			\itemch Mà $\dfrac{x}{y}=\dfrac{-a+\sqrt{b}}{2}$. Suy ra $\dfrac{x}{y}=\dfrac{-1+\sqrt{5}}{2}\Rightarrow \dfrac{y}{x}=\dfrac{2}{-1+\sqrt{5}}$. 
		\end{itemchoice}
	} 
\end{ex}


%%%%%%%%%%%%%%%% TRẢ LỜI NGẮN %%%%%%%%%%%%%%%%

% nguồn: Ok_TraLoiNgan_Lop11_Lan1
\begin{ex}%[1D6H2-1]%[ma:bf1c2b55]
	Cho các  số thực dương $x, y$ thoả mãn $x^2+y^2=14 x y$. Xác định giá trị $a$ thỏa mãn $\log _2(x+y)=a+\dfrac{\log _2 x y}{a}$.
	\shortans{2} 
	 \loigiai{
	Ta  có
	\[\begin{aligned}
	 x^2+y^2=14 x y &\Leftrightarrow(x+y)^2=16 x y \\
	&\Leftrightarrow \log _2(x+y)^2=\log  _2(16 x y)\\
	&\Leftrightarrow 2  \log _2(x+y)=4+\log _2(x y) \\
	&\Leftrightarrow \log  _2(x+y)=2+\frac{\log _2(x y)}{2}.
\end{aligned}\]
Vậy $a=2$. 
} 
\end{ex}

% nguồn: Ok_TraLoiNgan_Lop11_Lan1
\begin{ex}%[1D6H2-1]%[ma:386d9070]
	Tính giá trị  của biểu thức
	$B=\log\dfrac{1}{1000}+3\cdot \log_{\frac{1}{10}}100-10^{1+\log  2}$.
	 \shortans{-29}
	\loigiai{ 
		Ta có \begin{itemize}
		\item $\log\dfrac{1}{1000}=\log 10^{-3}=-3$. 
		\item  $\log_{\frac{1}{10}}100=\log_{10^{-1}}10^2=-2$.
		\item $10^{1+\log 2}=10^{\log10+\log 2}=10^{\log 20}=20$. 
	\end{itemize}
	Vậy $B=-3+3(-2)-20=-29$. 
 }
\end{ex}

% nguồn: Ok_TraLoiNgan_Lop11_Lan1
\begin{ex}%[1D6H2-1]%[ma:91a08836]
	Tính giá  trị biểu thức $A=\log_2\left(2\sin\dfrac{\pi}{12}\right)+\log_2\left(\cos\dfrac{\pi}{12}\right)$.
	 \shortans{-1}
	\loigiai{Ta có \begin{align*}
		A&\ =\log_2\left(2  \sin\dfrac{\pi}{12}\right)+\log_2\left(\cos\dfrac{\pi}{12}\right)\\
		&\ =\log_2\left(2  \sin\dfrac{\pi}{12} \cdot\cos\dfrac{\pi}{12}\right)\\
		&\ =\log_2\left(\sin\dfrac{\pi}{6}\right)\\ 
		&\ =\log_2\dfrac{1}{2}=-1. 
\end{align*}}
\end{ex}

% nguồn: Ok_TraLoiNgan_Lop11_Lan1
\begin{ex}%[1D6H2-1]%[ma:cf42dc65]
	Tính giá trị của biểu  thức $A=9^{\log_3 5 +\log_3 2}$.
	 \shortans{100}
	\loigiai{Ta có \begin{align*}
	A&\ =9^{\log_3 5 +\log_3 2}=3^{2\log_3 5}\cdot  3^{2\log_3 2}\\
	&\  =3^{\log_3 5^2}\cdot 3^{\log_3 2^2}\\
	 &\ =5^2\cdot 2^2 =100.
\end{align*}}
\end{ex}

% nguồn: Ok_TraLoiNgan_Lop11_Lan1
\begin{ex}%[1D6H2-1]%[ma:62218973]
	Cho $a=\log _2 3$, $b=\log _5 2$, $c=\log _2 7$. Biểu diễn $\log _{42} 15$ theo  $a$, $b$, $c$ ta thu được kết quả dạng $\dfrac{a b+m}{h \cdot b(a+c+n)}$ với $h$, $m$, $n$ là các số tự nhiên. Tính giá trị $P=m \cdot n+h$.
	 \shortans{2}
	 \loigiai{
	 Ta có $\log _{42} 15=\dfrac{\log _2 15}{\log _2 42}=\dfrac{\log _2 3+\log _2 5}{\log _2 2+\log _2 3+\log _2 7}=\dfrac{a+\frac{1}{b}}{1+a+c}=\dfrac{a b+1}{b(a+c+1)}$.\\
	Khi đó $\heva{ & m=1 \\ & n=1 \\ & h=1} \Rightarrow  P=2$.
	} 
\end{ex}

% nguồn: Ok_TraLoiNgan_Lop11_Lan2
\begin{ex}%[1D6H2-1]%[ma:d6fc7cd0]
	 Nếu $x=\log_47+\log_{16}7$ thì $4^{2x}$ có giá trị bằng
	 \shortans{343}
	\loigiai{ 
		Ta có 
		\begin{align*}
			x=\log_47+\log_{16}7 \Leftrightarrow x=\dfrac{3}{2}\log_47  \Rightarrow 4^{2x}=4^{2\left(\tfrac{3}{2}\log_47\right)}=7^3=343.
		\end{align*}
	 }
\end{ex}

% nguồn: Ok_TraLoiNgan_Lop11_Lan2
\begin{ex}%[1D6H2-1]%[1D6H2-2]%[1D6H2-3]%[ma:a811fe89]
	Cho các  số thực dương $a$, $b$, $c$ với $c\ne 1$ thoả mãn $\log_ab=3$, $\log_ac=-2$. Khi đó $\log_a\left(a^3b^2\sqrt{c}\right)$ bằng
	 \shortans{8}
	 \loigiai{
		Ta  có
		\begin{align*}
			&\log_a\left(a^3b^2\sqrt{c}\right) =\log_aa^3+\log_ab^2+\log_a\sqrt{c}.\\ 
			&\Rightarrow  \log_a\left(a^3b^2\sqrt{c}\right) =3+2\log_ab+\dfrac{1}{2}\log_ac.\\
			 &\Rightarrow \log_a\left(a^3b^2\sqrt{c}\right)=8.
		\end{align*}
	} 
\end{ex}

% nguồn: Ok_TraLoiNgan_Lop11_Lan2
\begin{ex}%[1D6H2-1]%[ma:c8c3bb90]
	Nếu  $x=\log_34+\log_94$ thì $3^x$ có giá trị bằng
	\shortans{8} 
	 \loigiai{
		 Ta có
		\begin{align*}
			x=\log_34+\log_94 \Leftrightarrow x=\log_34+\dfrac{1}{2}\log_34  \Leftrightarrow x=\dfrac{3}{2}\log_34 \Leftrightarrow 3^x=(4)^{\tfrac{3}{2}}=8.
		\end{align*}
	 }
\end{ex}

% nguồn: Ok_TraLoiNgan_Lop11_Lan2
\begin{ex}%[1D6H2-1]%[ma:f0835e58]
	Cho $\log a=10$; $\log b=100$. Khi đó $\log (a\cdot b^3)$ bằng 
	 \shortans{310}
	\loigiai{ 
		Ta có $$\log (a\cdot b^3)=\log a+\log b^3=\log a  +3\log b =10+3\cdot100=310.$$
	} 
\end{ex}

% nguồn: Ok_TraLoiNgan_Lop11_Lan2
\begin{ex}%[1D6H2-1]%[1D6H2-2]%[ma:eda285e9]
	Cho $\log_{ab}b=3$, $a>0$, $b>0$, $ab\ne 1$. Tính $\log_{\sqrt{ab}}\left(\dfrac{a}{b^2}\right)$. 
	 \shortans{-16}
	\loigiai{ 
		Ta có 
		\begin{align*}
			\bullet \log_{ab}b=3 &\Leftrightarrow \dfrac{1}{\log_b(ab)}=3 \\ 
			 &\Leftrightarrow 3(\log_ba+\log_bb)=1\\
			&\Leftrightarrow  3\log_ba+3=1\\
			&\Leftrightarrow \log_ba=-\dfrac{2}{3}\\ 
			 &\Rightarrow \log_ab=-\dfrac{3}{2}.
		\end{align*}
		\begin{align*}
			 \bullet \log_{\sqrt{ab}}\left(\dfrac{a}{b^2}\right) &=\log_{\sqrt{ab}}a-\log_{\sqrt{ab}}b^2 =2\log_{ab}a-4\log_{ab}b\\
			 &=\dfrac{2}{\log_aab}-4\cdot3 =\dfrac{2}{1+\log_ab}-12\\
			&=\dfrac{2}{1-\tfrac{3}{2}}-12=-16. 
		\end{align*}
	 }
\end{ex}

% nguồn: Ok_TraLoiNgan_Lop11_Lan2
\begin{ex}%[1D6H2-1]%[ma:7c2a1dca]
	Cho $a$, $b$ là các số dương, $a\ne 1$ sao cho $\log_ab=2$, giá  trị của $\log_a(a^3b)$ bằng
	\shortans{5} 
	\loigiai{ 
		Ta có $\log_a(a^3b)=3+\log_ab=3+2=5$. 
	} 
\end{ex}

% nguồn: Ok_TraLoiNgan_Lop11_Lan1
\begin{ex}%[1D6V2-1]%[ma:9bfdd1fb]
	Với $a$, $b$ là các số thực dương và thoả mãn $a b \neq 1$. Ta rút gọn biểu thức \[P=\left(\log _a b+\log _b a+2\right)\left(\log _a b-\log _{a  b} b\right) \log _b a-1\] thu được kết quả $(2 m-1) \log _a b^n$ với $m$, $n$ là các số tự nhiên. Tính giá trị $S=m-n$.
	\shortans{0} 
	 \loigiai{
	$$ 
	\begin{aligned}
		P&=\left(\log _a b+\log _b  a+2\right)\left(\log _a b-\log _{a b} b\right) \log _b a-1 \\
		& =\left(\log _a b+\frac{1}{\log _a b}+2\right)\left(\log _a b-\frac{\log _a b}{\log _a(a b)}\right) \log _b a-1  \\
		& =\frac{\log _a^2  b+2 \log _a b+1}{\log _a b}\left(\log _a b-\frac{\log _a b}{1+\log _a b}\right) \log _b a-1 \\
		& =\frac{\left(\log _a b+1\right)^2}{\log _a b} \cdot\frac{\log _a^2 b}{1+\log _a b} \cdot\log _b a-1\\ 
		&=\left(\log _a b+1\right) \log _a b\cdot \log _b a-1\\ 
		&=\log _a b. 
	\end{aligned}
	$$ 
	Khi đó $\heva{&2m - 1 =1\\&n=1}\Leftrightarrow \heva{&m=1\\&n=1} \Rightarrow S  = m-n = 0$.
	} 
\end{ex}

% nguồn: Ok_TraLoiNgan_Lop11_Lan1
\begin{ex}%[1D6V2-1]%[ma:91b4ecf9]
	Cho $a$, $b>0$ và đều khác $1$ thoả mãn $\ln a+\ln (8 b)=2 \ln (a+2 b)$. Rút gọn biểu thức $P=\log _b(2 a)+\log _{\frac{a}{2}}(2 b)-\dfrac{1}{\log _8 b}$ được kết quả bằng  bao nhiêu?
	\shortans{2} 
	 \loigiai{
	Với  $a$, $b$ là các số thực dương khác $1$, ta có
	 $$
	\begin{aligned}
		 &\ln a+\ln (8 b)=2 \ln (a+2 b) \\
		\Leftrightarrow\;& \ln (8 a  b)=\ln (a+2 b)^2 \\
		\Leftrightarrow\;& 8 a b=(a+2 b)^2 \\ 
		\Leftrightarrow\;&\left(a+2 b\right)^2=0 \\ 
		 \Leftrightarrow\;& a=2 b.
	\end{aligned}
	$$ 
	Khi đó $P=\log _b(2 a)+\log _{\frac{a}{2}}(2 b)-\dfrac{1}{\log _8 b}=\log _b(4 b)+\log _b(2 b)-\log  _b 8
	=\log  _b \dfrac{8 b^2}{8}=\log _b b^2=2$.
	} 
\end{ex}

% nguồn: Ok_TraLoiNgan_Lop11_Lan1
\begin{ex}%[1D6V2-1]%[ma:af6e7733]
	Cho $a$, $b$ là các số thực dương khác $1$ và thoả mãn $a b \neq 1$. Biết rằng \[\log _a \frac{3}{b}=\left(\log _a b+\log  _b a+2\right)\left(\log _a b-\log _{a b} b\right) \log _b a-1.\]
	Tìm $b$ (làm tròn kết quả đến hàng phần chục). 
	 \shortans{1,7}
	 \loigiai{
	Ta có 
	 $$
	\begin{aligned}
		P 
		&=\left(\log _a b+\log _b a+2\right)\left(\log _a b-\log _{a b} b\right)  \log _b a-1 \\
		&=\left(\log _a b+\frac{1}{\log _a  b}+2\right)\left(\log _a b-\frac{\log _a b}{\log _a(a b)}\right) \log _b a-1 \\
		&=\frac{\log _a^2 b+2 \log _a b+1}{\log _a b}\left(\log _a b-\frac{\log _a b}{1+\log _a b}\right) \log _b a-1  \\
		&=\frac{\left(\log _a b+1\right)^2}{\log _a b}\cdot\frac{\log _a^2  b}{1+\log _a b}\cdot \log _b a-1=\left(\log _a b+1\right) \log _a b\cdot \log _b a-1=\log _a b.
	\end{aligned}
	 $$
	Khi đó $\log _a \dfrac{3}{b}=P \Leftrightarrow \log  _a \dfrac{3}{b}=\log _a b \Leftrightarrow \dfrac{3}{b}=b \Leftrightarrow b^2=3 \Leftrightarrow b=\sqrt{3} \approx 1{,}7$.
	} 
\end{ex}

% nguồn: Ok_TraLoiNgan_Lop11_Lan1
\begin{ex}%[1D6V2-1]%[ma:41101898]
	Cho $x \in \left(0;\dfrac{\pi}{2}\right)$ biết $\log \sin x + \log  \cos x=-1$ và $\log \left(\sin x+\cos x\right)=\dfrac{1}{2} \left(\log n-1\right)$. Giá trị của $n$ là
	\shortans{12} 
	\loigiai{Ta có $$\log \sin x + \log  \cos x=-1 \Leftrightarrow \log \left(\sin x \cos x\right)=-1 \Leftrightarrow \sin x \cos x=10^{-1} \Leftrightarrow \sin x \cos x = \dfrac{1}{10}.$$
	Do đó 
	\allowdisplaybreaks 
	\begin{align*}
		&\log  \left(\sin x+\cos x\right)=\dfrac{1}{2} \left(\log n-1\right)\\
		& \Leftrightarrow 2\log \left(\sin x+\cos x\right)=  \log n-1 \\
		&\Leftrightarrow \log \left(\sin x+\cos x\right)^2= \log n  -\log 10\\
		&\Leftrightarrow \log \left(\sin^2 x+\cos^2 x+ 2 \sin x  \cos x \right) = \log \dfrac{n}{10}\\
		&\Leftrightarrow \log \left(1+2 \cdot  \dfrac{1}{10}\right)=\log \dfrac{n}{10} \\
		&\Leftrightarrow  \dfrac{6}{5}=\dfrac{n}{10}\\
		&\Leftrightarrow  n=12.
	\end{align*}
	 Vậy $n=12$.
	} 
\end{ex}

% nguồn: Ok_TraLoiNgan_Lop11_Lan1
\begin{ex}%[1D6V2-1]%[ma:2e4de241]
	Cho $x=2024!$. Tính giá trị của  biểu thức $A=\dfrac{1}{\log_2x}+\dfrac{1}{\log_3x}+\cdots +\dfrac{1}{\log_{2024}x}$.
	\shortans{1} 
	\loigiai{ 
		Theo bài $x=2024!\Rightarrow x>0,x\neq 1$. 
		\allowdisplaybreaks	\begin{eqnarray*}
		A&=&\dfrac{1}{\log_2x}+\dfrac{1}{\log_3x}+\cdots +\dfrac{1}{\log_{2024}x}=\log_x2+\log_x3+\cdots +\log_x2024\\ 
		&=&\log_x\left(1\cdot  2\cdot 3\cdots 2024\right)=\log_x2024!.
	\end{eqnarray*}
 Với $x=2024!\Rightarrow A=\log_{2024!}2024!=1$.}
\end{ex}

% nguồn: Ok_TraLoiNgan_Lop11_Lan2
\begin{ex}%[1D6V2-1]%[ma:af0331c5]
	Cho số thực $a$, $b$ thỏa mãn $a>b>1$ và $\dfrac{1}{\log_b a}+\dfrac{1}{\log_b b}=\sqrt{2024}$. Giá trị biểu thức $P=\dfrac{1}{\log_{ab} b}-\dfrac{1}{\log_{ab} a}$  bằng (kết quả làm tròn đến hàng phần mười)
	 \shortans{44,9}
	\loigiai{ 
		Ta có 
		$$ 
		\dfrac{1}{\log_a b} +  \dfrac{1}{\log_b a} = \sqrt{2024}
		 \Leftrightarrow
		\log_a b + \log_b a = \sqrt{2024} \quad  (1)
		 $$
		 và
		 $$
		P =  \dfrac{1}{\log_{ab} b} - \dfrac{1}{\log_{ab} a}
		= \log_b(ab) - \log_a(ab) 
		= (\log_b a+1) -(\log_a b+1) 
		 = \log_b a-\log_a b \quad (2)
		 $$
		 Từ (1), suy ra
		 $$
		\log_a^2 b + \log_b^2 a + 2\log_b b.\log_b a =  2024
		\Leftrightarrow 
		 \log_a^2 b + \log_b^2 a = 2022.
		 $$
		Từ (2),  suy ra
		 $$
		P^2 = \log_b^2 b + \log_b^2 a  - 2\log_b b.\log_b a
		=  2022-2
		 =2020.
		$$ 
		Do $a>b>1$ nên $\log_b b<1$ và $\log_b  a>1$ nên $P>0$. \\
		Vậy,  $P = \sqrt{2020}$.
	} 
\end{ex}

% nguồn: Ok_TraLoiNgan_Lop11_Lan2
\begin{ex}%[1D6V2-1]%[ma:84b8e4b9]
	Cho các số thực dương  $x$, $y$ với $x\ne3$ và $xy^2=3$. Giá trị của $\log_{y^2}\dfrac{x^3}{27}$ bằng
	 \shortans{-3}
	\loigiai{ 
		Ta có $xy^2=3 \Leftrightarrow  y^2=\dfrac{3}{x}$ $\Rightarrow \log_{y^2}\dfrac{x^3}{27} =\log_{\tfrac{3}{x}}\left(\dfrac{x}{3}\right)^3=-3$.
	 }
\end{ex}

% nguồn: Ok_TraLoiNgan_Lop11_Lan2
\begin{ex}%[1D6V2-1]%[ma:38edb8d2]
	Cho $2$ số thực dương $a$, $b$ thỏa mãn $\sqrt{a}\ne b$, $a\ne 1$, $\log_ab=3$. Tính $T={\log_{\tfrac{\sqrt{a}}{b}}}\sqrt[3]{ba}$. (kết quả  làm tròn đến hàng phần chục)
	\shortans{-0,5} 
	 \loigiai{
		Ta có $\log_ab=3\Rightarrow \log_ba=\dfrac{1}{3}$.\\ 
		Khi đó 
		\begin{align*}
			T&=\log_{\tfrac{\sqrt{a}}{b}}\sqrt[3]{ba} =\log_{\tfrac{\sqrt{a}}{b}}\sqrt[3]{b} +\log_{\tfrac{\sqrt{a}}{b}}\sqrt[3]{a}\\ 
			&=\dfrac{1}{\log_{\sqrt[3]{b}}\dfrac{\sqrt{a}}{b}} +\dfrac{1}{\log_{\sqrt[3]{a}}\dfrac{\sqrt{a}}{b}}  =\dfrac{1}{\log_{\sqrt[3]{b}}\sqrt{a}-\log_{\sqrt[3]{b}}b} +\dfrac{1}{\log_{\sqrt[3]{a}}\sqrt{a}-\log_{\sqrt[3]{a}}b}\\
			&=\dfrac{1}{\tfrac{3}{2}\log_ba-3} +\dfrac{1}{\tfrac{3}{2}-3\log_ab}  =\dfrac{1}{\tfrac{3}{2}\cdot\tfrac{1}{3}-3}+\dfrac{1}{\tfrac{3}{2}-3\cdot3} =-\dfrac{8}{15}.
		\end{align*}
	 }
\end{ex}

% nguồn: Ok_TraLoiNgan_Lop11_Lan2
\begin{ex}%[1D6V2-1]%[ma:2fd89578]
	Cho $x, y$ là  các số thực lớn hơn $1$ thoả mãn $x^2+9y^2=6xy$. Tính $M=\dfrac{1+\log_{12} x+\log_{12} y}{2\log_{12}(x+3y)}$.
	\shortans{1} 
	 \loigiai{
		Ta có 
		\begin{align*}
			x^2 + 9y^2 =  6xy
			&\Leftrightarrow x^2 +(3y)^2  + 2x.3y = 6xy + 6xy \\
			&\Leftrightarrow(x+3y)^2  = 12xy\\
			 &\Leftrightarrow \log_{12} (x+3y)^2 = \log_{12}(12xy) \\
			&\Leftrightarrow 2\log_{12}(x+3y) = 1+\log_{12} x+\log_{12}  y \\
			&\Leftrightarrow \dfrac{1+\log_{12} x+\log_{12} y}{2\log_{12}(x+3y)}  = 1.
		\end{align*}
	 Vậy, $M = 1$.}
\end{ex}

% nguồn: Ok_TraLoiNgan_Lop11_Lan2
\begin{ex}%[1D6V2-1]%[ma:165a632f]
	Giá trị của biểu thức $N=\dfrac{1}{\log_{2^{2023}} 2023!}+\dfrac{1}{\log_{3^{2023}} 2023!}+\dfrac{1}{\log_{4^{2023}} 2023!}+\dots+\dfrac{1}{\log_{2023^{2023}} 2023!}$ là 
	 \shortans{2023}
	\loigiai{ 
		 Ta có
		\begin{align*}
			N &= \dfrac{1}{\log_{2^{2023}} 2023!} +  \dfrac{1}{\log_{3^{2023}} 2023!} + \dfrac{1}{\log_{4^{2023}} 2023!} + \dots + \dfrac{1}{\log_{2023^{2023}} 2023!}\\
			&= \log_{2023!} 2^{2023}  + \log_{2023!} 3^{2023} + \log_{2023!} 4^{2023} + \dots + \log_{2023!} 2023^{2023}\\
			&= 2023(\log_{2023!} 2 + \log_{2023!} 3 + \log_{2023!} 4  + \dots + \log_{2023!} 2023)\\
			&=  2023\log_{2023!}(2.3.4.\dots.2023)\\
			&=  2023\log_{2023!}(2023!) \\
			&= 2023. 
		\end{align*}
	} 
\end{ex}

% nguồn: Ok_TraLoiNgan_Lop11_Lan2
\begin{ex}%[1D6V2-1]%[ma:9e0e9f5a]
	Cho $\log_ab=3,\log_ac=-2$. Khi đó $\log_a\left(a^3b^2\sqrt{c}\right)$  bằng bao nhiêu?
	\shortans{8} 
	 \loigiai{
		Ta có 
		\begin{align*}
			\log_a\left(a^3b^2\sqrt{c}\right) =\log_aa^3+\log_ab^2+\log_a\sqrt{c} =3+2\log_ab+\dfrac{1}{2}\log_ac =3+2\cdot 3+\dfrac{1}{2}\cdot  (-2)=8.
		\end{align*}
	} 
\end{ex}


%%%%%%%%%%%%%%%% TỰ LUẬN %%%%%%%%%%%%%%%%

% nguồn: Ok_TuanLuan_Toan11_Lan1
\begin{bt}%[1D6N2-1]%[ma:1f71db66]
	Cho số thực $a$ thõa mãn $0<a\neq 1$. Tính giá  trị của biểu thức $A=\log_812-\log_815+\log_820$.
	\loigiai{ 
		Ta có  $A=\log_812-\log_815+\log_820=\log_8\dfrac{12\cdot 20}{15}=\log_816=\dfrac{4}{3}$.
	} 
\end{bt}

% nguồn: Ok_TuanLuan_Toan11_Lan1
\begin{bt}%[1D6N2-1]%[ma:e54bf620]
	 Tính $M=\log_4{1\,250}$ theo $a$ biết $a=\log_25$.
	\loigiai{Ta  có $M=\log_4{1250}=\dfrac{1}{2}\log_2{\left(5^4\cdot 2\right)}=\dfrac{1}{2}\left(4\log_25+1\right)=\dfrac{1}{2}+2a$.}
\end{bt}

% nguồn: Ok_TuanLuan_Toan11_Lan1
\begin{bt}%[1D6N2-1]%[ma:b769ab6c]
	Cho $\log_32=a$ và $\log_35$.  Tính $\log_{10}{60}$ theo $a$ và $b$.
	\loigiai{$\log_{10}{60}=\dfrac{\log_3{60}}{\log_3{10}}=\dfrac{\log_3\left(4\cdot 3\cdot 5\right)}{\log_3\left(2\cdot  5\right)}=\dfrac{2\log_32+\log_35+1}{\log_32+\log_35}=\dfrac{2a+b+1}{a+b}$.}
\end{bt}

% nguồn: Ok_TuanLuan_Toan11_Lan2
\begin{ex}%[1D6N2-1]%[ma:69a543e8]
	Cho số thực $a$ thõa mãn $0<a\neq 1$. Tính giá trị  của biểu thức $A=\log_812-\log_815+\log_820$.
	\loigiai{ 
		 Ta có $A=\log_812-\log_815+\log_820=\log_8\dfrac{12\cdot 20}{15}=\log_816=\dfrac{4}{3}$.
	 }
\end{ex}

% nguồn: Ok_TuanLuan_Toan11_Lan2
\begin{ex}%[1D6N2-1]%[ma:c1df20d4]
	Tính $M=\log_4{1\,250}$ theo $a$ biết  $a=\log_25$.
	 \loigiai{Ta có $M=\log_4{1250}=\dfrac{1}{2}\log_2{\left(5^4\cdot 2\right)}=\dfrac{1}{2}\left(4\log_25+1\right)=\dfrac{1}{2}+2a$.}
\end{ex}

% nguồn: Ok_TuanLuan_Toan11_Lan2
\begin{ex}%[1D6N2-1]%[ma:3c89a728]
	Cho $\log_32=a$ và $\log_35$. Tính $\log_{10}{60}$ theo  $a$ và $b$.
	\loigiai{$\log_{10}{60}=\dfrac{\log_3{60}}{\log_3{10}}=\dfrac{\log_3\left(4\cdot 3\cdot  5\right)}{\log_3\left(2\cdot 5\right)}=\dfrac{2\log_32+\log_35+1}{\log_32+\log_35}=\dfrac{2a+b+1}{a+b}$.}
\end{ex}

% nguồn: Ok_TuanLuan_Toan11_Lan1
\begin{bt}%[1D6H2-1]%[ma:b1f0d3bf]
	Cho $\log_a x=3,\log_b x=4$ với $a > 1$, $b > 1$, $x >  1$. Giá trị của $P=\log _{ab} x$ bằng bao nhiêu?
	 \loigiai{
		Ta có $P=\log _{ab} x=\dfrac{1}{\log_x  (ab)}=\dfrac{1}{\log_x a+\log_x b}=\dfrac{1}{\dfrac{1}{3}+\dfrac{1}{4}}=\dfrac{12}{7}$.
	 }
\end{bt}

% nguồn: Ok_TuanLuan_Toan11_Lan1
\begin{bt}%[1D6H2-1]%[ma:8a9edcda]
	Cho $f(x)=\mathrm e ^{\sqrt{1+\frac{1}{x^2}+\frac{1}{(x+1)^2}}}$. Biết rằng $f(1) \cdot f(2) \cdot f(3) \cdots f(2025)=\mathrm e ^{\frac{m}{n}}$  với $m, n$ là các số tự nhiên và $\frac{m}{n}$ là phân số tối giản. Tính $m-n^2$.
	 \loigiai{ \\
		Đặt  $g(x)=\sqrt{1+\dfrac{1}{x^2}+\dfrac{1}{(1+x)^2}}$ \\
		Với ${x>0}$ ta  có\\
		\begin{eqnarray*}
			 g(x)&=&\sqrt{1+\dfrac{1}{x^2}+\frac{1}{(1+x)^2}}=\dfrac{\sqrt{x^2+(x+1)^2+x^2\cdot (x+1)^2}}{x(x+1)}=\dfrac{\sqrt{\left(x^2+x+1\right)^2}}{x \left(x+1 \right)} \\
			&=&\dfrac{{{x}^{2}}+x+1}{x(x+1)}=1+\dfrac{1}{x(x+1)}=1+\dfrac{1}{x}-\dfrac{1}{x+1}.\\ 
		\end{eqnarray*}
		 Suy ra
		 ${g(1)+g(2)+g(3)+\ldots+g(2025)}$\\
		 ${=\left(1+\dfrac{1}{1}-\dfrac{1}{2}\right)+\left(1+\dfrac{1}{2}-\dfrac{1}{3}\right)+\left(1+\dfrac{1}{3}-\dfrac{1}{4}\right)+\cdots+\left(1+\dfrac{1}{2025}-\dfrac{1}{2026}\right)=2026-\dfrac{1}{2026}}$.\\
		 Khi đó:\\
		\begin{eqnarray*}
			f(1)  \cdot f(2) \cdot f(3) \cdots f(2025)&=&\mathrm e ^ {g(1)+g(2)+g(3)+\cdots+g(2025)} \\
			&=&\mathrm  e ^{2026-\frac{1}{2026}}\\
			&=&\mathrm  e ^{\frac{2026^2-1}{2026}} \\
			&=&\mathrm e  ^{\frac{m}{n}}.\\
		\end{eqnarray*}
		Do  đó ${m=2026^2-1, n=2026}$.\\
		Vậy  ${m-n^2=2026^2-1-2026^2=-1}$.
	} 
\end{bt}

% nguồn: Ok_TuanLuan_Toan11_Lan1
\begin{bt}%[1D6H2-1]%[ma:5860d684]
	\begin{listEX}[2]
		\item $A=2^{2-4\sqrt{5}} \cdot  16^{\sqrt{5}}$.
		\item $B=2  \log _9 18+\log _{\frac{1}{3}} 6$.
	\end{listEX}
	 \loigiai{
		\begin{enumerate}
			\item  $A=2^{2-4\sqrt{5}} \cdot 16^{\sqrt{5}}=2^{2-4\sqrt{5}} \cdot 2^{4\sqrt{5}}=2^{2-4\sqrt{5}+4\sqrt{5}}=2^2=4$.
			\item $B=2 \log _9  18+\log _{\frac{1}{3}}6=2 \log _{3^2} 18+\log _{3^{-1}}6=\log_3 18-\log _3 6=\log _3 3=1$.
		\end{enumerate}
	 }
\end{bt}

% nguồn: Ok_TuanLuan_Toan11_Lan1
\begin{bt}%[1D6H2-1]%[ma:253252c2]
	Cho $\log _a b=3$ và $\log _a c=4$  với $a; b; c>0; a \neq 1$. Tính giá trị của
	 $$
	P=\log _a\left(\dfrac{a^2  \cdot \sqrt{b}}{c^3}\right)
	 $$
	\loigiai{ 
		Ta  có
		 $$
		\begin{aligned}
			P & =\log _a\left(\dfrac{a^2 \cdot \sqrt{b}}{c^3}\right)=\log _a a^2+\log _a \sqrt{b}-\log _a c^3=2+\log _a b^{\tfrac{1}{2}}-3 \log _a c  \\
			& =2+\dfrac{1}{2} \log _a b-3 \log _a c=2+\dfrac{3}{2}-12=\dfrac{-17}{2} . 
		\end{aligned}
		$$ 
	 }
\end{bt}

% nguồn: Ok_TuanLuan_Toan11_Lan1
\begin{bt}%[1D6H2-1]%[5-KNTT]%[ma:430e201a]
	Cho số thực  $a$ thõa mãn $0<a \neq 1$. Tính giá trị của biểu thức $A=2 \log_{2} 12+3 \log_{2} 5-\log_{2} 15-\log_{2} 150$.
	 \loigiai{
		 Ta có $A=\log_{2} 12^2+\log_{2} 5^3-\left(\log_{2} 15+\log_{2} 150\right)=\log_{2} 12^2\cdot 5^3-\log_{2} 15\cdot150=\log_{2} \dfrac{18000}{2250}=\log_{2} 8=3$.
	 }
\end{bt}

% nguồn: Ok_TuanLuan_Toan11_Lan1
\begin{bt}%[1D6H2-1]%[ma:d43de599]
	Cho $\log_ab=2$ và $\log_ac=3$. Tính $  Q=\log_a\left(b^2c^3\right) $.
	\loigiai{Ta có $ Q=\log_a\left(b^2c^3\right)=\log_ab^2+\log_ac^3=2\log_ab+3\log_ac=2\cdot 2+3\cdot3=13 $.} 
\end{bt}

% nguồn: Ok_TuanLuan_Toan11_Lan1
\begin{bt}%[1D6H2-1]%[ma:356234dc]
	Tính giá  trị của biểu thức $A=\log_{8}12-\log_{8}15+\log_{8} 20$.
	 \loigiai{Ta có $A=\log_{8}12-\log_{8}15+\log_{8}20=\log_{8}\left(\dfrac{12\cdot 20}{15} \right)=\log_{8}16=\dfrac{4}{3}$.}
\end{bt}

% nguồn: Ok_TuanLuan_Toan11_Lan2
\begin{ex}%[1D6H2-1]%[ma:609c2de1]
	Cho $\log_a x=3,\log_b x=4$ với $a > 1$, $b > 1$, $x > 1$. Giá trị  của $P=\log _{ab} x$ bằng bao nhiêu?
	\loigiai{ 
		Ta có $P=\log _{ab} x=\dfrac{1}{\log_x  (ab)}=\dfrac{1}{\log_x a+\log_x b}=\dfrac{1}{\dfrac{1}{3}+\dfrac{1}{4}}=\dfrac{12}{7}$.
	 }
\end{ex}

% nguồn: Ok_TuanLuan_Toan11_Lan2
\begin{ex}%[1D6H2-1]%[ma:9e53199c]
	Cho $f(x)=\mathrm e ^{\sqrt{1+\frac{1}{x^2}+\frac{1}{(x+1)^2}}}$. Biết rằng $f(1) \cdot f(2) \cdot f(3) \cdots f(2025)=\mathrm  e ^{\frac{m}{n}}$ với $m, n$ là các số tự nhiên và $\frac{m}{n}$ là phân số tối giản. Tính $m-n^2$.
	\loigiai{  \\
		Đặt  $g(x)=\sqrt{1+\dfrac{1}{x^2}+\dfrac{1}{(1+x)^2}}$ \\
		Với ${x>0}$ ta có\\ 
		\begin{eqnarray*}
			 g(x)&=&\sqrt{1+\dfrac{1}{x^2}+\frac{1}{(1+x)^2}}=\dfrac{\sqrt{x^2+(x+1)^2+x^2\cdot (x+1)^2}}{x(x+1)}=\dfrac{\sqrt{\left(x^2+x+1\right)^2}}{x \left(x+1 \right)} \\
			 &=&\dfrac{{{x}^{2}}+x+1}{x(x+1)}=1+\dfrac{1}{x(x+1)}=1+\dfrac{1}{x}-\dfrac{1}{x+1}.\\
		\end{eqnarray*}
		Suy  ra
		 ${g(1)+g(2)+g(3)+\ldots+g(2025)}$\\
		${=\left(1+\dfrac{1}{1}-\dfrac{1}{2}\right)+\left(1+\dfrac{1}{2}-\dfrac{1}{3}\right)+\left(1+\dfrac{1}{3}-\dfrac{1}{4}\right)+\cdots+\left(1+\dfrac{1}{2025}-\dfrac{1}{2026}\right)=2026-\dfrac{1}{2026}}$.\\ 
		 Khi đó:\\
		\begin{eqnarray*}
			f(1) \cdot f(2) \cdot f(3) \cdots f(2025)&=&\mathrm e  ^ {g(1)+g(2)+g(3)+\cdots+g(2025)} \\
			&=&\mathrm  e ^{2026-\frac{1}{2026}}\\
			&=&\mathrm e  ^{\frac{2026^2-1}{2026}} \\
			&=&\mathrm e  ^{\frac{m}{n}}.\\
		\end{eqnarray*}
		Do đó  ${m=2026^2-1, n=2026}$.\\
		Vậy  ${m-n^2=2026^2-1-2026^2=-1}$.
	} 
\end{ex}

% nguồn: Ok_TuanLuan_Toan11_Lan2
\begin{ex}%[1D6H2-1]%[ma:058d54b7]
	\begin{listEX}[2]
		\item  $A=2^{2-4\sqrt{5}} \cdot 16^{\sqrt{5}}$.
		\item $B=2 \log _9 18+\log  _{\frac{1}{3}} 6$.
	\end{listEX}
	 \loigiai{
		\begin{enumerate}
			\item  $A=2^{2-4\sqrt{5}} \cdot 16^{\sqrt{5}}=2^{2-4\sqrt{5}} \cdot 2^{4\sqrt{5}}=2^{2-4\sqrt{5}+4\sqrt{5}}=2^2=4$.
			\item $B=2 \log _9 18+\log _{\frac{1}{3}}6=2 \log _{3^2} 18+\log _{3^{-1}}6=\log_3 18-\log  _3 6=\log _3 3=1$.
		\end{enumerate}
	 }
\end{ex}

% nguồn: Ok_TuanLuan_Toan11_Lan2
\begin{ex}%[1D6H2-1]%[ma:3dc93904]
	Cho $\log _a b=3$ và $\log _a c=4$ với $a; b; c>0;  a \neq 1$. Tính giá trị của
	$$ 
	P=\log  _a\left(\dfrac{a^2 \cdot \sqrt{b}}{c^3}\right)
	 $$
	\loigiai{ 
		 Ta có
		 $$
		\begin{aligned}
			P & =\log _a\left(\dfrac{a^2 \cdot \sqrt{b}}{c^3}\right)=\log _a a^2+\log _a \sqrt{b}-\log _a c^3=2+\log _a b^{\tfrac{1}{2}}-3  \log _a c \\
			& =2+\dfrac{1}{2} \log _a b-3 \log _a c=2+\dfrac{3}{2}-12=\dfrac{-17}{2} . 
		\end{aligned}
		 $$
	 }
\end{ex}

% nguồn: Ok_TuanLuan_Toan11_Lan2
\begin{ex}%[1D6H2-1]%[5-KNTT]%[ma:3ffbca17]
	Cho  số thực $a$ thõa mãn $0<a \neq 1$. Tính giá trị của biểu thức $A=2 \log_{2} 12+3 \log_{2} 5-\log_{2} 15-\log_{2} 150$.
	\loigiai{ 
		Ta có $A=\log_{2}  12^2+\log_{2} 5^3-\left(\log_{2} 15+\log_{2} 150\right)=\log_{2} 12^2\cdot 5^3-\log_{2} 15\cdot150=\log_{2} \dfrac{18000}{2250}=\log_{2} 8=3$.
	 }
\end{ex}

% nguồn: Ok_TuanLuan_Toan11_Lan2
\begin{ex}%[1D6H2-1]%[ma:eee16493]
	Cho  $\log_ab=2$ và $\log_ac=3$. Tính $ Q=\log_a\left(b^2c^3\right) $.
	\loigiai{Ta có $  Q=\log_a\left(b^2c^3\right)=\log_ab^2+\log_ac^3=2\log_ab+3\log_ac=2\cdot 2+3\cdot3=13 $.}
\end{ex}

% nguồn: Ok_TuanLuan_Toan11_Lan2
\begin{ex}%[1D6H2-1]%[ma:ebdcb5d6]
	Tính giá trị của biểu thức  $A=\log_{8}12-\log_{8}15+\log_{8} 20$.
	\loigiai{Ta có  $A=\log_{8}12-\log_{8}15+\log_{8}20=\log_{8}\left(\dfrac{12\cdot 20}{15} \right)=\log_{8}16=\dfrac{4}{3}$.}
\end{ex}

% nguồn: Ok_TuanLuan_Toan11_Lan2
\begin{ex}%[1D6H2-1]%[ma:3420daef]
	 Tính giá trị của biểu thức $P=\log _28+\log _{\sqrt{3}}9$.
	 \loigiai{
		Ta có $P=\log _28+\log _{\sqrt{3}}9=\log _22^3+\log _{3^{\tfrac{1}{2}}}3^2=3\log _22+4\log  _33=3+4=7$.
	 }
\end{ex}

% nguồn: Ok_TuanLuan_Toan11_Lan2
\begin{ex}%[1D6H2-1]%[ma:247f0b7a]
	Tính giá trị  các biểu thức.
	\begin{enumEX}{3}
		\item $\log_2\dfrac{1}{8}$. 
		\item  $\log_{\sqrt{3}}9$.
		\item  $\log_33\sqrt{3}$.
		 \item $\log_{\tfrac{1}{2}}32$.
		 \item $\log_42+\log_432$.
		 \item $\log_280-\log_25$.
	\end{enumEX}
	 \loigiai{
		\begin{enumerate}
			\item $\log_2\dfrac{1}{8}=\log_2{2^{-3}}=-3$. 
			\item $\log_{\sqrt{3}}9=\log_{\sqrt{3}}\left(\sqrt{3}\right)^4=4$. 
			\item $\log_3\left(3\sqrt{3}\right)=\log_3 3+\log_3\sqrt{3}=1+\log_3 3^{\frac{1}{2}}=1+\dfrac{1}{2}=\dfrac{3}{2}$. 
			\item  $\log_{\tfrac{1}{2}}32=\dfrac{\log_232}{\log_2\dfrac{1}{2}}=\dfrac{5}{-1}=-5$.
			\item  $\log_42+\log_432=\log_4\left(2\cdot32\right)=\log_464=\log_4{4^3}=3\log_44=3$.
			\item 
			$\log_280-\log_25=\log_2\dfrac{80}{5}=\log_216=\log_2{2^4}=4\log_22=4$. 
		\end{enumerate}
	} 
\end{ex}

% nguồn: Ok_TuanLuan_Toan11_Lan2
\begin{ex}%[1D6H2-1]%[ma:6bcbd4d6]
	Tính giá trị biểu thức trong các  trường hợp sau
	\begin{enumEX}{4}
		\item  $\log_22^{-12}$.
		\item $\ln  \mathrm{e}^{\sqrt{2}}$.
		\item $\log_8 16-\log_8 2$. 
		\item $\log_26\cdot\log_68$. 
	\end{enumEX}
	\loigiai{ 
		\begin{enumerate}
			\item  $\log_22^{-12}=-12\log_22=-12$.
			\item  $\ln{\mathrm{e}^{\sqrt{2}}}=\sqrt{2}\ln\mathrm{e}=\sqrt{2}\cdot 1=\sqrt{2}$.
			\item $\log_816-\log_82=\log_8\dfrac{16}{2}=\log_88=1$. 
			 \item $\log_26\cdot\log_68=\log_26\cdot\dfrac{\log_28}{\log_26}=\log_28=\log_22^3=3\log_22=3$.
		\end{enumerate}
	 }
\end{ex}

% nguồn: Ok_TuanLuan_Toan11_Lan2
\begin{ex}%[1D6H2-1]%[ma:0c7f24c2]
	Tính  $(\sqrt[5]{8})^{\log _2243}$.
	\loigiai{ 
		Ta có  $(\sqrt[5]{8})^{\log _2243}=8^{\tfrac{1}{5}\log _23^5}=8^{\log _23}=\left(2^{\log _23}\right)^3=3^3=27$.
	} 
\end{ex}

% nguồn: Ok_TuanLuan_Toan11_Lan2
\begin{ex}%[1D6H2-1]%[ma:8e642bae]
	 Cho $a$ là số thực dương và $a\ne 1$. Tính giá trị của các biểu thức sau:
	\begin{enumEX}{2}
		\item  $T=\log_a\left(a^3\right)$.
		\item  $I=\log_{\tfrac{a}{4}}\left(\dfrac{a^3}{64}\right)$ với $a\ne 4$.
		 \item $T=\log_a\left(\dfrac{a^2\sqrt[3]{a}\sqrt[5]{a^3}}{\sqrt[15]{a^4}}\right)$.
		\item  $\log_{a^2}\left(a\sqrt{a}\right)$.
		 \item $a^{\log_{\sqrt{a}}4}$.
		 \item $P=\log_{\sqrt[3]{a}}\dfrac{1}{a^3}$.
	\end{enumEX}
	\loigiai{ 
		\begin{enumerate}
			\item Ta có $T=\log_a\left(a^3\right)=3\log_aa=3$.\\ 
			\item Ta có $I=\log_{\tfrac{a}{4}}\left(\dfrac{a^3}{64}\right)=\log_{\tfrac{a}{4}}\left(\dfrac{a}{4}\right)^3=3$.\\ 
			\item Ta  có
			\allowdisplaybreaks 
			\begin{align*}
				T&=\log_a\left(\dfrac{a^2\cdot \sqrt[3]{a}\cdot\sqrt[5]{a^3}}{\sqrt[15]{a^4}}\right)\\ 
				&=\log_a\left(\dfrac{a^2\cdot a^{\tfrac{1}{3}}\cdot a^{\tfrac{3}{5}}}{a^{\tfrac{4}{15}}}\right)\\ 
				&=\log_a\left(\dfrac{a^2\cdot a^{\tfrac{1}{3}}\cdot a^{\tfrac{3}{5}}}{a^{\tfrac{4}{15}}}\right)\\ 
				 &=\log_a{a^{\tfrac{8}{3}}}=\dfrac{8}{3}.
			\end{align*}
			\item Ta có $\log_{a^2}\left(a\sqrt{a}\right)=\dfrac{1}{2}{\log_a}{a^{\tfrac{3}{2}}}=\dfrac{1}{2}\cdot\dfrac{3}{2}{\log_a}a=\dfrac{3}{4}$. 
			\item Ta có $a^{\log_{\sqrt{a}}4}=a^{2\log_a4}=a^{\log_a{4^2}}=4^2=16$. 
			\item  Ta có $P=\log_{\sqrt[3]{a}}\dfrac{1}{a^3}=\log_{a^{\tfrac{1}{3}}}a^{-3}=-3\cdot3\log_aa=-9$.
		\end{enumerate}
	 }
\end{ex}

% nguồn: Ok_TuanLuan_Toan11_Lan2
\begin{ex}%[1D6H2-1]%[ma:4e650a5f]
	Cho $a>0$, $a\neq  1$. Tính giá trị của biểu thức $P=\log _{\sqrt[3]{a}}\left(\dfrac{1}{a^3}\right)$.
	 \loigiai{
		Ta có $P=\log _{\sqrt[3]{a}}\left(\dfrac{1}{a^3}\right)=\log  _{a^{\tfrac{1}{3}}}a^{-3}=-9\log _aa=-9$.
	} 
\end{ex}

% nguồn: Ok_TuanLuan_Toan11_Lan2
\begin{ex}%[1D6H2-1]%[ma:efaa9067]
	Cho số thực dương $a$ khác $1$. Tính giá  trị của biểu thức $\log _2(4a)$.
	\loigiai{ 
		Ta có $\log _2(4a)=\log _24+\log _2a=2+\log _2a$. 
	} 
\end{ex}

% nguồn: Ok_TuanLuan_Toan11_Lan2
\begin{ex}%[1D6H2-1]%[ma:ad79aeba]
	Cho số thực $0<a\neq 1$. Tính giá trị của biểu thức $\log _{\sqrt{a}}\left(a^2\sqrt[3]{a}\right)$. 
	 \loigiai{
		Ta  có $\log _{\sqrt{a}}\left(a^2\sqrt[3]{a}\right)=\log _{a^{\tfrac{1}{2}}}\left(a^2\cdot a^{\tfrac{1}{3}}\right)=\log _{a^{\tfrac{1}{2}}}\left(a^{\tfrac{7}{3}}\right)=\dfrac{\dfrac{7}{3}}{\dfrac{1}{2}}=\dfrac{14}{3}$.
	 }
\end{ex}

% nguồn: Ok_TuanLuan_Toan11_Lan1
\begin{bt}%[1D6V2-1]%[ma:619ba33a]
	Cho số thực $a$ thỏa mãn $0<a\neq1$. Tính giá trị của biểu thức $T=\log_a\left(\dfrac{a^2\cdot\sqrt[3]{a^2}\cdot\sqrt[5]{a^4}}{\sqrt[15]{a^7}}\right)$. 
	 \loigiai{
		 $T=\log_a\left(\dfrac{a^2\cdot\sqrt[3]{a^2}\cdot\sqrt[5]{a^4}}{\sqrt[15]{a^7}}\right)=\log_a\dfrac{a^{2+\tfrac{2}{3}+\tfrac{4}{5}}}{a^{\tfrac{7}{15}}}=\log_aa^{2+\tfrac{2}{3}+\tfrac{4}{5}-\tfrac{7}{15}}=\log_aa^3=3$.
	 }
\end{bt}

% nguồn: Ok_TuanLuan_Toan11_Lan1
\begin{bt}%[1D6V2-1]%[ma:0a28f48b]
	Cho $a=\log_5{18}$ và $b=\log_5{60}$. Tính $\log_32$ theo $a$ và  $b$.
	\loigiai{Đầu tiên ta có hệ $\heva{&a=\log_5{18}=\log_52+2\log_53\\&b=\log_5{60}=2\log_52+\log_53+1.}$  \\
		Đặt  $x=\log_52$ và $y=\log_53$ từ đó ta có hệ phương trình bậc nhất hai ẩn
		\begin{eqnarray*}
			 && \heva{&x+2y=a\\&2x+y=b-1}\\
			&\Leftrightarrow&  \heva{&x=\log_52=\dfrac{-a+2b-2}{3}\\&y=\log_53=\dfrac{2a-b+1}{3}.}
		\end{eqnarray*}
		Vậy $\log_32=\dfrac{\log_52}{\log_5  3}=\dfrac{-a+2b-2}{2a-b+1}$.
	} 
\end{bt}

% nguồn: Ok_TuanLuan_Toan11_Lan2
\begin{ex}%[1D6V2-1]%[ma:3fa0e88a]
	Cho  số thực $a$ thỏa mãn $0<a\neq1$. Tính giá trị của biểu thức $T=\log_a\left(\dfrac{a^2\cdot\sqrt[3]{a^2}\cdot\sqrt[5]{a^4}}{\sqrt[15]{a^7}}\right)$.
	\loigiai{ 
		$T=\log_a\left(\dfrac{a^2\cdot\sqrt[3]{a^2}\cdot\sqrt[5]{a^4}}{\sqrt[15]{a^7}}\right)=\log_a\dfrac{a^{2+\tfrac{2}{3}+\tfrac{4}{5}}}{a^{\tfrac{7}{15}}}=\log_aa^{2+\tfrac{2}{3}+\tfrac{4}{5}-\tfrac{7}{15}}=\log_aa^3=3$. 
	 }
\end{ex}

% nguồn: Ok_TuanLuan_Toan11_Lan2
\begin{ex}%[1D6V2-1]%[ma:24fd0521]
	Cho  $a=\log_5{18}$ và $b=\log_5{60}$. Tính $\log_32$ theo $a$ và $b$.
	\loigiai{Đầu tiên ta  có hệ $\heva{&a=\log_5{18}=\log_52+2\log_53\\&b=\log_5{60}=2\log_52+\log_53+1.}$ \\
		Đặt $x=\log_52$ và $y=\log_53$ từ  đó ta có hệ phương trình bậc nhất hai ẩn
		\begin{eqnarray*}
			 && \heva{&x+2y=a\\&2x+y=b-1}\\
			&\Leftrightarrow& \heva{&x=\log_52=\dfrac{-a+2b-2}{3}\\&y=\log_53=\dfrac{2a-b+1}{3}.} 
		\end{eqnarray*}
		Vậy $\log_32=\dfrac{\log_52}{\log_5 3}=\dfrac{-a+2b-2}{2a-b+1}$. 
	 }
\end{ex}

% nguồn: Ok_TuanLuan_Toan11_Lan2
\begin{ex}%[1D6V2-1]%[ma:9ce5d563]
	Cho $\log  _23=a$, $\log _25=b$. Tính $\log _{15}8$.
	\loigiai{ 
		Ta có $\log _{15}8=\dfrac{1}{\log _815}=\dfrac{1}{\log _{2^3}(3\cdot 5)}=\dfrac{1}{\dfrac{1}{3}\left(\log _23+\log _25\right)}=\dfrac{3}{\log _23+\log _25}=\dfrac{3}{a+b}$.\\ 
		Vậy $\log _{15}8=\dfrac{3}{a+b}$.\\ 
	} 
\end{ex}

% nguồn: Ok_TuanLuan_Toan11_Lan2
\begin{ex}%[1D6V2-1]%[ma:5e892126]
	Cho ba số thực dương $a$, $b$, $c$ theo thứ tự  lập thành một cấp số nhân và $a+b+c=64$. Tính giá trị của biểu thức $P=3\log_{2}\left(ab+bc+ca\right)-\log_{2}\left(abc\right)$.
	 \loigiai{
		\begin{itemize}
			\item  Do $a$, $b$, $c$ theo thứ tự lập thành một cấp số nhân nên ta có $ac = b^2$.
			\item Biểu thức $P$ được tính  như sau
			\begin{eqnarray*}
				P&=&3\log_{2}\left(ab+bc+ca\right)-\log_{2}\left(abc\right)=3\log_{2}\left(ab+bc+b^{2}\right)-\log_{2}\left(b  \cdot b^{2}\right)\\
				 &=&3\log_{2}b\left(a+b+c\right)-\log_{2}\left(b^{3}\right)=3\log_{2}\left(64b\right)-\log_{2}\left(b^{3}\right)\\
				 &=&\log_{2}\left(64b\right)^{3}-\log_{2}\left(b^{3}\right)=\log_{2}\dfrac{\left(2^{6}b\right)^{3}}{b^{3}}\\
				&=&\log_{2}\dfrac{2^{18}b^{3}}{b^{3}}=\log_{2}2^{18}=18. 
			\end{eqnarray*}
		\end{itemize}
	 }
\end{ex}

% nguồn: Ok_TuanLuan_Toan11_Lan2
\begin{ex}%[1D6V2-1]%[ma:400c0dc8]
	Cho các số dương $a$, $b$, $c$ khác  $1$ thoả mãn $\log _a(b c)=3$ và $\log _b(a c)=4$. Tính giá trị của biểu thức $\log _c(a b)$.
	\loigiai{ 
		 Ta có
		 \allowdisplaybreaks
		\begin{eqnarray*}
			\heva{&\log _a(b c)=3 \\& \log _b(a c)=4} &\Leftrightarrow&\heva{&b c=a^3 \\& a c=b^4}\\ &\Leftrightarrow&\heva{&b=a^3 \cdot c^{-1} \\ &a  c=\left(a^3 \cdot c^{-1}\right)^4}\\ &\Leftrightarrow&\heva{&b=a^3 \cdot c^{-1} \\ &a c=a^{12} \cdot c^{-4}}\\ &\Leftrightarrow&\heva{&a=c^{\tfrac{5}{11}} \\& b=c^{\tfrac{4}{11}}}.
		\end{eqnarray*}
		 Lúc đó
		$ \log  _c(a b)=\log _c\left(c^{\tfrac{5}{11}}\cdot c^{\tfrac{4}{11}}\right)=\log _c c^{\tfrac{9}{11}}=\dfrac{9}{11}
		$. 
	 }
\end{ex}

% nguồn: Ok_TuanLuan_Toan11_Lan2
\begin{ex}%[1D6V2-1]%[ma:414ddf4c]
	Cho $a$, $b$ là các số thực dương thỏa mãn $\log _4 a=\log _{25} b=\log \dfrac{4 b-a}{4}$. Tính giá trị của biểu thức  $\log _{\sqrt{6}}\left(\dfrac{a}{2}+4 b \sqrt{2}\right)-\log_{\sqrt{6}} b$.
	 \loigiai{
		Đặt $\log _4 a=\log _{25} b=\log \dfrac{4 b-a}{4}=t \Leftrightarrow\heva{&a=4^t \\&  b=25^t \\& \dfrac{4 b-a}{4}=10^t}$.\\
		 Suy ra
		\allowdisplaybreaks 
		\begin{eqnarray*}
			\dfrac{4\cdot 25^t-4^t}{4}=10^t &\Leftrightarrow& 4\cdot 25^t-4^t=4\cdot  10^t\\
			&  \Leftrightarrow& 4 \cdot\left(\dfrac{5}{2}\right)^{2 t}-4 \cdot\left(\dfrac{5}{2}\right)^t-1=0\\
			 & \Leftrightarrow&\left(\dfrac{5}{2}\right)^t=\dfrac{1+\sqrt{2}}{2}\\
			& \Rightarrow& \dfrac{a}{b}=\left(\dfrac{2}{5}\right)^{2  t}=12-8 \sqrt{2}.
		\end{eqnarray*}
		 Lúc đó
		\allowdisplaybreaks 
		\begin{eqnarray*}
			\log _{\sqrt{6}}\left(\dfrac{a}{2}+4 b \sqrt{2}\right)-\log  _{\sqrt{6}} b&=&\log _{\sqrt{6}}\left(\dfrac{\dfrac{a}{2}+4 b \sqrt{2}}{b}\right)\\
			&=&\log _{\sqrt{6}}\left(\dfrac{a}{2 b}+4 \sqrt{2}\right)\\ 
			&=&\log _{\sqrt{6}}(6-4 \sqrt{2}+4  \sqrt{2})\\
			&=&\log  _{\sqrt{6}}(6)=2.
		\end{eqnarray*}
	} 
\end{ex}

% nguồn: Ok_TuanLuan_Toan11_Lan2
\begin{ex}%[1D6V2-1]%[ma:67ff3dad]
	Cho các số $a$, $b$, $c$ thỏa mãn $\log _a 3=2$, $\log _b 3=\dfrac{1}{4}$, $\log _{abc}3=\dfrac{2}{15}$. Tính  giá trị của $\log_c 3$.
	\loigiai{ 
		Ta  có $\log _a 3=2 \Rightarrow \log _3 a=\dfrac{1}{2} \Leftrightarrow \log _b 3=\frac{1}{4} \Rightarrow \log _3 b=4$.\\
		Lại có $\log _{a b c} 3=\dfrac{2}{15} \Leftrightarrow \log _3abc=\dfrac{15}{2} \Leftrightarrow \log _3 a+\log  _3 b+\log _3 c=\dfrac{15}{2}$.\\
		 Suy ra
		$$ 
		\log _3  c=\dfrac{15}{2}-\dfrac{1}{2}-4=3 \Leftrightarrow \log _c 3=\dfrac{1}{3}.
		$$ 
	} 
\end{ex}

% nguồn: Ok_TuanLuan_Toan11_Lan2
\begin{ex}%[1D6C2-1]%[ma:ec10973e]
	Cho $a$, $b$, $c$  là ba số thực dương khác $1$ và $abc\ne1$. Biết $\log_{a}5=3$, $\log_{b}5=4$, $\log_{abc}5=\dfrac{10}{17}$. Khi đó giá trị của $\log_c 5$ bằng bao nhiêu?
	\loigiai{ 
		\begin{itemize}
			\item Từ giả thiết  ta có\\
			 $\log_{a}5=3\Rightarrow \log_{5}a=\dfrac{1}{3}$.\\
			$\log_{b}5=4\Rightarrow  \log_{5}b=\dfrac{1}{4}$.\\
			 $\log_{abc}5=\dfrac{10}{17}\Rightarrow \log_{5}\left(abc\right)=\dfrac{17}{10}$.
			\item Ta có $\log_{5}\left(abc\right)=\log_{5}a+\log_{5}b+\log_{5}c\Rightarrow  \log_{5}c=\log_{5}\left(abc\right)-\log_{5}a-\log_{5}b=\dfrac{17}{10}-\dfrac{1}{3}-\dfrac{1}{4}=\dfrac{67}{60}$.
			\item Vậy  $\log_{c}5=\dfrac{60}{67}$.
		\end{itemize}
	} 
\end{ex}
